\documentclass[12pt,a4paper]{article}

% --- packages ---
\usepackage{amsmath,amssymb,mathtools}
\usepackage{geometry}
\usepackage{microtype}
\usepackage[shortlabels]{enumitem} % <— only enumitem, shortlabels enables [(a)] if you want it
\usepackage{lmodern}
\usepackage[T1]{fontenc}
\usepackage{graphicx}
\usepackage{subcaption}
\graphicspath{{figures/}}
\usepackage{tikz}
\usetikzlibrary{arrows.meta}
\usepackage{centernot}
\newcommand{\qed}{\hfill$\square$}
\usepackage{booktabs,tabularx}

\usepackage{fancyhdr}
\pagestyle{fancy}


\usepackage[most]{tcolorbox}

\usepackage{tikz-cd} 

\usepackage{fontspec}

\usepackage{amsmath,amssymb,amsthm}

\newenvironment{statement}{\par\noindent\textbf{Statement.}\ }{\par}

% --- theorem setups ---
\theoremstyle{plain}
\newtheorem{theorem}{Theorem}[section]
\newtheorem{proposition}[theorem]{Proposition}
\newtheorem{lemma}[theorem]{Lemma}
\newtheorem{corollary}[theorem]{Corollary}

\theoremstyle{definition}
\newtheorem{definition}[theorem]{Definition}
\newtheorem{example}[theorem]{Example}

\theoremstyle{remark}
\newtheorem*{remark}{Remark} % unnumbered


\theoremstyle{remark}
\newtheorem*{counterexample}{Counterexample}

\DeclareMathOperator{\sgn}{sgn}
\DeclareMathOperator{\csch}{csch}
\DeclareMathOperator{\Arg}{Arg}
\DeclareMathOperator{\Res}{Res}
\DeclareMathOperator{\Log}{Log}

\newcommand{\C}{\mathbb{C}}
\newcommand{\R}{\mathbb{R}}

 \newcommand{\gap}{\vspace{6pt}}
 
\newcommand{\axes}[3]{%
	\draw[->,thin] (#1) -- ++(2.0,0) node[right] {#2};
	\draw[->,thin] (#1) -- ++(0,2.0) node[above] {#3};
}

\setlength{\parindent}{0pt}
\setlength{\parskip}{6pt}
\setlist{itemsep=6pt, topsep=6pt, parsep=2pt, partopsep=2pt}

% SCOPED header: bold ends immediately after the line
\newcommand{\probheader}[1]{%
	\par\vspace{8pt}%
	{\Large\bfseries #1\par}% <- braces scope the bold
	\vspace{6pt}%
}

% Thin separator (optional)
\newcommand{\sep}{\vspace{6pt}\hrule\vspace{8pt}}


\fancyhf{}

% --- HEADER ---
\fancyhead[L]{\textit{Analysis II extended}}
\fancyhead[R]{\textit{Exercises on Topologies of Function Spaces}}


\fancyfoot[C]{\thepage}

\geometry{margin=1in}
\setlength{\parskip}{0.9em}
\setlength{\parindent}{0pt}
\linespread{1.12}
\setlist[itemize]{topsep=0.3em,itemsep=0.2em}
\setlist[enumerate]{topsep=0.3em,itemsep=0.2em}

\title{\Large Analysis II}
\author{}
\date{}


\begin{document}
	
	\maketitle
	
	
\section*{Problem 1 — Quickies (uniform convergence \& regularity)}

\paragraph{(a)} If $(f_n)$ is a sequence of real functions on $[0,1]$ converging uniformly to a function $f$ on $[0,1]$, and if $f_n$ is continuous at $x_n\in[0,1]$ with $x_n\to x$, does it follow that $f$ is continuous at $x$?

\paragraph{(b)} If $(f_n)$ is a sequence of continuous real functions on $[-1,1]$ converging pointwise to a continuous function $f$ on $[-1,1]$, and if the convergence is uniform on $[-r,r]$ for every $r\in(0,1)$, does it follow that the convergence is uniform on $[-1,1]$?

\paragraph{(c)} If $(f_n)$ is a sequence of real functions on the interval $[0,1]$ converging uniformly to a function $f$ on $[0,1]$, and if each $f_n$ is continuous except at countably many points, does it follow that there exists a point at which $f$ is continuous?

\paragraph{(d)$^\ast$} If $(f_n)$ is a sequence of differentiable functions on $[0,1]$ converging uniformly to a function $f$ on $[0,1]$, does it follow that there exists a point at which $f$ is differentiable?

\bigskip

\subsection*{Theory used}
\begin{itemize}
	\item \textbf{Continuity at a point is preserved by uniform limits:} If $f_n\to f$ uniformly on a set $E$ and each $f_n$ is continuous at $x_0\in E$, then $f$ is continuous at $x_0$.
	\item \textbf{Uniform limit of continuous functions is continuous.}
	\item \textbf{Countable union of countable sets is countable.}
	\item \textbf{Local vs global uniformity:} Uniform convergence on $[-r,r]$ for every $r<1$ need not imply uniform convergence on $[-1,1]$.
	\item \textbf{Weierstrass phenomenon:} There exist continuous nowhere differentiable functions that are uniform limits of differentiable functions (e.g., partial sums of the Weierstrass series).
\end{itemize}

\subsection*{Solutions}

\paragraph{(a) No.}
Let
\[
f(x)=\begin{cases}
	0,& x<\tfrac12,\\
	1,& x\ge \tfrac12,
\end{cases}
\qquad \text{and set } f_n=f \text{ for all } n.
\]
Then $f_n\to f$ uniformly. Define $x_n=\tfrac12+\tfrac1n\to x=\tfrac12$. For each $n$, $f_n$ is continuous at $x_n$ (indeed $f$ is continuous at all $x\ne \tfrac12$), but $f$ is discontinuous at $x=\tfrac12$. Hence the asserted implication fails.

\paragraph{(b) No.}
For $n\ge1$, define
\[
f_n(x)=
\begin{cases}
	n(1-x),& x\in[1-\tfrac1n,\,1],\\
	0,& x\in[-1,\,1-\tfrac1n),
\end{cases}
\qquad\text{and let } f\equiv 0.
\]
Each $f_n$ is continuous on $[-1,1]$. For any fixed $x\in[-1,1)$, $f_n(x)=0$ for all large $n$, so $f_n(x)\to 0=f(x)$. Also $f_n(1)=0\to 0=f(1)$, hence $f_n\to f$ pointwise and $f$ is continuous. Given any $r\in(0,1)$, if $n>\frac1{1-r}$ then $[1-\frac1n,1]\subset (r,1]$, so $f_n\equiv 0$ on $[-r,r]$; thus convergence is uniform on $[-r,r]$. However,
\[
\sup_{x\in[-1,1]}|f_n(x)-f(x)|=\max_{x\in[-1,1]} f_n(x)=1 \quad\text{for all }n,
\]
so the convergence is not uniform on $[-1,1]$.

\paragraph{(c) Yes.}
Let $D_n$ be the (countable) set of discontinuities of $f_n$ and $D=\bigcup_{n=1}^\infty D_n$, which is countable. Choose $x_0\in[0,1]\setminus D$. Then each $f_n$ is continuous at $x_0$. Since $f_n\to f$ uniformly, the continuity-at-a-point lemma implies $f$ is continuous at $x_0$. Therefore $f$ is continuous at (at least) one point.

\paragraph{(d)$^\ast$ No.}
Consider the Weierstrass function
\[
W(x)=\sum_{k=0}^\infty a^k\cos(b^k\pi x),
\]
with $0<a<1$, $b\in\mathbb{N}$ odd, and $ab>1$. Let $S_n(x)=\sum_{k=0}^n a^k\cos(b^k\pi x)$. Then each $S_n$ is differentiable on $[0,1]$, and $S_n\to W$ uniformly on $[0,1]$, while $W$ is continuous everywhere and nowhere differentiable. Thus there need not exist any point where $f$ is differentiable.
	
	
\section*{Problem 2 — Which $(f_n)$ converge uniformly on $X$?}

\begin{enumerate}[(a)]
	\item $f_n(x)=x^n$ on $X=(0,1)$.
	\item $f_n(x)=x^n$ on $X=(0,\tfrac12)$.
	\item $f_n(x)=x e^{-n x}$ on $X=[0,\infty)$.
	\item $f_n(x)=e^{-x^2}\sin(x/n)$ on $X=\mathbb{R}$.
\end{enumerate}

\subsection*{Theory}
Uniform convergence on $X$ means $\sup_{x\in X}|f_n(x)-f(x)|\to0$. If $|f_n(x)|\le M_n$ for all $x\in X$ with $M_n\to0$, then $f_n\to0$ uniformly. We shall also use $|\sin t|\le|t|$ and the calculus fact $\sup_{x\ge0} x e^{-a x}=1/(a e)$, attained at $x=1/a$.

\subsection*{Solutions}

\paragraph{(a) Not uniform.}
For each $x\in(0,1)$ we have $x^n\to0$, but
\[
\sup_{x\in(0,1)} |x^n|=\sup_{x\in(0,1)} x^n=1,
\]
so the supremum does not tend to $0$. Therefore the convergence is not uniform on $(0,1)$.

\paragraph{(b) Uniform.}
For $x\in(0,\tfrac12)$,
\[
0\le x^n\le \Big(\tfrac12\Big)^n,
\]
hence
\[
\sup_{x\in(0,1/2)}|x^n|\le \Big(\tfrac12\Big)^n\to0.
\]
Thus $x^n\to0$ uniformly on $(0,\tfrac12)$.

\paragraph{(c) Uniform.}
For fixed $n$, the maximum of $x\mapsto x e^{-n x}$ on $[0,\infty)$ occurs at $x=1/n$, with value $1/(n e)$. Consequently
\[
\sup_{x\in[0,\infty)} |x e^{-n x}|=\frac{1}{n e}\to0,
\]
so $f_n\to0$ uniformly on $[0,\infty)$.

\paragraph{(d) Uniform.}
Using $|\sin t|\le |t|$,
\[
|f_n(x)|\le e^{-x^2}\frac{|x|}{n}.
\]
The function $g(x)=|x|e^{-x^2}$ is bounded on $\mathbb{R}$ with $\sup g=1/\sqrt{2e}$ (attained at $|x|=1/\sqrt2$). Hence
\[
\sup_{x\in\mathbb{R}} |f_n(x)|
\le \frac{1}{n}\sup_{x\in\mathbb{R}} |x|e^{-x^2}
= \frac{1}{n\sqrt{2e}}\to0,
\]
and $f_n\to0$ uniformly on $\mathbb{R}$.

\subsection*{Summary}
(a) not uniform; (b) uniform; (c) uniform; (d) uniform.


\section*{Problem 3 — Uniform limits: sums vs products}

Let $(f_n)$ and $(g_n)$ be sequences of real-valued functions on a set $E$ converging uniformly to $f$ and $g$, respectively.

\subsection*{(1) Sums converge uniformly}
By the triangle inequality,
\[
\sup_{x\in E}\big|(f_n+g_n)-(f+g)\big|
\le \sup_{x\in E}|f_n-f|+\sup_{x\in E}|g_n-g|\xrightarrow[n\to\infty]{}0.
\]
Thus $f_n+g_n\to f+g$ uniformly on $E$.

\subsection*{(2) Products need not converge uniformly}
Consider $E=\mathbb{R}$ and define $f_n(x)\equiv \frac1n$, $f(x)\equiv 0$, and $g_n(x)\equiv g(x)=x$. Then
\[
\sup_{x\in\mathbb{R}}|f_n(x)-f(x)|=\frac1n\to0,\qquad
\sup_{x\in\mathbb{R}}|g_n(x)-g(x)|=0,
\]
so $f_n\to f$ and $g_n\to g$ uniformly. However
\[
f_n(x)g_n(x)=\frac{x}{n},\qquad
\sup_{x\in\mathbb{R}}\Big|\frac{x}{n}-0\Big|=\infty,
\]
so $(f_ng_n)$ does not converge uniformly to $fg$ on $\mathbb{R}$ (even though it converges pointwise to $0$).

\subsection*{(3) If $f$ and $g$ are bounded, then products converge uniformly}
Suppose $|f|\le M$ and $|g|\le L$ on $E$. Since $g_n\to g$ uniformly, there exists $N_0$ such that for all $n\ge N_0$,
\[
\sup_{x\in E}|g_n(x)-g(x)|<1\quad\Rightarrow\quad
\sup_{x\in E}|g_n(x)|\le L+1.
\]
Then for $n\ge N_0$,
\[
\begin{aligned}
	\sup_{x\in E}|f_n(x)g_n(x)-f(x)g(x)|
	&\le \sup_{x\in E}|f_n(x)-f(x)|\,\sup_{x\in E}|g_n(x)|
	+ \sup_{x\in E}|f(x)|\,\sup_{x\in E}|g_n(x)-g(x)| \\
	&\le (L+1)\sup_{x\in E}|f_n-f| + M\sup_{x\in E}|g_n-g|\xrightarrow[n\to\infty]{}0.
\end{aligned}
\]
Hence $f_ng_n\to fg$ uniformly on $E$.

\subsection*{(4) If $f$ is bounded but $g$ is not}
Uniform convergence of $(f_ng_n)$ to $fg$ need not hold. The example in (2) has bounded $f\equiv 0$, unbounded $g(x)=x$, uniform convergence $f_n\to f$ and $g_n\to g$, yet $(f_ng_n)$ fails to converge uniformly to $fg$.

\section*{Problem 4 — Uniform limit of bounded functions; pointwise limit can be unbounded}

\paragraph{Statement.}
Let $(f_n)$ be a sequence of bounded, real-valued functions on a set $E$ converging uniformly to a function $f$.
Show that $f$ must be bounded. Give an example of a sequence $(g_n)$ of bounded, real-valued
functions on $[-1,1]$ converging pointwise to a function $g$ which is not bounded.

\subsection*{Solution}

\paragraph{(1) Uniform limit of bounded functions is bounded.}
Since $f_n\to f$ uniformly on $E$, there exists $N$ such that
\[
\sup_{x\in E}|f(x)-f_N(x)|<1.
\]
Because $f_N$ is bounded, let $B=\sup_{x\in E}|f_N(x)|<\infty$. Then, for all $x\in E$,
\[
|f(x)|\le |f(x)-f_N(x)|+|f_N(x)|\le 1+B,
\]
so $f$ is bounded.

\paragraph{(2) Pointwise limit that is unbounded.}
Define
\[
g(x)=
\begin{cases}
	\frac{1}{|x|}, & x\neq 0,\\
	0, & x=0,
\end{cases}
\qquad x\in[-1,1].
\]
Then $g$ is real-valued and unbounded on $[-1,1]$.
For $n\in\mathbb{N}$, define
\[
g_n(x)=
\begin{cases}
	\min\!\big(n,\tfrac{1}{|x|}\big), & x\neq 0,\\
	0, & x=0.
\end{cases}
\]
Each $g_n$ is bounded on $[-1,1]$ by $n$. For $x\neq 0$ we have $g_n(x)\to \frac{1}{|x|}$, and for $x=0$ we have $g_n(0)=0=g(0)$. Hence $g_n\to g$ pointwise on $[-1,1]$, while $g$ is unbounded.

\section*{Problem 5 — Moving points and uniform convergence}

Let $(f_n)$ be a sequence of real-valued continuous functions on a closed, bounded interval $[a,b]$,
and suppose that $f_n$ converges pointwise to a continuous function $f$.

\subsection*{(1) If $f_n\to f$ uniformly and $x_n\to x$, then $f_n(x_n)\to f(x)$}
For any sequence $x_n\in[a,b]$ with $x_n\to x$,
\[
|f_n(x_n)-f(x)|
\le |f_n(x_n)-f(x_n)|+|f(x_n)-f(x)|
\le \sup_{y\in[a,b]}|f_n(y)-f(y)|+|f(x_n)-f(x)|.
\]
The first term tends to $0$ by uniform convergence, the second by continuity of $f$ at $x$.
Hence $f_n(x_n)\to f(x)$.

\subsection*{(2) If $f_n$ does not converge uniformly to $f$, there exists $x_n\to x$ with $f_n(x_n)\not\to f(x)$}
Assume $f_n\nrightarrow f$ uniformly. Then there exists $\varepsilon>0$ such that for every $N$
one can find $n\ge N$ and $x\in[a,b]$ with
\[
|f_n(x)-f(x)|\ge \varepsilon.
\]
Inductively choose $n_k\uparrow\infty$ and $x_k\in[a,b]$ so that $|f_{n_k}(x_k)-f(x_k)|\ge \varepsilon$ for all $k$.
By compactness of $[a,b]$ there is a subsequence (relabeled) with $x_k\to x\in[a,b]$.
Since $f$ is continuous, $f(x_k)\to f(x)$, and therefore
\[
\liminf_{k\to\infty} |f_{n_k}(x_k)-f(x)|
\ge \liminf_{k}\Big(|f_{n_k}(x_k)-f(x_k)|-|f(x_k)-f(x)|\Big)
\ge \varepsilon.
\]
Thus $f_{n_k}(x_k)\not\to f(x)$. Defining $x_n$ by $x_{n_k}=x_k$ (and arbitrarily elsewhere) yields a sequence $x_n\to x$
with $f_n(x_n)\not\to f(x)$.

\section*{Problem 5 — Moving points vs.\ uniform convergence on $[a,b]$}

\paragraph{Statement.}
Let $(f_n)$ be a sequence of real-valued continuous functions on the compact interval $[a,b]$,
and suppose $f_n \to f$ pointwise with $f$ continuous on $[a,b]$.
\begin{enumerate}
	\item[(i)] If $f_n \to f$ uniformly and $x_n \to x$ in $[a,b]$, prove that $f_n(x_n)\to f(x)$.
	\item[(ii)] Conversely, if $f_n \not\to f$ uniformly on $[a,b]$, prove that there exist $x_n\to x$ in $[a,b]$
	such that $f_n(x_n)\not\to f(x)$.
\end{enumerate}

\subsection*{Theory}
We collect the tools used in the proofs.

\begin{definition}[Uniform convergence]
	A sequence $(f_n)$ of functions on $E$ converges uniformly to $f$ if
	\[
	\sup_{y\in E} |f_n(y)-f(y)| \xrightarrow[n\to\infty]{} 0.
	\]
\end{definition}

\begin{lemma}[Uniform convergence controls moving points]\label{lem:move}
	Let $E$ be a metric space. If $f_n\to f$ uniformly on $E$ and $x_n\to x$ in $E$ with $f$ continuous at $x$, then
	\[
	|f_n(x_n)-f(x)| \le \underbrace{|f_n(x_n)-f(x_n)|}_{\le \sup_E |f_n-f|} + \underbrace{|f(x_n)-f(x)|}_{\to 0}
	\longrightarrow 0.
	\]
\end{lemma}

\begin{lemma}[Compactness selection]\label{lem:compact}
	If $(x_n)$ is a sequence in a compact set $K$, then there exists a convergent subsequence $x_{n_k}\to x\in K$.
\end{lemma}

\begin{lemma}[Witness for failure of uniform convergence]\label{lem:witness}
	Let $K$ be compact. If $f_n \not\to f$ uniformly on $K$, then there exists $\varepsilon>0$, a strictly increasing $n_k$,
	and points $x_k\in K$ such that $|f_{n_k}(x_k)-f(x_k)|\ge \varepsilon$ for all $k$.
\end{lemma}

\begin{remark}[Equivalence on compact sets]\label{rem:equiv}
	Assume each $f_n$ is continuous on $[a,b]$ and $f$ is continuous. Then the following are equivalent:
	\begin{enumerate}
		\item $f_n \to f$ uniformly on $[a,b]$.
		\item For every sequence $x_n\to x$ in $[a,b]$, we have $f_n(x_n)\to f(x)$.
	\end{enumerate}
	The direction (1)$\Rightarrow$(2) is Lemma~\ref{lem:move}. The converse follows from (ii) below.
\end{remark}

\subsection*{Proofs}

\paragraph{(i) Uniform convergence $\Rightarrow$ moving values converge.}
Let $x_n\to x$ in $[a,b]$. By Lemma~\ref{lem:move},
\[
|f_n(x_n)-f(x)| \le \sup_{y\in[a,b]}|f_n(y)-f(y)| + |f(x_n)-f(x)| \xrightarrow[n\to\infty]{} 0,
\]
because the first term $\to 0$ by uniform convergence and the second by continuity of $f$ at $x$.

\paragraph{(ii) Failure of uniform convergence $\Rightarrow$ a bad moving sequence.}
Assume $f_n \not\to f$ uniformly on $[a,b]$. By Lemma~\ref{lem:witness} there exist $\varepsilon>0$, indices $n_k\uparrow\infty$,
and points $x_k\in[a,b]$ such that
\[
|f_{n_k}(x_k)-f(x_k)| \ge \varepsilon \qquad \text{for all } k.
\]
By compactness (Lemma~\ref{lem:compact}), pass to a subsequence (not relabeled) with $x_k\to x\in[a,b]$.
Since $f$ is continuous, $f(x_k)\to f(x)$. Then
\[
\liminf_{k\to\infty} |f_{n_k}(x_k)-f(x)|
\;\ge\; \liminf_{k\to\infty}\Big(|f_{n_k}(x_k)-f(x_k)| - |f(x_k)-f(x)|\Big)
\;\ge\; \varepsilon - 0 \;=\; \varepsilon.
\]
Hence $f_{n_k}(x_k)\not\to f(x)$. Defining a sequence $(x_n)$ by $x_{n_k}=x_k$ and setting $x_n=x$ for other $n$
gives $x_n\to x$ with $f_n(x_n)\not\to f(x)$.

\subsection*{Examples}

\paragraph{Example 1 (Spike/hat example: pointwise convergence without uniformity).}
Let
\[
\phi(t) = \max\{1-|t|,\,0\}, \qquad t\in\mathbb{R}.
\]
Define $f_n:[0,1]\to\mathbb{R}$ by
\[
f_n(x) = \phi(nx-1).
\]
Then each $f_n$ is continuous, $f_n(x)\to 0$ for every $x\in[0,1]$ (so $f\equiv 0$), but
\[
\sup_{x\in[0,1]}|f_n(x)-f(x)|=\sup_{x\in[0,1]} f_n(x)=1 \quad \text{for all } n,
\]
so the convergence is \emph{not} uniform. Taking $x_n=1/n\to 0$ yields $f_n(x_n)=1\nrightarrow 0=f(0)$,
which witnesses the conclusion of (ii).

\paragraph{Example 2 (Uniform convergence implies moving-point convergence).}
On $[0,1]$, let $f_n(x)=x^n$ and $f\equiv 0$. Then for any $\delta\in(0,1)$, $f_n\to f$ uniformly on $[0,1-\delta]$ with
\[
\sup_{x\in[0,1-\delta]}|f_n(x)-f(x)| \le (1-\delta)^n \to 0.
\]
Thus for any $x_n\to x\in[0,1-\delta]$, Lemma~\ref{lem:move} gives $f_n(x_n)\to f(x)=0$.
(Contrast with the boundary case $x_n\to 1$, where uniformity fails on $[0,1]$ and moving-point convergence can fail.)

\paragraph{Example 3 (Uniform in the tail: decaying envelopes).}
Let $f_n(x)=x e^{-n x}$ on $[0,\infty)$. Then
\[
\sup_{x\ge 0} |f_n(x)| = \frac{1}{n e}\to 0,
\]
so $f_n\to 0$ uniformly. Hence for any $x_n\to x\ge 0$ we have $f_n(x_n)\to 0$.

\subsection*{Remarks}
\begin{itemize}
	\item The compactness of $[a,b]$ is crucial in (ii) to extract a convergent subsequence $(x_{n_k})$.
	\item The continuity of $f$ at the limit point is used twice: in Lemma~\ref{lem:move} and to conclude $f(x_k)\to f(x)$ in (ii).
	\item On compact domains, Remark~\ref{rem:equiv} provides a practical test for uniform convergence: \emph{check moving points}.
\end{itemize}

\section*{Problem 6 — Discontinuities under uniform limits}

Let $(f_n)$ be a sequence of real-valued functions on $[0,1]$ converging uniformly to a function $f$.
For each $n$, let $D_n$ be the set of discontinuities of $f_n$, and let $D$ be the set of discontinuities of $f$.

\subsection*{(a) Inclusion $D \subseteq \limsup D_n$}
We show that
\[
D \subseteq \bigcap_{n=1}^{\infty}\ \bigcup_{j=n}^{\infty} D_j
\qquad\text{(the limsup of the sets $D_j$).}
\]
Equivalently, we prove the contrapositive. Suppose $x\notin \limsup D_j$. Then there exists $N$ such that
$x\notin D_j$ for all $j\ge N$. Thus, for all $j\ge N$, the function $f_j$ is continuous at $x$.
Since $f_j \to f$ uniformly on $[0,1]$, the uniform-limit lemma implies that $f$ is continuous at $x$.
Hence $x\notin D$, proving the desired inclusion.

\subsection*{(b) A uniform bound on the number of discontinuities passes to the limit}
Assume there exists a finite $k$ such that each $f_n$ has at most $k$ discontinuities.
We claim that $f$ has at most $k$ discontinuities.

Suppose, for a contradiction, that $f$ has at least $k+1$ distinct discontinuities $x_1,\dots,x_{k+1}\in[0,1]$.
Let the oscillation of a function $g$ at $x$ be
\[
\omega_g(x)=\lim_{r\downarrow 0}\left(\sup_{|t-x|<r} g(t) - \inf_{|t-x|<r} g(t)\right).
\]
Then $g$ is continuous at $x$ iff $\omega_g(x)=0$. Moreover, if $\|g-h\|_\infty\le \varepsilon$, then
$\omega_g(x)\le\omega_h(x)+2\varepsilon$ and $\omega_h(x)\le\omega_g(x)+2\varepsilon$.

Since $x_i$ is a discontinuity of $f$, we have $\omega_f(x_i)>0$ for each $i$.
Choose $0<\varepsilon<\frac12\min_{1\le i\le k+1}\omega_f(x_i)$.
Because $f_n\to f$ uniformly, there exists $N$ such that $\|f_n-f\|_\infty<\varepsilon$ for all $n\ge N$.
For each $i$ and $n\ge N$,
\[
\omega_{f_n}(x_i)\ \ge\ \omega_f(x_i) - 2\varepsilon\ >\ 0,
\]
so $f_n$ is discontinuous at $x_i$. Hence $f_n$ has at least $k+1$ discontinuities for all $n\ge N$,
contradicting the hypothesis that each $f_n$ has at most $k$ discontinuities.
Therefore $f$ has at most $k$ discontinuities.

\subsection*{Remarks and example}

(1) From (a) we always have $D \subseteq \limsup D_n$, so if each $D_n$ is finite then $D$ is at most countable.

(2) The conclusion in (b) is stronger: if each $f_n$ has at most $k$ discontinuities for a fixed $k$, then $f$ has at most $k$ discontinuities.

(3) Example (sharpness for $k=1$).

Fix $c\in(0,1)$ and define $f=\mathbf{1}_{[c,1]}$ and $f_n=\mathbf{1}_{[c+1/n,\,1]}$.
Then $f_n\to f$ uniformly, each $f_n$ has exactly one discontinuity, and $f$ has exactly one discontinuity.

\paragraph{Quick example (sharpness).}
Let $k=1$. Fix a point $c\in(0,1)$ and define
\[
f(x) = \mathbf{1}_{[c,1]}(x) =
\begin{cases}
	0, & x < c,\\
	1, & x \ge c,
\end{cases}
\qquad\text{and}\qquad
f_n(x) = \mathbf{1}_{[\,c+1/n,\,1\,]}(x).
\]
Each $f_n$ has exactly one discontinuity (at $x=c+1/n$);
$f_n \to f$ uniformly on $[0,1]$,
and $f$ has exactly one discontinuity (at $x=c$).
This shows that the bound ``$\le k$'' is \emph{sharp}.
	
	
\section*{Problem 7 — Absolutely Convergent Coefficients and Trig Series}

\paragraph{Statement.}
Let $\sum_{m=1}^\infty a_m$ be absolutely convergent.
(a) Define $f_n(x)=\sum_{m=1}^n a_m\sin(mx)$ on $[-\pi,\pi]$. Show $f_n'(x)=\sum_{m=1}^n m a_m\cos(mx)$.
(b) Show $f(x)=\sum_{m=1}^\infty a_m\sin(mx)$ defines a continuous function on $[-\pi,\pi]$, but the series
$\sum_{m=1}^\infty m a_m\cos(mx)$ need not converge.

\paragraph{Solution.}
(a) Finite sums differentiate termwise:
\[
f_n'(x)=\sum_{m=1}^n a_m(\sin mx)' = \sum_{m=1}^n m a_m\cos(mx).
\]
(b) Since $|a_m\sin(mx)|\le|a_m|$ and $\sum|a_m|<\infty$, the M–test yields uniform convergence of
$\sum a_m\sin(mx)$ on $[-\pi,\pi]$, hence $f$ is continuous.
However, with $a_m=\dfrac{1}{m(\log(m+1))^2}$ we have $\sum|a_m|<\infty$, yet at $x=0$,
\[
\sum_{m=1}^\infty m a_m\cos(mx)=\sum_{m=1}^\infty \frac{1}{(\log(m+1))^2}=\infty,
\]
so $\sum m a_m\cos(mx)$ may diverge.

\bigskip
\section*{Problem 8 — Uniformity and Continuity of $\displaystyle f(x)=\sum_{m=1}^\infty (x-m)^{-2}$ on $X$}

\paragraph{Statement.}
Let $X=\mathbb R\setminus\{1,2,3,\dots\}$, $f_n(x)=\sum_{m=1}^n (x-m)^{-2}$ and $f(x)=\sum_{m=1}^\infty (x-m)^{-2}$.
Does $f_n\to f$ uniformly on $X$? Is $f$ continuous?

\paragraph{Theory.}
\begin{itemize}
	\item \textbf{Uniform convergence on all of $X$.}
	Uniform convergence on $X=\mathbb{R}\setminus\mathbb{N}$ would require
	\[
	\sup_{x\in X}\ \sum_{m>n} (x-m)^{-2}\ \longrightarrow\ 0\qquad (n\to\infty).
	\]
	This fails because near $x=n+1$ the tail contains $(x-(n+1))^{-2}$, which can be arbitrarily large.
	
	\item \textbf{Uniform convergence on compact subsets.}
	For any compact $K\Subset X$ there is a positive distance $\delta=\mathrm{dist}(K,\mathbb{N})>0$
	from $K$ to the pole set. Then for all $x\in K$ and all $m$,
	\[
	|x-m|\ge\delta \quad\text{except for finitely many $m$,}
	\]
	and the tail is dominated by a convergent $1/m^{2}$–type series. Hence by the Weierstrass M–test,
	$\displaystyle \sum_{m=1}^\infty (x-m)^{-2}$ converges uniformly on $K$.
\end{itemize}


\paragraph{Solution.}
Uniform convergence on $X$ fails: given $n$, as $x\to(n+1)$ within $X$,
the tail contains $(x-(n+1))^{-2}\to\infty$, so $\sup_{x\in X}\sum_{m>n}(x-m)^{-2}=\infty$.
However, for each compact $K\Subset X$ there is $\delta>0$ with $\operatorname{dist}(K,\mathbb N)\ge\delta$.
Then $|(x-m)^{-2}|\le (|m|-C)^{-2}$ for $m$ large and $x\in K$, and $\sum (|m|-C)^{-2}$ converges; thus the series
converges uniformly on $K$. Hence $f$ is continuous on $X$.

\bigskip
\section*{Problem 9 — Power Series with $\sum a_n$ Convergent}

\paragraph{Statement.}
Let $(a_n)$ be real with $\sum_{n=0}^\infty a_n$ convergent.
(a) Show $\sum a_n x^n$ converges for $x\in(-1,1)$ and $f(x)=\sum a_n x^n$ is differentiable on $(-1,1)$.
(b) Show that $f$ extends continuously to $(-1,1]$ with $f(1)=\sum a_n$, and that for each $r\in(-1,1)$ the series
converges uniformly on $[r,1]$. Must $f'(1)$ exist?

\paragraph{Theory.}
\begin{itemize}
	\item \textbf{Dirichlet/Abel summation.}
	With \(s_n=\sum_{k=0}^n a_k\) bounded (since \(\sum a_k\) converges) and \(x^n\) monotone with \(x^n\to0\) for \(|x|<1\),
	the series \(\displaystyle \sum_{n=0}^\infty a_n x^n\) converges.
	
	\item \textbf{Abel transform.} For any \(N\ge0\),
	\[
	\sum_{n=0}^{N} a_n x^n
	\;=\;
	(1-x)\sum_{n=0}^{N} s_n x^n \;+\; s_N x^{N+1}.
	\]
\end{itemize}


\paragraph{Solution.}
(a) Let $s_n=\sum_{k=0}^n a_k$. Since $\sum a_n$ converges, $(s_n)$ is bounded. For $|x|<1$, $x^n\to0$ and is monotone in $n$ on $[0,1)$, hence by Dirichlet,
$\sum a_n x^n$ converges. Using Abel summation
\[
\sum_{n=0}^N a_n x^n=(1-x)\sum_{n=0}^N s_n x^n+s_N x^{N+1},
\]
one checks uniform convergence (and of the differentiated series) on $[-r,r]$, $0<r<1$, so $f$ is differentiable and
\[
f'(x)=\sum_{n=1}^\infty n a_n x^{n-1}\qquad (|x|<1).
\]
(b) From the Abel identity, as $x\to1^-$,
\(
f(x)=(1-x)\sum_{n=0}^\infty s_n x^n \to 0
\),
so $f(1)=\sum a_n$ and $f$ is continuous at $1$. For fixed $r<1$, uniform convergence on $[r,1]$ follows from Dirichlet’s test with uniformly bounded $(s_n)$ and monotone $x^n$.
In general $f'(1)$ \emph{need not} exist; e.g. $a_n=(-1)^n/(n+1)$ gives $\sum a_n$ convergent but $\sum n a_n$ divergent, so $f'(1^-)$ does not exist.

\bigskip
\section*{Problem 10 — Uniform Convergence on $(-1,1)$ with Radius $1$}

\paragraph{Question.}
Is there a real power series with radius of convergence $1$ that converges uniformly on $(-1,1)$?

\paragraph{Answer.}
Yes. For example,
\[
\sum_{n=1}^\infty \frac{x^n}{n^2}
\]
has radius $R=\big(\limsup (n^{-2})^{1/n}\big)^{-1}=1$, and by the M–test (since $\sum 1/n^2<\infty$ and $|x|^n\le1$ on $(-1,1)$) the series converges uniformly on $(-1,1)$.

% ===================== Abel–Dirichlet Cheat-Sheet =====================
\paragraph{Abel partial summation (discrete integration by parts).}
Let $(a_n)$ and $(b_n)$ be sequences, and put $S_n=\sum_{k=0}^n a_k$ with $S_{-1}:=0$.
For any integers $0\le m\le N$,
\[
\sum_{n=m}^N a_n b_n
= S_N b_N - S_{m-1} b_m \;-\; \sum_{n=m}^{N-1} S_n\,(b_{n+1}-b_n).
\]
\emph{Power series version (Abel transform).} For $x\in\mathbb{R}$ and $N\ge0$,
\[
\sum_{n=0}^N a_n x^n
= (1-x)\sum_{n=0}^N S_n x^n + S_N x^{N+1}.
\tag{A}
\]

\paragraph{Dirichlet test (series).}
If $(S_n)$ is bounded and $(b_n)$ is monotone with $b_n\to 0$, then $\sum a_n b_n$ converges.

\paragraph{Uniform Dirichlet on intervals.}
If $\sup_n |S_n|\le M$ and for each $x$ in a set $E$ the sequence $b_n(x)$ is monotone in $n$ with
$\sup_{x\in E} |b_n(x)|\to 0$, then $\sum a_n b_n(x)$ converges uniformly on $E$.
In particular, for fixed $r\in(0,1)$, the series $\sum a_n x^n$ converges uniformly on $[0,r]$ and on $[r,1)$.

\paragraph{Abel’s theorem (boundary continuity for power series).}
If $\sum_{n=0}^\infty a_n$ converges (so $S_n\to S$) and $f(x)=\sum_{n=0}^\infty a_n x^n$ for $x\in[0,1)$,
then $f$ extends continuously to $x=1$ with
\[
\lim_{x\to 1^-} f(x)=\sum_{n=0}^\infty a_n.
\]
\emph{Proof sketch via (A):} write $f(x)=(1-x)\sum_{n\ge0} S_n x^n$ and compare with the geometric kernel
$\sum_{n\ge0} x^n=(1-x)^{-1}$; boundedness of $(S_n)$ yields the limit $S$.

\paragraph{Differentiation inside $\sum a_n x^n$ on $(-1,1)$.}
If $\sum a_n$ converges then for each $r\in(0,1)$:
\[
\sum_{n=0}^\infty a_n x^n \text{ and } \sum_{n=1}^\infty n a_n x^{n-1}
\quad\text{converge uniformly on } [-r,r],
\]
hence $f(x)=\sum a_n x^n$ is $C^1$ on $(-1,1)$ with $f'(x)=\sum n a_n x^{n-1}$.
\emph{Reason:} apply (A) to control tails by bounded $S_n$ and the factors $(1-x)$, together with Dirichlet’s test
for $\sum S_n [(n+1)x^n-nx^{n-1}]$ on compact subintervals of $(-1,1)$.

\paragraph{Quick checklist.}
\begin{itemize}
	\item Use (A) to pass $x\to 1^-$ or to prove uniform convergence on $[r,1]$.
	\item If $(S_n)$ bounded and $x^n$ monotone $\downarrow 0$, then $\sum a_n x^n$ converges (Dirichlet).
	\item For termwise differentiation, restrict to $|x|\le r<1$ and use uniform Dirichlet.
	\item Need convergence of $\sum n|a_n|$ only if you want uniform convergence of $f'$ on all of $(-1,1)$.
\end{itemize}
% =================== End Abel–Dirichlet Cheat-Sheet ====================

% Preamble:
% \usepackage[most]{tcolorbox}

\begin{tcolorbox}[title=Abel–Dirichlet Cheat-Sheet,colback=white,colframe=black,
	fonttitle=\bfseries,boxsep=4pt,arc=2pt]
	
	\paragraph{Abel partial summation (discrete IBP).}
	For sequences $(a_n),(b_n)$ and $S_n=\sum_{k=0}^n a_k$ (with $S_{-1}:=0$), for $0\le m\le N$,
	\[
	\sum_{n=m}^N a_n b_n
	= S_N b_N - S_{m-1} b_m - \sum_{n=m}^{N-1} S_n\,(b_{n+1}-b_n).
	\]
	\textit{Power–series form (Abel transform).} For $x\in\mathbb R$,
	\[
	\sum_{n=0}^N a_n x^n=(1-x)\sum_{n=0}^N S_n x^n + S_N x^{N+1}. \tag{A}
	\]
	
	\paragraph{Dirichlet test (series).}
	If $(S_n)$ is bounded and $(b_n)$ is monotone with $b_n\!\to\!0$, then $\sum a_n b_n$ converges.\\
	\textit{One-line proof via Abel:} apply the identity above with $m=0$, $B_N:=b_N\!\to\!0$ and
	$\sum S_n(b_{n+1}-b_n)$ absolutely bounded by $\sup|S_n|\,\sum|b_{n+1}-b_n| \le \sup|S_n|\,(|b_0|+|B_N|)$.
	
	\paragraph{Uniform Dirichlet (on sets).}
	If $\sup_n|S_n|\le M$ and for each $x\in E$ the sequence $b_n(x)$ is monotone in $n$ with
	$\sup_{x\in E}|b_n(x)|\to0$, then $\sum a_n b_n(x)$ converges \emph{uniformly} on $E$.
	In particular, for any fixed $r\in(0,1)$, $\sum a_n x^n$ converges uniformly on $[0,r]$ and on $[r,1)$.
	
	\paragraph{Abel’s theorem (boundary continuity of power series).}
	If $\sum_{n\ge0} a_n$ converges and $f(x)=\sum_{n\ge0} a_n x^n$ on $[0,1)$, then
	\[
	\lim_{x\to1^-} f(x)=\sum_{n=0}^\infty a_n.
	\]
	\textit{Sketch via (A):} $f(x)=(1-x)\sum S_n x^n$, with bounded $(S_n)$ and the kernel $(1-x)\sum x^n=1$.
	
	\paragraph{Differentiation inside $\textstyle\sum a_n x^n$ on $(-1,1)$.}
	If $\sum a_n$ converges, then for each $r\in(0,1)$ the series
	$\sum a_n x^n$ and $\sum n a_n x^{n-1}$ converge uniformly on $[-r,r]$.
	Hence $f(x)=\sum a_n x^n$ is $C^1$ on $(-1,1)$ with $f'(x)=\sum n a_n x^{n-1}$.
	
	\paragraph{Quick checklist.}
	Use (A) near $x=1$; Dirichlet for $|x|<1$; uniform Dirichlet on compact subintervals; require $\sum n|a_n|<\infty$
	only for \emph{uniform} convergence of $f'$ on all $(-1,1)$.
\end{tcolorbox}
	
	
\section*{Problem 11 — Uniform Continuity and Boundedness on $[0,\infty)$}

\paragraph{Statement.}
Which of the following $f:[0,\infty)\to\mathbb R$ are (a) uniformly continuous; (b) bounded?
\begin{enumerate}[label=(\roman*)]
	\item $f(x)=\sin(x^2)$,
	\item $f(x)=\inf\{|x-n^2|:n\in\mathbb N\}$,
	\item $f(x)=\dfrac{\sin(x^3)}{x+1}$.
\end{enumerate}

\paragraph{Solution.}
\begin{enumerate}[label=(\roman*)]
	\item \textbf{Bounded:} yes, $|\sin|\le1$. \textbf{Not uniformly continuous.}
	Let $a_k=\sqrt{k\pi}$, $b_k=\sqrt{k\pi+\pi/2}$. Then $|a_k-b_k|\to0$ but
	$|f(a_k)-f(b_k)|=1$.
	\item \textbf{Uniformly continuous:} yes. $f(x)=d(x,\{n^2\})$ is a distance to a closed set,
	hence $|f(x)-f(y)|\le|x-y|$ (1-Lipschitz).
	\textbf{Unbounded:} at $x=\tfrac{n^2+(n+1)^2}{2}$ we have $f(x)=\tfrac{2n+1}{2}\to\infty$.
	\item \textbf{Bounded:} $|f(x)|\le 1/(x+1)\le1$.
	\textbf{Uniformly continuous:} yes. Given $\varepsilon>0$, choose $R$ with $2/(R+1)<\varepsilon/2$.
	If $x,y\ge R$ then $|f(x)-f(y)|\le 2/(R+1)<\varepsilon/2$.
	On $[0,R+1]$, $f$ is continuous on a compact set, hence uniformly continuous;
	choose $\delta\le1$ to combine the cases.
\end{enumerate}

% (i) sin(x^2)
\begin{figure}[h]
	\centering
	\includegraphics[width=\textwidth]{figures/MATH-ALLS-2_sinx2_global.png}
	\caption{$f(x)=\sin(x^2)$ on $[-10,10]$.}
\end{figure}

\begin{figure}[h]
	\centering
	\includegraphics[width=\textwidth]{figures/MATH-ALLS-2_sinx2_zoom.png}
	\caption{Zoom of $f(x)=\sin(x^2)$ on $[-3,3]$.}
\end{figure}

% (ii) distance to squares
\begin{figure}[h]
	\centering
	\includegraphics[width=\textwidth]{figures/MATH-ALLS-2_dist_n2_global.png}
	\caption{$f(x)=\inf\{|x-n^2|:n\in\mathbb N\}$ on $[0,50]$. Piecewise-linear “sawtooth” with zeros at perfect squares.}
\end{figure}

\begin{figure}[h]
	\centering
	\includegraphics[width=\textwidth]{figures/MATH-ALLS-2_dist_n2_zoom.png}
	\caption{Zoom on $[0,20]$ highlighting minima at $n^2$ and local maxima midway between consecutive squares.}
\end{figure}

% (iii) sin(x^3)/(x+1)
\begin{figure}[h]
	\centering
	\includegraphics[width=\textwidth]{figures/MATH-ALLS-2_sinx3_over_xp1_global.png}
	\caption{$f(x)=\sin(x^3)/(x+1)$ on $[-1,20]$ (note the removable pole at $x=-1$).}
\end{figure}

\begin{figure}[h]
	\centering
	\includegraphics[width=\textwidth]{figures/MATH-ALLS-2_sinx3_over_xp1_tail.png}
	\caption{Tail behaviour of $f(x)=\sin(x^3)/(x+1)$: oscillations with amplitude $\sim 1/x$ decay.}
\end{figure}


\bigskip
\section*{Problem 12 — Uniform Limit of Uniformly Continuous Functions}

\paragraph{Statement.}
If $(f_n)$ are uniformly continuous on $\mathbb R$ and $f_n\to f$ uniformly, show $f$ is uniformly continuous.
Give $(f_n)$ uniformly continuous with pointwise $f_n\to f$ where $f$ is continuous but not uniformly continuous.

\paragraph{Proof.}
Fix $\varepsilon>0$. Take $N$ with $\|f_N-f\|_\infty<\varepsilon/3$ and $\delta>0$ for uniform continuity of $f_N$
so that $|x-y|<\delta\Rightarrow |f_N(x)-f_N(y)|<\varepsilon/3$.
Then for $|x-y|<\delta$,
\[
|f(x)-f(y)|\le |f(x)-f_N(x)|+|f_N(x)-f_N(y)|+|f_N(y)-f(x)|<\varepsilon.
\]

\paragraph{Example.}
Let $\mathrm{clip}_n(x)=\max(-n,\min(x,n))$ and
\[
f_n(x)=\sin\!\big(\mathrm{clip}_n(x)^2\big).
\]
Then $f_n$ is globally Lipschitz (hence uniformly continuous), and for each fixed $x$, if $n>|x|$ then $f_n(x)=\sin(x^2)$.
Thus $f_n\to f$ pointwise with $f(x)=\sin(x^2)$, which is continuous but not uniformly continuous on $\mathbb R$.

\begin{center}
	\includegraphics[width=0.86\textwidth]{figures/MATH-ALLS-2_clipn_sin2_plot.png}
\end{center}
\noindent
\textit{Figure.} The clipped functions $f_n(x)=\sin(\mathrm{clip}_n(x)^2)$.
Each $f_n$ agrees with $\sin(x^2)$ on $[-n,n]$ and becomes constant outside, so $f_n$ are uniformly continuous
yet converge pointwise to the non–uniformly continuous limit $\sin(x^2)$.


\bigskip
\section*{Problem 13 — Limit at Infinity Implies Uniform Continuity on $[0,\infty)$}

\paragraph{Statement.}
Let $f:[0,\infty)\to\mathbb R$ be continuous and suppose $\lim_{x\to\infty} f(x)=L\in\mathbb R$.
Must $f$ be uniformly continuous on $[0,\infty)$?

\paragraph{Proof.}
\paragraph{Answer.} \textbf{Yes, it must.}
Fix $\varepsilon>0$. Choose $R>0$ so that $\lvert f(x)-L\rvert<\varepsilon/3$ for all $x\ge R$.
Continuity on the compact $[0,R+1]$ gives $\delta_0>0$ with
\[
\lvert x-y\rvert<\delta_0 \;\Rightarrow\; \lvert f(x)-f(y)\rvert<\varepsilon
\quad\text{for } x,y\in[0,R+1].
\]
Set $\delta:=\min\{\delta_0,1\}$. For $\lvert x-y\rvert<\delta$:
\begin{itemize}
	\item if $x,y\le R+1$: done by compact uniform continuity on $[0,R+1]$;
	\item if $x,y\ge R$: then
	\[
	\lvert f(x)-f(y)\rvert
	\le \lvert f(x)-L\rvert+\lvert f(y)-L\rvert
	< \tfrac{\varepsilon}{3}+\tfrac{\varepsilon}{3}
	< \varepsilon;
	\]
	\item the “mixed” case cannot occur, because $\lvert x-y\rvert<1$ forbids one point $\ge R+1$
	and the other $<R$.
\end{itemize}
Thus $f$ is uniformly continuous on $[0,\infty)$.

	
\section*{Problem 14 — Bounded Derivative and a Classical Oscillating Example}

\paragraph{(a) Bounded derivative $\Rightarrow$ uniform continuity.}
If $|f'(x)|\le M$ on $\mathbb R$, then by the Mean Value Theorem
\[
|f(x)-f(y)|=|f'(\xi)||x-y|\le M|x-y|,
\]
so $f$ is Lipschitz, hence uniformly continuous.

\paragraph{(b) $g(x)=x^2\sin(1/x^2)$ with $g(0)=0$.}
For $x\ne 0$,
\[
g'(x)=2x\sin(1/x^2)-\frac{2\cos(1/x^2)}{x},\qquad
g'(0)=\lim_{x\to0}\frac{x^2\sin(1/x^2)}{x}=0.
\]
Thus $g$ is differentiable, but $g'$ is unbounded near $0$.
Since $g$ is continuous on the compact interval $[-1,1]$, it is uniformly continuous there.


\begin{figure}[h]
	\centering
	\includegraphics[width=\textwidth]{figures/MATH-ALLS-2_g_global.png}
	\caption{Global view of $g(x)=x^2\sin(1/x^2)$ with $g(0)=0$. Oscillations concentrate near $0$ but the amplitude $x^2$ tends to $0$.}
\end{figure}

\begin{figure}[h]
	\centering
	\includegraphics[width=\textwidth]{figures/MATH-ALLS-2_g_zoom.png}
	\caption{Zoom near $x=0$: $g(x)$ oscillates with shrinking amplitude and increasing frequency, consistent with $g'(0)=0$.}
\end{figure}

\begin{figure}[h]
	\centering
	\includegraphics[width=\textwidth]{figures/MATH-ALLS-2_gprime_zoom.png}
	\caption{Derivative $g'(x)=2x\sin(1/x^2)-2\cos(1/x^2)/x$ near $0$ (vertical scale clipped to show oscillations). The $-2\cos(1/x^2)/x$ term causes blow-up.}
\end{figure}

\begin{figure}[h]
	\centering
	\includegraphics[width=\textwidth]{figures/MATH-ALLS-2_gprime_semilog.png}
	\caption{Semilog plot of $|g'(x)|$ on $(0,0.2]$: the growth as $x\to0^+$ is unbounded, proving $g'$ is not bounded near $0$.}
\end{figure}


\bigskip
\section*{Problem 15 — Approximations of a Riemann–Integrable Function on $[0,1]$}

Let $f$ be bounded and Riemann integrable on $[0,1]$.
\begin{enumerate}[label=(\alph*)]
	\item \textbf{No.} A uniform limit of continuous functions is continuous; a Riemann–integrable $f$ can be discontinuous.
	For example, $f=\mathbf 1_C$ (indicator of the Cantor set) is Riemann integrable with $\int f=0$ but is not continuous on $C$.
	\item \textbf{Yes.} Since $f\in\mathcal R[0,1]$, upper/lower sums approximate $f$ in $L^1$. Step functions are $L^1$–approximable by continuous functions (smooth their jumps in small neighborhoods). Hence there exist continuous $f_n$ with $\int_0^1|f_n-f|\to0$.
	\item \textbf{Yes.} From (b) approximate $f$ in $L^1$ by a continuous $g$. By Weierstrass, there exist polynomials $p_n$ with $\|p_n-g\|_\infty\to0$, so $\int_0^1|p_n-f|\le\|p_n-g\|_\infty+\int_0^1|g-f|\to0$.
\end{enumerate}

\bigskip
\section*{Problem 16$^\ast$ — A Nowhere Differentiable Continuous Function}

Let $\varphi(x)=|x|$ for $x\in[-1,1]$ and extend to $\mathbb R$ by $\varphi(x+2)=\varphi(x)$.
\begin{enumerate}[label=(\roman*)]
	\item \textbf{$\varphi$ is 1-Lipschitz.} On each interval between consecutive integers $\varphi$ is linear with slope $\pm1$, hence
	$|\varphi(s)-\varphi(t)|\le|s-t|$ for all $s,t$.
	\item Define $f(x)=\sum_{n=0}^\infty (3/4)^n\varphi(4^n x)$. Since $0\le\varphi\le1$, the Weierstrass M–test shows uniform convergence, so $f$ is continuous.
	\item Fix $x\in\mathbb R$ and $m\in\mathbb N$. Put $\delta_m=\pm \tfrac12 4^{-m}$ with the sign so that no integer lies between $4^m x$ and $4^m(x+\delta_m)$. Then
	\[
	\big|\varphi(4^m(x+\delta_m))-\varphi(4^m x)\big|=|4^m\delta_m|=\tfrac12,
	\]
	while for $n<m$,
	\[
	\big|\varphi(4^n(x+\delta_m))-\varphi(4^n x)\big|\le |4^n\delta_m|=\tfrac12\,4^{\,n-m},
	\]
	and for $n\ge m+1$ we have $4^n\delta_m\in 2\mathbb Z$ so by $2$–periodicity the difference is $0$.
	Therefore,
	\[
	\begin{aligned}
		\big|f(x+\delta_m)-f(x)\big|
		&\ge \Big(\frac{3}{4}\Big)^m\frac12
		- \sum_{n=0}^{m-1}\Big(\frac{3}{4}\Big)^n \frac12\,4^{\,n-m} \\
		&= \frac{1}{2}\,4^{-m}\!\left(3^m-\sum_{n=0}^{m-1}3^n\right)
		= \frac{1}{2}\,4^{-m}\cdot\frac{3^m+1}{2}.
	\end{aligned}
	\]
	Dividing by $|\delta_m|=\tfrac12\,4^{-m}$ yields
	\[
	\left|\frac{f(x+\delta_m)-f(x)}{\delta_m}\right|\ge \frac12\,(3^m+1).
	\]
	Letting $m\to\infty$ shows $f$ is not differentiable at $x$. Since $x$ was arbitrary, $f$ is continuous and nowhere differentiable.
\end{enumerate}

\begin{center}
	\includegraphics[width=0.8\textwidth]{figures/MATH-ALLS-2_phi_periodic_plot.png}
\end{center}
\noindent
\textit{Figure.} The triangular wave $\varphi(x)=|x|$ periodically extended by
$\varphi(x+2)=\varphi(x)$.  This Lipschitz function forms the building block in the
definition of the Weierstrass–type nowhere differentiable function
$f(x)=\sum_{n=0}^\infty \big(\tfrac34\big)^n \varphi(4^n x)$.

\begin{figure}[h]
	\centering
	\includegraphics[width=\textwidth]{figures/MATH-ALLS-2_weierstrass_partial_sums.png}
	\caption{Partial sums $f_N(x)=\sum_{n=0}^{N-1}\big(\tfrac{3}{4}\big)^n \varphi(4^n x)$ for several $N$.
		Increasing $N$ adds finer triangular features, suggesting fractal-like roughness of the limit.}
\end{figure}

\begin{figure}[h]
	\centering
	\includegraphics[width=\textwidth]{figures/MATH-ALLS-2_weierstrass_zoom.png}
	\caption{Zoom of $f_{10}(x)$ near $x=0$, illustrating rapid oscillations that preclude differentiability.}
\end{figure}
	
	\section*{Problem 1: Quickies on norms}
	
	\paragraph{(a) Ball inclusion and Lipschitz equivalence.}
	Two norms $\|\cdot\|,\|\cdot\|'$ on $V$ are Lipschitz equivalent iff there exist $r,R>0$ with
	$B_r\subseteq B_1'\subseteq B_R$, where $B_\rho=\{x:\|x\|<\rho\}$, $B'_\rho=\{x:\|x\|'<\rho\}$.
	Indeed, $c\|x\|\le \|x\|'\le C\|x\|$ implies the inclusions with $r=1/C$, $R=1/c$;
	conversely the inclusions give $\|x\|'\le (1/r)\|x\|$ and $\|x\|\le R\|x\|'$.
	
	\paragraph{(b) Same convergent sequences.}
	The bounds $c\|x\|\le \|x\|'\le C\|x\|$ imply
	$\|x_n-x\|\to0 \Leftrightarrow \|x_n-x\|'\to0$. Conversely, if this fails, one can find
	$\|x_k\|=1$ with $\|x_k\|'\to\infty$; then $y_k=x_k/k$ satisfies $y_k\to0$ in $\|\cdot\|$
	but not in $\|\cdot\|'$.
	
	\paragraph{(c) The norm $\|x\|_\varphi=\|x\|+|\varphi(x)|$.}
	It is a norm. If $\varphi$ is not continuous, there are $x_k$ with $\|x_k\|=1$ and
	$|\varphi(x_k)|\to\infty$, hence $\|x_k\|_\varphi\to\infty$, so no inequality
	$\|x\|_\varphi\le C\|x\|$ holds; thus $\|\cdot\|_\varphi$ is not Lipschitz equivalent to $\|\cdot\|$.
	
	\paragraph{(d$^\ast$) Finite dimensionality.}
	If every norm on $V$ is Lipschitz equivalent to $\|\cdot\|$, then $V$ must be finite-dimensional.
	Otherwise there exists a discontinuous linear functional $\varphi$ on $V$, and by (c) the norm
	$\|x\|+|\varphi(x)|$ would not be equivalent to $\|\cdot\|$.
	
	\bigskip
	\section*{Problem 2: Integrals of vector-valued maps}
	
	
	
	Let $f:[0,1]\to\mathbb R^n$ be continuous and $\displaystyle \int_0^1 f:=\big(\int_0^1 f_1,\dots,\int_0^1 f_n\big)$.
	
	\paragraph{(a)}
	With $v=\int_0^1 f$ and the Euclidean dot product,
	\[
	\|v\|_2^2=v\cdot v=v\cdot \int_0^1 f=\int_0^1 v\cdot f(x)\,dx
	\le \int_0^1 \|v\|_2\,\|f(x)\|_2\,dx
	=\|v\|_2\int_0^1 \|f(x)\|_2\,dx,
	\]
	hence $\Big\|\int_0^1 f\Big\|_2 \le \int_0^1 \|f(x)\|_2\,dx$.
	
% ---------- Problem 2(b) ----------
\paragraph{(b) Equality for every norm.}
Equality
\[
\Bigl\|\!\int_{0}^{1} f \Bigr\| \;=\; \int_{0}^{1} \|f\|
\]
holding for \emph{every} norm on $\mathbb{R}^n$ occurs precisely when there exist
\[
u\in\mathbb{R}^n\setminus\{0\}
\quad\text{and}\quad
\alpha:[0,1]\to\mathbb{R}\ \text{of constant sign}
\]
such that
\[
f(x)=\alpha(x)\,u \qquad (x\in[0,1]).
\]

\noindent\textit{Sufficiency.}
If $f(x)=\alpha(x)u$ and $\alpha$ has constant sign, then for any norm $\|\cdot\|$,
\[
\Bigl\|\!\int_{0}^{1} f \Bigr\|
=\bigl\| (\int_{0}^{1}\alpha)\,u \bigr\|
= \bigl|\!\int_{0}^{1}\alpha \bigr|\,\|u\|
= \int_{0}^{1} |\alpha|\,\|u\|
= \int_{0}^{1} \|f\|.
\]

\noindent\textit{Necessity.}
Fix a strictly convex norm (e.g.\ Euclidean). Equality in the triangle inequality forces
$f(x)$ to be positively colinear a.e.; hence $f(x)=\alpha(x)u$ for some $u$ and scalar $\alpha$.
Requiring equality for \emph{every} norm rules out sign changes of $\alpha$, so $\alpha$ has
constant sign on $[0,1]$.
\medskip

% ---------- Problem 3 ----------
\section*{Problem 3: Continuity of the pointwise scalar product}



\paragraph{Problem 3.}

Let $(x^{(m)})$ and $(y^{(m)})$ be sequences in $\mathbb{R}^n$ converging to $x$ and $y$ respectively.
Show that $x^{(m)}\cdot y^{(m)}$ converges to $x\cdot y$. Deduce that if
$f:\mathbb{R}^n\to\mathbb{R}^p$ and $g:\mathbb{R}^n\to\mathbb{R}^p$ are continuous at
$x\in\mathbb{R}^n$, then so is the pointwise scalar product
$f\cdot g:\mathbb{R}^n\to\mathbb{R}$.

\paragraph{Solution}

Let $x^{(m)}\to x$ and $y^{(m)}\to y$ in $\mathbb{R}^n$. Then
\begin{align*}
	x^{(m)}\!\cdot y^{(m)} - x\!\cdot y
	&= (x^{(m)}-x)\!\cdot y^{(m)} \;+\; x\!\cdot (\,y^{(m)}-y\,) \\
	&\xrightarrow[m\to\infty]{} 0,
\end{align*}
by Cauchy--Schwarz and boundedness near the limits. Hence the dot product is continuous.

If $f,g:\mathbb{R}^n\to\mathbb{R}^p$ are continuous at $x$, then the pointwise scalar product
\[
z \;\longmapsto\; f(z)\cdot g(z)
\]
is continuous at $x$ as a composition of continuous maps.
	


%====================  Problem 4  ====================%
\section*{Problem 4. {L1} on $C([0,1])$ and on bounded Riemann–integrable functions}

\paragraph{Problem (exact).}
\begin{enumerate}[(a)]
	\item Show that $\|f\|_{1}=\int_{0}^{1}|f|$ defines a norm on the vector space $C([0,1])$.
	Is it Lipschitz equivalent to the uniform norm? Is $C([0,1])$ with norm $\|\cdot\|_{1}$ complete?
	
	\item Let $R([0,1])$ denote the vector space of all bounded Riemann integrable functions on $[0,1]$.
	Does $\|f\|_{1}=\int_{0}^{1}|f|$ define a norm on $R([0,1])$? If so, is $R([0,1])$ complete with this norm?
	What if we replace $\|\cdot\|_{1}$ with $\|f\|_{\infty}=\sup\{|f(x)|:x\in[0,1]\}$?
\end{enumerate}

\noindent\rule{\linewidth}{0.6pt}

\paragraph{Theory you’ll use.}
\begin{itemize}
	\item On a finite–measure set (here length $1$): $\displaystyle \|f\|_{1}\le \|f\|_{\infty}$.
	\item A norm must satisfy positivity, homogeneity, triangle inequality, and \emph{definiteness}
	($\|f\|=0 \Rightarrow f=0$ in the underlying space).
	\item Uniform limits of bounded Riemann–integrable functions are Riemann–integrable (boundedness preserved).
	\item $C([0,1])$ is dense in $L^{1}([0,1])$; the $\|\cdot\|_{1}$–completion of $C([0,1])$ is $L^{1}$.
\end{itemize}

\noindent\rule{\linewidth}{0.6pt}

\paragraph{Solutions.}

\subsection*{(a) On $C([0,1])$}
\begin{enumerate}[1)]
	\item \textbf{$\|\cdot\|_{1}$ is a norm on $C([0,1])$.} For $f,g\in C([0,1])$ and $\alpha\in\mathbb{R}$,
	\[
	\|f\|_{1}=\int_{0}^{1}|f|\ge 0,\qquad
	\|\alpha f\|_{1}=|\alpha|\,\|f\|_{1},\qquad
	\|f+g\|_{1}\le \|f\|_{1}+\|g\|_{1}.
	\]
	If $\|f\|_{1}=0$, then $|f|=0$ a.e.; by continuity, $f\equiv 0$. Hence $\|\cdot\|_{1}$ is a norm.
	
	\item \textbf{Not Lipschitz equivalent to $\|\cdot\|_{\infty}$.}
	Always $\|f\|_{1}\le \|f\|_{\infty}$. A reverse uniform bound
	$c\|f\|_{\infty}\le \|f\|_{1}$ fails: take continuous “spikes”
	\[
	f_n(x)=\begin{cases}
		1-\dfrac{|x-a_n|}{\ell_n}, & |x-a_n|\le \ell_n,\\[4pt]
		0, & \text{otherwise},
	\end{cases}
	\quad a_n\to 0,\ \ell_n\downarrow 0.
	\]
	Then $\|f_n\|_{\infty}=1$ but $\|f_n\|_{1}=\int_{a_n-\ell_n}^{a_n+\ell_n}\!\!\big(1-\frac{|x-a_n|}{\ell_n}\big)\,dx
	=\ell_n\to 0$. No $c>0$ works, so the norms are not Lipschitz equivalent.
	
	\item \textbf{Not complete in $\|\cdot\|_{1}$.}
	Let $h=\mathbf{1}_{[0,1/2]}$ (discontinuous). Choose $h_k\in C([0,1])$ that converge to $h$ in $L^{1}$,
	e.g.\ smooth ramps replacing the jump over width $1/k$. Then $\|h_k-h\|_{1}\to 0$, so $(h_k)$ is Cauchy in
	$\|\cdot\|_{1}$ but its limit is not continuous. Thus $\big(C([0,1]),\|\cdot\|_{1}\big)$ is not complete.
\end{enumerate}

\subsection*{(b) On $R([0,1])$ (bounded Riemann–integrable functions)}
\begin{enumerate}[1)]
	\item \textbf{$\|\cdot\|_{1}$ is not a norm (only a seminorm).}
	Definiteness fails: let
	\[
	f(x)=\begin{cases}
		1, & x=0,\\
		0, & x\in (0,1],
	\end{cases}
	\]
	which is bounded and Riemann–integrable with $\int_{0}^{1}|f|=0$ but $f\not\equiv 0$.
	Hence $\|f\|_{1}=0\not\Rightarrow f=0$ on $R([0,1])$.
	
	\item \textbf{Completeness with $\|\cdot\|_{1}$.}
	Since $\|\cdot\|_{1}$ is not a norm on the pointwise space $R([0,1])$, completeness as a normed space
	is not applicable. Passing to equivalence classes a.e.\ gives $L^{1}$, which is complete.
	
	\item \textbf{With the sup norm $\|\cdot\|_{\infty}$.}
	$\|\cdot\|_{\infty}$ is a norm on $R([0,1])$. If $(f_n)$ is $\|\cdot\|_{\infty}$–Cauchy, then
	$f_n\to f$ uniformly for some bounded $f$, and uniform limits of Riemann–integrable functions are
	Riemann–integrable; thus $f\in R([0,1])$. Therefore $R([0,1])$ is closed in the complete space of bounded
	functions and is itself complete in $\|\cdot\|_{\infty}$.
\end{enumerate}
%======================================================



\section*{Problem 5: Norms on $C^1([0,1])$ and $C^1_0([0,1])$}


\paragraph{Problem (exact).}
\begin{enumerate}[(a)]
	\item Let $C^{1}([0,1])$ be the vector space of real continuous functions on $[0,1]$ with
	continuous first derivatives. Define functions $\alpha,\beta,\gamma,\delta:C^{1}([0,1])\to\mathbb{R}$ by
	\[
	\alpha(f)=\sup_{x\in[0,1]}|f(x)|+\sup_{x\in[0,1]}|f'(x)|,\qquad
	\beta(f)=\sup_{x\in[0,1]} \big(|f(x)|+|f'(x)|\big),
	\]
	\[
	\gamma(f)=\sup_{x\in[0,1]}|f(x)|,\qquad
	\delta(f)=\sup_{x\in[0,1]}|f'(x)|.
	\]
	Which of these define norms on $C^{1}([0,1])$?
	Out of those that define norms, which pairs are Lipschitz equivalent?
	
	\item Let $C^{1}_{0}([0,1])$ be the set of functions $f\in C^{1}([0,1])$ such that
	$f(x)=0$ for $x$ in some neighborhood of the end points $0$ and $1$. Verify that
	$C^{1}_{0}([0,1])$ is a vector space. How would your answers in (a) change if we replace
	$C^{1}([0,1])$ by $C^{1}_{0}([0,1])$ (equivalently $C^{1}_{c}((0,1))$)?
\end{enumerate}

\noindent\rule{\linewidth}{0.6pt}

\paragraph{Theory you’ll use.}
\begin{itemize}
	\item A norm requires positivity, absolute homogeneity, triangle inequality, and
	\emph{definiteness} ($\|f\|=0\Rightarrow f=0$ in the underlying space).
	\item For any functions $A(x),B(x)$:
	$\sup(A+B)\le \sup A+\sup B$ and $\sup A\le \sup(A+B)$, $\sup B\le \sup(A+B)$.
	\item Mean–Value Theorem: $|f(x)-f(t)|\le |x-t|\;\sup|f'|$.
	\item Lipschitz equivalence of norms $\|\cdot\|_a,\|\cdot\|_b$:
	$\exists\,c,C>0$ with $c\|f\|_a\le \|f\|_b\le C\|f\|_a$ for all $f$.
\end{itemize}

\paragraph{(a) $C^1([0,1])$.}
Define
\[
\alpha(f)=\|f\|_\infty+\|f'\|_\infty,\qquad
\beta(f)=\big\|\,|f|+|f'|\,\big\|_\infty,\qquad
\gamma(f)=\|f\|_\infty,\qquad
\delta(f)=\|f'\|_\infty.
\]
Then $\alpha,\beta,\gamma$ are norms on $C^1([0,1])$, while $\delta$ is not (constants have
$\delta=0$ but are nonzero). Moreover,
\[
\beta(f)\le \alpha(f)\le 2\beta(f),
\]
so $\alpha\sim\beta$. The uniform norm $\gamma$ is not equivalent to $\alpha$ or $\beta$:
take smooth spikes $f_n$ with $\|f_n\|_\infty=1$ and $\|f'_n\|_\infty\to\infty$, giving
$\gamma(f_n)=1$ yet $\alpha(f_n)\to\infty$.

\paragraph{(b) $C^1_0([0,1])$.}
The set $C^1_0([0,1])=\{f\in C^1([0,1]): f=0\ \text{on neighborhoods of }0,1\}$ is a vector space.
Now $\delta$ is a norm: if $\|f'\|_\infty=0$, then $f$ is constant; since $f=0$ near an endpoint,
$f\equiv0$. For $f\in C^1_0$ we have $\|f\|_\infty\le \|f'\|_\infty$ (choose $t$ with $f(t)=0$ and apply MVT),
hence
\[
\delta(f)\le \beta(f)\le \alpha(f)\le 2\,\delta(f).
\]
Therefore $\alpha,\beta,\delta$ are pairwise equivalent on $C^1_0([0,1])$, while $\gamma$
is dominated by $\delta$ but not equivalent to it (spike example).
Replacing $C^1_0([0,1])$ by $C^1_c((0,1))$ yields the same conclusions.


\section*{Problem 6: Open and Closed Subsets of $\mathbb{R}^2$}

\paragraph{Statement.}
Which of the following subsets of $\mathbb{R}^2$ with the Euclidean norm are open? Which are closed? (And why?)
\begin{enumerate}[(i)]
	\item $\{(x,0): 0\le x\le 1\}$;
	\item $\{(x,0): 0< x< 1\}$;
	\item $\{(x,y): y\ne 0\}$;
	\item $\{(x,y): x\in\mathbb{Q}\ \text{or}\ y\in\mathbb{Q}\}$;
	\item $\{(x,y): y=nx \text{ for some } n\in\mathbb{N}\}\ \cup\ \{(x,y): x=0\}$;
	\item $\{(x,f(x)) : x\in\mathbb{R}\}$, where $f:\mathbb{R}\to\mathbb{R}$ is a continuous function.
\end{enumerate}

\paragraph{Theory.}
A set $U\subset\mathbb{R}^2$ is open iff every $p\in U$ admits $\varepsilon>0$ with $B_\varepsilon(p)\subset U$.
A set $F\subset\mathbb{R}^2$ is closed iff it contains all its limit points (equivalently, $\mathbb{R}^2\setminus F$ is open).
If $g$ is continuous, then $g^{-1}(C)$ is closed whenever $C$ is closed, and $g^{-1}(O)$ is open whenever $O$ is open.
Graphs of continuous functions are closed: $\{(x,y): y=f(x)\}=\{(x,y): y-f(x)=0\}$ is the zero set of a continuous map.

\paragraph{Solutions.}
\begin{enumerate}[(i)]
	\item Not open (a ball around any point leaves the segment). Closed: image of compact $[0,1]$ under $x\mapsto(x,0)$, hence compact in $\mathbb{R}^2$, thus closed. \emph{Answer:} closed, not open.
	
	\item Not open (no interior). Not closed: the limit points $(0,0)$ and $(1,0)$ are missing. \emph{Answer:} neither.
	
	\item Complement is $\{(x,0)\}$, which is closed as $\{(x,y):y=0\}=(\pi_2)^{-1}(\{0\})$ where $\pi_2(x,y)=y$ is continuous. Hence the set is open. It is not closed since points with $y\to 0$ converge to the $x$-axis. \emph{Answer:} open, not closed.
	
	\item Not open: every open ball contains a point with both coordinates irrational. The complement 
	$(\mathbb{R}\setminus\mathbb{Q})\times(\mathbb{R}\setminus\mathbb{Q})$ is not open for the same reason, so the set is not closed. \emph{Answer:} neither.
	
	\item Each line $\{(x,y):y=nx\}$ and $\{x=0\}$ is closed. Let $(x_k,n_k x_k)\to(x_0,y_0)$ with $n_k\in\mathbb{N}$ or $x_k=0$.
	If $x_0=0$, the limit lies on $\{x=0\}$. If $x_0\neq 0$, then $n_k=\frac{n_k x_k}{x_k}\to \frac{y_0}{x_0}$, forcing $\frac{y_0}{x_0}\in\mathbb{N}$; hence $(x_0,y_0)$ lies on one of the listed lines. Thus the union is closed. It is not open (lines have empty interior). \emph{Answer:} closed, not open.
	
	\item The graph $\{(x,f(x))\}$ equals $\{(x,y): y-f(x)=0\}$, the zero set of the continuous map $(x,y)\mapsto y-f(x)$. Hence it is closed. It has empty interior in $\mathbb{R}^2$, so it is not open. \emph{Answer:} closed, not open.
\end{enumerate}

\paragraph{Intuition/Examples.}
Line segments and lines have empty interior; removing a closed line (the $x$-axis) leaves an open set;
in (iv) both the set and its complement are dense with empty interior; graphs of continuous functions are closed.

\section*{Problem 7: Closed sets in $C([0,1])$ under $\|\cdot\|_\infty$ vs.\ $\|\cdot\|_1$}

\paragraph{Statement.}
Is the set $\{f: f(1/2)=0\}$ closed in the space $\big(C([0,1]),\|\cdot\|_\infty\big)$?
What about the set $\{f:\int_0^1 f=0\}$?
In each case, does the answer change if we replace the uniform norm with the norm $\|\cdot\|_1$?

\paragraph{Theory.}
In a normed space, the zero set of a continuous linear functional is closed.
For $C([0,1])$,
the evaluation functional $E_{x_0}(f)=f(x_0)$ is continuous under $\|\cdot\|_\infty$ (since $|f(x_0)|\le\|f\|_\infty$),
but not under $\|\cdot\|_1$.
The integral functional $I(f)=\int_0^1 f$ is continuous under both norms, because
$|I(f)|\le\int_0^1 |f|\le \|f\|_\infty$ and also $|I(f)|\le \|f\|_1$.

\paragraph{Solutions.}
\begin{enumerate}
	\item[\textbf{A.}] $A=\{f\in C([0,1]): f(1/2)=0\}$.
	\begin{itemize}
		\item Under $\|\cdot\|_\infty$: $A=E_{1/2}^{-1}(\{0\})$ with $E_{1/2}$ continuous, hence $A$ is closed.
		\item Under $\|\cdot\|_1$: $A$ is not closed. Define $g\equiv 1$. For $n\ge 2$ let $f_n$ equal $1$ outside $[1/2-1/n,\,1/2+1/n]$, and on this interval let $f_n$ be the linear tent with $f_n(1/2)=0$ and $f_n(1/2\pm 1/n)=1$. Then $f_n\in A$ and
		$\|f_n-g\|_1$ equals the area of the tent, $\frac12\cdot \frac{2}{n}\cdot 1=\frac1n\to 0$, so $f_n\to g$ in $\|\cdot\|_1$ while $g\notin A$.
	\end{itemize}
	
	\item[\textbf{B.}] $B=\{f\in C([0,1]): \int_0^1 f=0\}$.
	\begin{itemize}
		\item Under $\|\cdot\|_\infty$: $B=I^{-1}(\{0\})$ with $I$ continuous since $|\int_0^1 f|\le \|f\|_\infty$, hence $B$ is closed.
		\item Under $\|\cdot\|_1$: likewise $|\int_0^1 f|\le \|f\|_1$, so $I$ is continuous and $B$ is closed.
	\end{itemize}
\end{enumerate}

\paragraph{Summary.}
\[
\begin{array}{c|c|c}
	\text{Set} & (C([0,1]),\|\cdot\|_\infty) & (C([0,1]),\|\cdot\|_1)\\ \hline
	\{f: f(1/2)=0\} & \text{closed} & \text{not closed}\\
	\{f: \int_0^1 f=0\} & \text{closed} & \text{closed}
\end{array}
\]

\section*{Problem 8: Continuity of Linear Operators between Normed Spaces}

\paragraph{Statement.}
Let $X,Y$ be normed spaces and $T:X\to Y$ a linear map. Prove that the following are equivalent:
\begin{enumerate}
	\item $T$ is continuous;
	\item $T$ is continuous at $0$;
	\item $T$ is bounded: $\exists M>0$ such that $\|Tx\|_Y\le M\|x\|_X$ for all $x\in X$;
	\item $T$ is uniformly continuous on $X$.
\end{enumerate}

\paragraph{Theory.}
For a linear map $T$, continuity at a single point implies continuity everywhere, since
$T(x+h)-T(x)=T(h)$. A linear operator is called \emph{bounded} if
$\|Tx\|\le M\|x\|$ for some $M$. The smallest such $M$ is
\[
\|T\|=\sup_{\|x\|\le 1}\|Tx\|.
\]
Boundedness is equivalent to continuity, and implies uniform continuity:
\(\|Tx-Ty\|=\|T(x-y)\|\le\|T\|\|x-y\|\).

\paragraph{Proof of equivalences.}
\begin{enumerate}
	\item[(1)$\Rightarrow$(2)] Immediate.
	
	\item[(2)$\Rightarrow$(3)] Assume $T$ is continuous at $0$.
	Then $\exists\delta>0$ such that $\|x\|<\delta\Rightarrow\|Tx\|<1$.
	For arbitrary $x\ne0$ let $y=\frac{\delta}{2\|x\|}x$.
	Then $\|y\|<\delta$ and $\|Ty\|<1$, so
	\[
	\|Tx\|=\frac{2\|x\|}{\delta}\|Ty\|<\frac{2}{\delta}\|x\|.
	\]
	Hence $T$ is bounded with $M=2/\delta$.
	
	\item[(3)$\Rightarrow$(1)] If $T$ is bounded, for all $x,h$ we have
	$\|T(x+h)-T(x)\|=\|Th\|\le M\|h\|$, so $T$ is continuous everywhere.
	
	\item[(3)$\Leftrightarrow$(4)] From boundedness, $\|Tx-Ty\|\le M\|x-y\|$, so $T$ is uniformly continuous.
	Uniform continuity implies continuity at $0$, hence boundedness as before.
\end{enumerate}

\paragraph{Examples.}
\begin{itemize}
	\item $T:\mathbb{R}^2\to\mathbb{R}$, $T(x,y)=x+y$ is bounded with $\|T\|\le\sqrt2$, so continuous.
	\item $T:\ell^1\to\ell^\infty$, $T((x_n))=(n x_n)$ is linear but unbounded, hence not continuous.
\end{itemize}


\section*{Problem 9: Equivalent Norms and Their Topologies}

\paragraph{Statement.}
Let $\|\cdot\|_1$ and $\|\cdot\|_2$ be two norms on a vector space $X$.
Show that the following are equivalent:
\begin{enumerate}
	\item The norms are equivalent: $\exists a,b>0$ such that $a\|x\|_1\le\|x\|_2\le b\|x\|_1$ for all $x\in X$;
	\item The topologies induced by $\|\cdot\|_1$ and $\|\cdot\|_2$ are the same;
	\item The convergent sequences are the same for both norms.
\end{enumerate}

\paragraph{Theory.}
Each norm induces a metric $d_i(x,y)=\|x-y\|_i$ and hence a topology.
Two norms are called \emph{equivalent} if there exist constants $a,b>0$
with $a\|x\|_1\le\|x\|_2\le b\|x\|_1$ for all $x$.
Then both norms yield the same open sets, convergence, and continuity.
Conversely, identical convergence or topology implies the existence of such $a,b$.



\paragraph{Solution (step by step).}
We prove the cycle
\[
(1)\Rightarrow(2)\Rightarrow(3)\Rightarrow(1).
\]

\subparagraph{(1)$\Rightarrow$(2).}
Assume $a\|x\|_1 \le \|x\|_2 \le b\|x\|_1$.
Let $B_i(x_0,r)$ denote the open ball in $(X,\|\cdot\|_i)$.
\begin{itemize}
	\item If $y\in B_1(x_0,r)$, then $\|y-x_0\|_1 < r$.  
	Hence $\|y-x_0\|_2 < br$, so $y\in B_2(x_0,br)$.  
	Thus $B_1(x_0,r)\subseteq B_2(x_0,br)$.
	\item Conversely, from $a\|x\|_1\le\|x\|_2$, we get  
	$B_2(x_0,r)\subseteq B_1(x_0,r/a)$.
\end{itemize}
So open balls in one norm contain open balls of the other.
Hence both topologies have the same open sets.

\subparagraph{(2)$\Rightarrow$(3).}
If the two topologies coincide, then convergence (defined by open neighborhoods)
is identical. Thus a sequence converges in one norm if and only if it converges in the other.

\subparagraph{(3)$\Rightarrow$(1).}
Assume the same convergent sequences.  
We show existence of $a,b>0$.

Define $f(x)=\|x\|_2$.
If no $b$ exists with $\|x\|_2\le b\|x\|_1$,  
we could find a sequence $x_n$ with $\|x_n\|_1=1$ and $\|x_n\|_2>n$.  
Then $x_n/n\to 0$ in $\|\cdot\|_1$ but not in $\|\cdot\|_2$, contradicting (3).  
Thus $\exists b>0$ s.t.\ $\|x\|_2\le b\|x\|_1$.

Similarly, interchanging norms gives $a>0$ s.t.\ $a\|x\|_1\le\|x\|_2$.  
Hence the norms are equivalent.


\paragraph{Examples.}
\begin{enumerate}
	\item On $\mathbb{R}^n$:
	\[
	\|x\|_1=\sum_{i=1}^n |x_i|,\qquad
	\|x\|_2=\Big(\sum_{i=1}^n x_i^2\Big)^{1/2},\qquad
	\|x\|_\infty=\max_{1\le i\le n}|x_i|.
	\]
	There exist constants $a,b>0$ such that
	\[
	\|x\|_\infty \le \|x\|_2 \le \|x\|_1 \le \sqrt{n}\,\|x\|_2.
	\]
	Thus all norms on $\mathbb{R}^n$ are equivalent.
	
	\item On infinite-dimensional spaces (e.g.\ $C([0,1])$):
	The norms $\|f\|_1=\int_0^1 |f|$ and $\|f\|_\infty=\max_{x\in[0,1]}|f(x)|$ are not equivalent.
	Example: let $f_n$ be a narrow spike of height $1$ and width $1/n$. Then
	\[
	\|f_n\|_\infty=1 \quad \text{but} \quad \|f_n\|_1=\frac{1}{n}\to 0.
	\]
	Hence convergence in $\|\cdot\|_1$ does not imply convergence in $\|\cdot\|_\infty$.
\end{enumerate}

\section*{Problem 10: The space $\ell^1$ is a Banach space}

\paragraph{Statement.}
Let $\ell^1$ denote the set of real sequences $(x_n)$ such that $\sum_{n=1}^{\infty}|x_n|$ is convergent.
Show that, with addition and scalar multiplication defined termwise, $\ell^1$ is a vector space.
Define $\|\cdot\|_1:\ell^1\to\mathbb{R}$ by $\|x\|_1=\sum_{n=1}^{\infty}|x_n|$.
Show that $\|\cdot\|_1$ is a norm on $\ell^1$, and that $(\ell^1,\|\cdot\|_1)$ is complete.

\paragraph{Theory.}
Absolutely convergent series are stable under termwise addition and scalar multiplication.
A norm must satisfy positivity, definiteness, homogeneity, and triangle inequality.
Completeness means every Cauchy sequence converges in the space.
We will use that for a Cauchy sequence $(x^{(m)})$ in $\ell^1$,
each coordinate $(x^{(m)}_k)$ is Cauchy in $\mathbb{R}$, and Fatou's lemma for series:
$\sum_{k=1}^\infty |x_k| \le \liminf_{m}\sum_{k=1}^\infty |x^{(m)}_k|$ when $x_k=\lim_m x^{(m)}_k$.

\paragraph{Solution.}

\textbf{(1) $\ell^1$ is a vector space.}
If $x,y\in\ell^1$, then
\[
\sum_{n=1}^\infty |x_n+y_n|\le \sum_{n=1}^\infty (|x_n|+|y_n|)<\infty,
\]
so $x+y\in\ell^1$.
If $\alpha\in\mathbb{R}$, then
\[
\sum_{n=1}^\infty |\alpha x_n|=|\alpha|\sum_{n=1}^\infty |x_n|<\infty,
\]
so $\alpha x\in\ell^1$.

\textbf{(2) $\|\cdot\|_1$ is a norm.}
For $x\in\ell^1$, positivity and definiteness are clear.
Homogeneity:
$\|\alpha x\|_1=\sum|\alpha x_n|=|\alpha|\sum|x_n|=|\alpha|\,\|x\|_1$.
Triangle inequality:
$\|x+y\|_1=\sum|x_n+y_n|\le \sum(|x_n|+|y_n|)=\|x\|_1+\|y\|_1$.

\textbf{(3) Completeness.}
Let $(x^{(m)})$ be Cauchy in $\ell^1$.
Then there exists $B<\infty$ with $\|x^{(m)}\|_1\le B$ for all $m$ large.
For each $k$, $(x^{(m)}_k)$ is Cauchy in $\mathbb{R}$; set $x_k=\lim_m x^{(m)}_k$.
For $N\in\mathbb{N}$,
\[
\sum_{k=1}^N |x_k| \le \liminf_{m\to\infty} \sum_{k=1}^N |x^{(m)}_k|
\le \liminf_{m\to\infty} \|x^{(m)}\|_1 \le B.
\]
Letting $N\to\infty$ shows $x\in\ell^1$ and $\|x\|_1\le B$.

To prove $\|x^{(m)}-x\|_1\to 0$, fix $\varepsilon>0$.
Choose $m_0$ such that $\|x^{(m)}-x^{(n)}\|_1<\varepsilon/3$ for all $m,n\ge m_0$.
Fix $n\ge m_0$ and choose $N$ with $\sum_{k>N}|x^{(n)}_k|<\varepsilon/6$.
Then for $m\ge m_0$,
\[
\sum_{k>N}|x^{(m)}_k|
\le \sum_{k>N}|x^{(m)}_k-x^{(n)}_k| + \sum_{k>N}|x^{(n)}_k|
\le \|x^{(m)}-x^{(n)}\|_1 + \varepsilon/6 < \varepsilon/3.
\]
By pointwise convergence, choose $m_1$ so that $\sum_{k=1}^N |x^{(m)}_k-x_k|<\varepsilon/3$ for $m\ge m_1$.
Also $\sum_{k>N}|x_k|\le \liminf_{m}\sum_{k>N}|x^{(m)}_k|\le \varepsilon/3$.
Hence for $m\ge \max\{m_0,m_1\}$,
\[
\|x^{(m)}-x\|_1
\le \sum_{k=1}^N |x^{(m)}_k-x_k| + \sum_{k>N}|x^{(m)}_k| + \sum_{k>N}|x_k| < \varepsilon.
\]
Therefore $(\ell^1,\|\cdot\|_1)$ is complete.
\qed

\paragraph{Examples/Comments.}
$e^{(j)}=(0,\dots,0,1,0,\dots)$ has $\|e^{(j)}\|_1=1$; $(1/n)\notin\ell^1$ while $(1/n^2)\in\ell^1$.
Thus $\ell^1$ is a Banach space.

\section*{Problem 11$^*$: Absolute convergence characterizes completeness}

\paragraph{Statement.}
Let $(V,\|\cdot\|)$ be a normed space. Show that $V$ is complete if and only if
for every sequence $(x_n)$ in $V$ with $\sum_{j=1}^\infty \|x_j\|$ convergent,
the series $\sum_{n=1}^\infty x_n$ is convergent.
\emph{Hint.} If $(x_n)$ is Cauchy, then there is a subsequence $(x_{n_j})$
such that $\sum_j \|x_{n_{j+1}}-x_{n_j}\|<\infty$.

\paragraph{Theory.}
A series $\sum x_n$ converges iff its partial sums $s_N=\sum_{n=1}^N x_n$ converge.
If $\sum \|x_n\|<\infty$, then $(s_N)$ is Cauchy since
$\|s_m-s_n\|\le\sum_{k=n+1}^m \|x_k\|$.
In a complete space, every Cauchy sequence converges.
From a Cauchy sequence one can extract a subsequence with summable differences.

\paragraph{Proof.}

\textbf{($\Rightarrow$)} Suppose $V$ is complete and $\sum_{n=1}^\infty \|x_n\|<\infty$.
Then for $m>n$,
\[
\|s_m-s_n\|=\Big\|\sum_{k=n+1}^m x_k\Big\|
\le \sum_{k=n+1}^m \|x_k\|\xrightarrow[n,m\to\infty]{}0,
\]
so $(s_N)$ is Cauchy and converges in $V$. Hence $\sum x_n$ converges.

\textbf{($\Leftarrow$)} Assume every absolutely convergent series converges in $V$.
Let $(y_n)$ be Cauchy. Choose a subsequence $(y_{n_j})$ with
$\|y_{n_{j+1}}-y_{n_j}\|\le 2^{-j}$.
Then $\sum_{j=1}^\infty \|y_{n_{j+1}}-y_{n_j}\|\le 1$, so by the assumption the series
$\sum_{j=1}^\infty (y_{n_{j+1}}-y_{n_j})$ converges.
Its $m$-th partial sum equals $y_{n_m}-y_{n_1}$; hence $y_{n_j}\to y$ for some $y\in V$.

To prove $y_n\to y$, fix $\varepsilon>0$. Since $(y_n)$ is Cauchy, there exists $N$ such that
$\|y_n-y_m\|<\varepsilon/2$ for all $n,m\ge N$.
Choose $j$ with $n_j\ge N$ and $\|y_{n_j}-y\|<\varepsilon/2$.
Then for all $n\ge N$,
\[
\|y_n-y\|\le \|y_n-y_{n_j}\|+\|y_{n_j}-y\|<\varepsilon/2+\varepsilon/2=\varepsilon.
\]
Hence $y_n\to y$ and $V$ is complete.
\qed

\paragraph{Example.}
Let $c_{00}$ be the space of finitely supported sequences with the $\ell^1$-norm.
Take $x_n=\frac{1}{2^n}e_n$. Then $\sum\|x_n\|_1=1$ but
$\sum x_n=(\frac{1}{2^n})_{n\ge1}\notin c_{00}$.
So absolute convergence does not imply convergence in $c_{00}$, and $c_{00}$ is not complete.

\paragraph{Remark. Absolute convergence may fail to imply convergence in incomplete spaces.}

In a normed space, an absolutely convergent series always has Cauchy partial sums, but it
\emph{converges in the space} if and only if the space is complete.

\textbf{Counterexample.}
Let $c_{00}$ be the space of finitely supported sequences with the $\ell^{1}$-norm.
Define
\[
x_n = \frac{1}{2^{n}} e_n, \qquad (n\ge1),
\]
where $e_n$ is the sequence with $1$ in position $n$ and zeros elsewhere.
Then
\[
\sum_{n=1}^\infty \|x_n\|_1 = \sum_{n=1}^\infty \frac{1}{2^n} = 1 < \infty,
\]
so the series $\sum x_n$ is \emph{absolutely convergent}.
Its partial sums are
\[
s_N = \sum_{n=1}^N x_n = \bigg(\frac12, \frac14, \dots, \frac{1}{2^N}, 0, 0, \dots\bigg),
\]
which converge (in $\ell^1$) to
\[
s = \bigg(\frac12, \frac14, \frac18, \dots\bigg).
\]
However, $s\notin c_{00}$ (it is not finitely supported), so $\sum x_n$
does \emph{not} converge in $c_{00}$.

\textbf{Conclusion.}
Absolute convergence does not guarantee convergence in an incomplete space.
In contrast, in any Banach space (complete normed space), every absolutely
convergent series converges in the space.

\section*{Problem 12: Every bounded sequence has a convergent subsequence}

\paragraph{Statement.}
Let $V$ be a normed space in which every bounded sequence has a convergent subsequence.
\begin{enumerate}[(a)]
	\item Show that this property is equivalent to the sequential compactness of the unit sphere
	$S=\{x\in V:\|x\|=1\}$.
	\item Show that $V$ must be complete.
	\item[\textnormal{(c)$^*$}] Show further that $V$ must be finite--dimensional.
\end{enumerate}
\emph{Hint for (c).} For every finite--dimensional subspace $V_0$ of $V$, there exists $x\in V$
with $\|x+y\|>\|x\|/2$ for each $y\in V_0$.

\paragraph{Theory.}
Sequential compactness means every sequence has a convergent subsequence.
Cauchy sequences are bounded; if a Cauchy sequence has one convergent subsequence, then the whole sequence converges.
A Riesz--lemma variant: if $V_0$ is finite--dimensional (hence closed) and $0<\alpha<1$, there exists $x$ with $\|x\|=1$ and
$\operatorname{dist}(x,V_0)>\alpha$; equivalently, $\|x+y\|>\alpha$ for all $y\in V_0$.

\paragraph{Solutions.}
\begin{enumerate}[(a)]
	\item
	Assume every bounded sequence in $V$ has a convergent subsequence. Then the unit sphere $S$ is bounded,
	so any sequence in $S$ has a convergent subsequence in $V$. If $x_{n_k}\to x$, then
	$\|x\|=\lim_k\|x_{n_k}\|=1$, hence $x\in S$. Thus $S$ is sequentially compact.
	
	Conversely, suppose $S$ is sequentially compact. Let $(x_n)$ be bounded in $V$ and set $r_n=\|x_n\|$.
	Passing to a subsequence we may assume $(r_n)$ converges to $r\ge 0$.
	For each $n$ with $x_n\neq 0$ define $u_n=x_n/\|x_n\|\in S$; for $x_n=0$ define $u_n$ arbitrarily in $S$.
	By sequential compactness of $S$, $(u_n)$ admits a convergent subsequence $u_{n_k}\to u\in S$.
	Then $x_{n_k}=r_{n_k}u_{n_k}\to r u$ in $V$. Hence every bounded sequence has a convergent subsequence.
	
	\item
	Let $(y_n)$ be Cauchy in $V$. Then $(y_n)$ is bounded, so it has a convergent subsequence $y_{n_k}\to y$.
	Fix $\varepsilon>0$. Choose $N$ such that $\|y_n-y_m\|<\varepsilon/2$ for all $m,n\ge N$,
	and choose $k$ with $n_k\ge N$ and $\|y_{n_k}-y\|<\varepsilon/2$.
	For $n\ge N$ we have
	\[
	\|y_n-y\|\le \|y_n-y_{n_k}\|+\|y_{n_k}-y\|<\varepsilon.
	\]
	Thus $y_n\to y$, so $V$ is complete.
	
	\item[\textnormal{(c)$^*$}]
	Suppose $V$ were infinite--dimensional. Construct inductively a sequence $(x_n)\subset V$ with
	$\|x_n\|=1$ and
	\[
	\operatorname{dist}(x_{n+1},V_n)>\tfrac12,
	\qquad
	V_n=\operatorname{span}\{x_1,\dots,x_n\}.
	\]
	This is possible by the Riesz--lemma variant (applied to the finite--dimensional subspace $V_n$ with $\alpha=\tfrac12$).
	Then for $j>i$ we have $x_i\in V_{j-1}$, hence
	\[
	\|x_j-x_i\|\ge \operatorname{dist}(x_j,V_{j-1})>\tfrac12 .
	\]
	Therefore $(x_n)$ is bounded (all have norm $1$) but has no Cauchy subsequence, because any two distinct
	terms of any tail are at distance $>1/2$.
	This contradicts the hypothesis that every bounded sequence has a convergent subsequence.
	Hence $V$ must be finite--dimensional.
\end{enumerate}

\section*{Problem 13: Coordinatewise subsequential convergence in $\ell^\infty$}

\paragraph{Statement.}
Let $(x^{(n)})_{n\ge1}$ be a bounded sequence in $\ell^\infty$.
Show that there is a subsequence $(x^{(n_j)})_{j\ge1}$ which converges in every coordinate; that is,
for each $i\in\mathbb{N}$, the sequence $(x^{(n_j)}_i)_{j\ge1}$ of real numbers converges for each $i$.
Why does this not show that every bounded sequence in $\ell^\infty$ has a convergent subsequence?

\paragraph{Solution.}
For each fixed $i$, $(x^{(n)}_i)_{n\ge1}$ is bounded in $\mathbb{R}$, hence has a convergent subsequence.
By iterating this extraction and taking the diagonal subsequence $x^{(m_j)}=x^{(n^{(j)}_j)}$,
we obtain that $x^{(m_j)}_i$ converges for each fixed $i$.
However, coordinatewise convergence does not imply convergence in the norm $\|\cdot\|_\infty$.
For example,
\[
x^{(n)}=(\underbrace{1,\dots,1}_{n\text{ times}},0,0,\dots)
\]
satisfies $x^{(n)}_i\to 1$ for each fixed $i$, but
$\|x^{(n)}-\mathbf{1}\|_\infty=1$ and $\|x^{(n)}-x^{(m)}\|_\infty=1$ for $m>n$,
so no subsequence is Cauchy (hence none is convergent) in $\ell^\infty$.

\bigskip
\section*{Problem 14$^*$: Two norms on $\mathcal{P}$ giving different limits}

\paragraph{Statement.}
Let $\mathcal{P}$ be the vector space of real polynomials on $[0,1]$.  
For any infinite set $I\subseteq[0,1]$, define $\|p\|_{I}:=\sup_{t\in I}|p(t)|$.
\begin{enumerate}
	\item[(a)] Show that $\|\cdot\|_{I}$ is a norm on $\mathcal{P}$.
	\item[(b)] Use this to produce a vector space, a single sequence in it, and two different norms on it
	such that the sequence converges to \emph{different} elements of the space with respect to the two norms.
	(Hint: Weierstrass approximation.)
	\item[(c)] Is such a phenomenon possible in $\ell^1$ or $\ell^2$ when the two norms are chosen
	from the standard norms on $\ell^1,\ell^2,\ell^\infty$?  
	What about in $C([0,1])$ when the two norms are chosen from $L^1,L^2,L^\infty$?
\end{enumerate}

\bigskip
\paragraph{Theory.}
\begin{itemize}
	\item If a real polynomial vanishes on an infinite subset of $[0,1]$ with a limit point in $[0,1]$,
	then it is identically zero (identity theorem for polynomials).
	\item On disjoint compact sets $K_1,K_2\subset[0,1]$ there exists $f\in C([0,1])$ with
	$f|_{K_1}\equiv0$ and $f|_{K_2}\equiv1$ (Urysohn on $[0,1]$).  
	By Weierstrass, there are polynomials $p_n\to f$ uniformly on $[0,1]$.
	\item In $\ell^p$ and in $C([0,1])$ (measure $1$),
	\[
	\|z\|_\infty \ \ge\ \|z\|_2 \ \ge\ \|z\|_1,
	\qquad
	\|h\|_{L^1}\ \le\ \|h\|_{L^2}\ \le\ \|h\|_{L^\infty}.
	\]
	Hence convergence in any stronger norm implies convergence in any weaker norm to the \emph{same} limit.
\end{itemize}

\bigskip
\paragraph{Solution.}

\subparagraph{(a) $\|\cdot\|_{I}$ is a norm on $\mathcal{P}$.}
Let $p,q\in\mathcal{P}$ and $\alpha\in\mathbb{R}$.
\begin{itemize}
	\item \textbf{Positivity/definiteness:} $\|p\|_{I}\ge0$.
	If $\|p\|_{I}=0$ then $p(t)=0$ for all $t\in I$.  
	Since $I$ is infinite and has a limit point in $[0,1]$, the identity theorem gives $p\equiv0$.
	\item \textbf{Homogeneity:} $\|\alpha p\|_{I}=\sup_{t\in I}|\alpha||p(t)|=|\alpha|\,\|p\|_{I}$.
	\item \textbf{Triangle inequality:}
	$\|p+q\|_{I}=\sup_{t\in I}|p(t)+q(t)|\le \sup_{t\in I}|p(t)|+\sup_{t\in I}|q(t)|=\|p\|_{I}+\|q\|_{I}$.
\end{itemize}
Hence $\|\cdot\|_{I}$ is a norm on $\mathcal{P}$.

\bigskip
\subparagraph{(b) Same sequence, two norms, two different limits.}
Let
\[
K_1=[0,\tfrac12],\qquad K_2=[\tfrac12,1],
\]
and pick infinite sets $I\subset K_1$, $J\subset K_2$.

By Urysohn on $[0,1]$ choose $f\in C([0,1])$ with $f|_{K_1}\equiv0$ and $f|_{K_2}\equiv1$.
By Weierstrass there exist polynomials $p_n\in\mathcal{P}$ such that
\[
\|p_n-f\|_{[0,1]}:=\sup_{t\in[0,1]}|p_n(t)-f(t)| \xrightarrow[n\to\infty]{} 0.
\]

Then
\[
\|p_n-0\|_{I}=\sup_{t\in I}|p_n(t)|
\ \le\ \sup_{t\in K_1}|p_n(t)-f(t)|
\ \xrightarrow[n\to\infty]{}\ 0,
\]
so $p_n\to 0$ in $(\mathcal{P},\|\cdot\|_{I})$.

Similarly,
\[
\|p_n-1\|_{J}=\sup_{t\in J}|p_n(t)-1|
\ \le\ \sup_{t\in K_2}|p_n(t)-f(t)|
\ \xrightarrow[n\to\infty]{}\ 0,
\]
so $p_n\to 1$ in $(\mathcal{P},\|\cdot\|_{J})$.

\medskip
\noindent
Thus the \emph{same} sequence $(p_n)$ converges to \emph{different} elements (0 and 1) under two different norms on the same vector space $\mathcal{P}$.

\bigskip
\subparagraph{(c) Impossibility in $\ell^1,\ell^2$ and in $C([0,1])$ with standard norms.}
No such example exists.
\begin{itemize}
	\item On $\ell^1,\ell^2,\ell^\infty$, for every $z$ we have
	$\|z\|_\infty\ge \|z\|_2\ge \|z\|_1$ (Hölder).
	If $x_n\to x$ in a stronger norm and $x_n\to y$ in a weaker norm,
	then $\|x-y\|\le \limsup_n \|x_n-x\|+\|x_n-y\|=0$, hence $x=y$.
	\item On $C([0,1])$,
	$\|h\|_{L^1}\le \|h\|_{L^2}\le \|h\|_{L^\infty}$ on a set of measure $1$.
	The same argument shows any two of these norm limits must coincide.
\end{itemize}

\paragraph{Remark. When two different limits are possible.}

The phenomenon of having \emph{two different limits} under two norms does not rely on the supremum
norm itself. It arises whenever we choose two \emph{non-equivalent} and \emph{non-comparable} norms
on the same vector space, so that the identity map between the two normed spaces is not continuous.

\medskip
\noindent\textbf{Example beyond supremum norms.}
Consider the vector space $\mathcal{P}$ of real polynomials on $[0,1]$.
Define two different weighted $L^2$ norms:
\[
\|p\|_{2,w_1}=\Big(\int_0^{1/2}|p(t)|^2\,dt\Big)^{1/2},
\qquad
\|p\|_{2,w_2}=\Big(\int_{1/2}^{1}|p(t)|^2\,dt\Big)^{1/2}.
\]
Let $f\in C([0,1])$ satisfy $f|_{[0,1/2]}=0$ and $f|_{[1/2,1]}=1$,
and choose polynomials $p_n\to f$ uniformly (by the Weierstrass theorem).
Then
\[
\|p_n\|_{2,w_1}\to 0,
\qquad
\|p_n-1\|_{2,w_2}\to 0.
\]
Hence the same sequence $(p_n)$ converges to $0$ under $\|\cdot\|_{2,w_1}$
and to $1$ under $\|\cdot\|_{2,w_2}$,
even though neither of these norms is a supremum norm.

\medskip
\noindent\textbf{What is special about the standard norms.}
On standard spaces such as $\ell^p$ or $C([0,1])$ with the usual $L^1,L^2,L^\infty$ norms, we have
\[
\|h\|_{1}\le\|h\|_{2}\le\|h\|_{\infty}.
\]
These norms are \emph{comparable and compatible}, so convergence in any two of them
forces convergence to the \emph{same} limit.
Therefore, different limits can occur only when the norms are \emph{not equivalent}
and not ordered by inequality, not because the supremum norm is involved.

\section*{Sequential Convergence: Definition, Theory, and Examples}

\paragraph{Definition (topological version).}
Let $(X,\tau)$ be a topological space. A sequence $(x_n)\subset X$ \emph{converges} to $x\in X$,
written $x_n\to x$, if for every neighbourhood $U\in\tau$ of $x$ there exists $N$ such that
$x_n\in U$ for all $n\ge N$.

\paragraph{Metric/normed reformulation.}
In a metric space $(X,d)$ (or a normed space $(X,\|\cdot\|)$) this is equivalent to:
\[
\forall\varepsilon>0\ \exists N\ \forall n\ge N\quad d(x_n,x)<\varepsilon
\qquad\big(\text{resp. }\|x_n-x\|<\varepsilon\big).
\]

\paragraph{Basic facts.}
\begin{itemize}
	\item Limits are unique in Hausdorff spaces.
	\item If $x_n\to x$, then every subsequence $x_{n_k}\to x$.
	\item The \emph{sequential closure} of $A\subset X$ is
	$\operatorname{scl}(A)=\{x\in X:\exists (a_n)\subset A,\ a_n\to x\}$.
	\item In every \emph{first countable} space (in particular, in all metric/normed spaces),
	\[
	\overline{A}=\operatorname{scl}(A).
	\]
	\item A map $f:X\to Y$ is \emph{sequentially continuous} if $x_n\to x$ implies $f(x_n)\to f(x)$.
	In first countable spaces, sequential continuity is equivalent to topological continuity.
	\item \emph{Sequential compactness}: every sequence has a convergent subsequence.
	In metric spaces: compact $\iff$ sequentially compact $\iff$ complete and totally bounded.
\end{itemize}

\paragraph{When sequences are not enough.}
In general topological spaces (not first countable) sequences may fail to detect closure or continuity.

\begin{itemize}
	\item \textbf{Ordinals example.}
	Let $X=\omega_1+1$ with the order topology and $A=\omega_1$.
	Then $\overline{A}=X$ (so $\omega_1\in\overline{A}$), but no sequence from $A$ converges to $\omega_1$,
	since any countable sequence of ordinals has countable supremum $<\omega_1$.
	Thus $\overline{A}\ne \operatorname{scl}(A)$; nets are required to capture convergence.
	
	\item \textbf{Cofinite topology.}
	On an infinite set $X$ with the cofinite topology, a sequence $(x_n)$ converges to $x$ iff
	for every $y\ne x$ the set $\{n:\ x_n=y\}$ is finite.
	In particular, a sequence listing each element at most once converges to \emph{every} point,
	so limits need not be unique (the space is not Hausdorff).
\end{itemize}

\paragraph{Canonical metric/normed examples.}
\begin{itemize}
	\item \textbf{$\mathbb{R}^n$ with the Euclidean norm.}
	Here $x_n\to x$ iff each coordinate converges. For subsets of $\mathbb{R}^n$,
	compactness $\Leftrightarrow$ closed and bounded $\Leftrightarrow$ sequential compactness.
	
	\item \textbf{$\ell^\infty$: coordinatewise vs.\ norm convergence.}
	Let $x^{(n)}=(\underbrace{1,\dots,1}_{n\ \text{times}},0,0,\dots)$.
	Then $x^{(n)}_i\to 1$ for each fixed $i$ (coordinatewise convergence),
	but $\|x^{(n)}-\mathbf 1\|_\infty=1$ for all $n$, so no subsequence converges in $\|\cdot\|_\infty$.
	
	\item \textbf{Sequential continuity in normed spaces.}
	For linear $T:X\to Y$ between normed spaces, $x_n\to x \Rightarrow T(x_n)\to T(x)$
	iff $T$ is bounded/continuous; in normed spaces sequential continuity $\Leftrightarrow$ continuity.
\end{itemize}

\section*{Supplement: A proof of Lebesgue's theorem on the Riemann integral}

\paragraph{Problem (exact).}
Let $f:[a,b]\to\mathbb{R}$ be bounded. Recall that Lebesgue's theorem says that $f$ is Riemann integrable on $[a,b]$ if and only if the set $D_f$ of points in $[a,b]$ where $f$ is discontinuous has Lebesgue measure zero. (By definition, a set $D\subset\mathbb{R}$ has Lebesgue measure zero if for every $\varepsilon>0$, there is a countable collection of open intervals $I_j=(a_j,b_j)$ such that $D\subset\bigcup_{j=1}^\infty I_j$ and $\sum_{j=1}^\infty |I_j|<\varepsilon$, where $|I_j|=b_j-a_j$.) As an optional exercise, prove this theorem by completing the outline below. We shall use the notation as in lectures, so $U(P,f),L(P,f)$ denote the upper and lower sums for $f$ relative to a partition $P$ of $[a,b]$.
\begin{enumerate}[(a)]
	\item Show that $y\in D_f\cap(a,b)$ (i.e.\ $y$ is an interior discontinuity) if and only if there exists $\varepsilon=\varepsilon_y>0$ such that $\sup_I f-\inf_I f>\varepsilon$ for every open interval $I\subset (a,b)$ with $y\in I$. Hence $D_f\cap(a,b)=\bigcup_{j=1}^\infty E_j$, where $E_j=\{y\in(a,b):\ \sup_I f-\inf_I f>j^{-1}\ \text{for every open interval }I\text{ with }y\in I\}$.
	\item Suppose that $f$ is Riemann integrable. It suffices to show that $E_j$ has Lebesgue measure zero for each $j$ (Why?). Fix $j$, let $\varepsilon>0$ and choose a partition $P=\{a=a_0<a_1<\cdots<a_n=b\}$ such that $U(P,f)-L(P,f)<j^{-1}\varepsilon$. Let $K=\{k:\ E_j\cap(a_k,a_{k+1})\neq\emptyset\}$. Then $E_j\setminus\{a_0,a_1,\dots,a_n\}\subset\bigcup_{k\in K}(a_k,a_{k+1})$. Show that $\sum_{k\in K}(a_{k+1}-a_k)<\varepsilon$. Deduce that $E_j$ has Lebesgue measure zero.
	\item Now suppose that $D_f$ has Lebesgue measure zero. Let $\varepsilon>0$, and choose open intervals $I_j\subset\mathbb{R}$, $j=1,2,\dots$, with $D_f\subset\bigcup_{j=1}^\infty I_j$ and $\sum_{j=1}^\infty |I_j|<\varepsilon$. Let $F=[a,b]\setminus\bigcup_{j=1}^\infty I_j$. Show that there exists $\delta>0$ such that the following holds: $x\in F,\ y\in F,\ |x-y|<\delta \Rightarrow |f(x)-f(y)|<\varepsilon$. [This is a strengthening of the theorem we proved in lecture that says that a continuous function on a closed, bounded interval (or more generally on a compact metric space) is uniformly continuous, but the same contradiction argument we used in fact works here.] Let $P=\{a=a_0<a_1<\cdots<a_n=b\}$ be any partition of $[a,b]$ such that $a_{j+1}-a_j<\delta$, and let $J=\{j:\ [a_j,a_{j+1}]\cap F\neq\emptyset\}$. Show that $\sup_{[a_j,a_{j+1}]} f-\inf_{[a_j,a_{j+1}]} f<2\varepsilon$ for each $j\in J$, and that $\bigcup_{j\notin J}(a_j,a_{j+1})\subset\bigcup_{j=1}^\infty I_j$. Conclude that $U(P,f)-L(P,f)<2(b-a+\sup_{[a,b]}|f|)\varepsilon$, and hence that $f$ is Riemann integrable on $[a,b]$.
\end{enumerate}

\paragraph{Theory (background and tools you will use).}

\paragraph{1) Upper/lower sums and Riemann integrability.}\mbox{}\par


For a partition $P=\{a=a_0<\cdots<a_n=b\}$ and bounded $f$,
\[
U(P,f)=\sum_{k=0}^{n-1} M_k\,(a_{k+1}-a_k),\qquad
M_k=\sup_{[a_k,a_{k+1}]} f,\qquad
L(P,f)=\sum_{k=0}^{n-1} m_k\,(a_{k+1}-a_k),\qquad
m_k=\inf_{[a_k,a_{k+1}]} f.
\]
Let the \emph{oscillation} on an interval $I$ be $\omega(f;I)=\sup_I f-\inf_I f$.
Then
\[
U(P,f)-L(P,f)=\sum_{k=0}^{n-1} \omega\!\big(f;[a_k,a_{k+1}]\big)\,(a_{k+1}-a_k).
\]
A bounded $f$ is Riemann integrable iff $\forall \eta>0$ there exists $P$ with $U(P,f)-L(P,f)<\eta$.

\paragraph{2) Oscillation at a point and discontinuities.}\mbox{}\par



The \emph{oscillation at a point} $x$ is
\[
\omega(f;x)=\inf_{\delta>0}\ \sup\{|f(u)-f(v)|:\ u,v\in (x-\delta,x+\delta)\cap[a,b]\}.
\]
Then $f$ is continuous at $x$ iff $\omega(f;x)=0$.
Equivalently: $x\in D_f$ iff $\exists\,\varepsilon>0$ such that $\omega(f;I)>\varepsilon$
for every interval $I\ni x$.
Define
\[
E_j=\Big\{x\in(a,b):\ \omega(f;I)>\tfrac{1}{j}\ \text{for every open }I\ni x\Big\}.
\]
Then $D_f\cap(a,b)=\bigcup_{j=1}^\infty E_j$.

\paragraph{3) Measure--zero coverings.}\mbox{}\par

A set $N\subset\mathbb{R}$ has Lebesgue measure zero iff
for every $\varepsilon>0$ there exist open intervals $I_j$ with
$N\subset \bigcup_{j=1}^\infty I_j$ and $\sum_{j=1}^\infty |I_j|<\varepsilon$.
If $E\subset \bigcup_{k\in K}(a_k,a_{k+1})$, then
\[
\lambda(E)\le \sum_{k\in K}(a_{k+1}-a_k).
\]

\paragraph{4) Compactness and uniform continuity on closed subsets.}\mbox{}\par
If $F\subset [a,b]$ is closed, then $F$ is compact.
If $f$ is continuous at each point of $F$, then $f|_F$ is continuous $F\to\mathbb{R}$,
hence (Heine--Cantor) uniformly continuous on $F$:
\[
\forall \varepsilon>0\ \exists\,\delta>0\ \forall x,y\in F\quad
|x-y|<\delta \ \Longrightarrow\ |f(x)-f(y)|<\varepsilon.
\]


\paragraph{(a) Interior discontinuity and oscillation.} \mbox{}\par
For $y\in(a,b)$ define the oscillation
\[
\omega(f;y):=\inf_{\delta>0}\Big(\sup_{(y-\delta,y+\delta)\cap[a,b]} f
-\inf_{(y-\delta,y+\delta)\cap[a,b]} f\Big).
\]
It is standard that $f$ is continuous at $y$ iff $\omega(f;y)=0$.

\emph{($\Rightarrow$)} Suppose $y\in D_f\cap(a,b)$. If no $\varepsilon_y>0$ satisfied
$\sup_I f-\inf_I f>\varepsilon_y$ for every open interval $I\ni y$, then for each $m\in\mathbb N$
we could choose an interval $I_m\ni y$ with $\sup_{I_m} f-\inf_{I_m} f\le 1/m$. Taking $I_m=(y-\delta_m,y+\delta_m)$ with $\delta_m\downarrow 0$ gives $\omega(f;y)\le 1/m$ for all $m$, hence $\omega(f;y)=0$, a contradiction. Thus there exists $\varepsilon_y>0$ such that
$\sup_I f-\inf_I f>\varepsilon_y$ for every open interval $I\ni y$.

\emph{($\Leftarrow$)} Conversely, if there exists $\varepsilon_y>0$ with
$\sup_I f-\inf_I f>\varepsilon_y$ for every open interval $I\ni y$, then for all $\delta>0$,
\[
\sup_{(y-\delta,y+\delta)} f-\inf_{(y-\delta,y+\delta)} f\ge \varepsilon_y,
\]
so $\omega(f;y)\ge \varepsilon_y>0$ and hence $f$ is discontinuous at $y$.

Consequently,
\[
D_f\cap(a,b)=\bigcup_{j=1}^\infty E_j,
\qquad
E_j:=\Big\{y\in(a,b):\ \sup_I f-\inf_I f>\tfrac1j \ \text{for every open interval } I\ni y\Big\}.
\]
Indeed, $y\in D_f\cap(a,b)$ iff $\omega(f;y)>0$, i.e.\ iff $1/j<\omega(f;y)$ for some $j$,
which is equivalent to $y\in E_j$ for some $j$.

\paragraph{(b) If $f$ is Riemann integrable, then $E_j$ has measure zero.} \mbox{}\par
Fix $j\in\mathbb N$ and $\varepsilon>0$. Take a partition
$P=\{a=a_0<\cdots<a_n=b\}$ with $U(P,f)-L(P,f)<\varepsilon/j$.
Let $K=\{k:\ E_j\cap(a_k,a_{k+1})\neq\emptyset\}$.
For $k\in K$, pick $y_k\in E_j\cap(a_k,a_{k+1})$; then with $I=(a_k,a_{k+1})$,
\[
\sup_I f-\inf_I f>\tfrac1j,
\]
hence $M_k-m_k\ge\tfrac1j$. Therefore
\[
\sum_{k\in K}(a_{k+1}-a_k)
\le j\sum_{k\in K}(M_k-m_k)(a_{k+1}-a_k)
\le j\big(U(P,f)-L(P,f)\big)<\varepsilon.
\]
Thus $E_j\setminus\{a_0,\dots,a_n\}$ is covered by $\bigcup_{k\in K}(a_k,a_{k+1})$ of
total length $<\varepsilon$, so $E_j$ has measure zero.

\paragraph{(c) If $D_f$ has measure zero, then $f$ is Riemann integrable.}\mbox{}\par
Let $\varepsilon>0$ and choose open intervals $I_\ell$ with
$D_f\subset\bigcup_\ell I_\ell$ and $\sum_\ell |I_\ell|<\varepsilon$.
Put $F=[a,b]\setminus\bigcup_\ell I_\ell$.
Then $F$ is compact and $f|_F$ is uniformly continuous:
$\exists\delta>0$ such that $x,y\in F$, $|x-y|<\delta\Rightarrow |f(x)-f(y)|<\varepsilon$.

Refine a partition $P=\{a=a_0<\cdots<a_n=b\}$ so that
(i) $a_{k+1}-a_k<\delta$ for all $k$, and
(ii) every endpoint of every $I_\ell$ is a partition point.
Let $J=\{k:\ [a_k,a_{k+1}]\subset F\}$.
For $k\in J$,
\[
\sup_{[a_k,a_{k+1}]} f - \inf_{[a_k,a_{k+1}]} f < \varepsilon.
\]
For $k\notin J$ we have $[a_k,a_{k+1}]\subset I_\ell$ for some $\ell$, hence
$\sum_{k\notin J}(a_{k+1}-a_k)\le \sum_\ell |I_\ell|<\varepsilon$.
Let $M=\sup_{[a,b]}|f|$. Then
\[
\begin{aligned}
	U(P,f)-L(P,f)
	&=\sum_{k=0}^{n-1}\big(\sup_{[a_k,a_{k+1}]} f-\inf_{[a_k,a_{k+1}]} f\big)\,(a_{k+1}-a_k)\\
	&< \varepsilon\sum_{k\in J}(a_{k+1}-a_k)+2M\sum_{k\notin J}(a_{k+1}-a_k)\\
	&\le \varepsilon(b-a)+2M\,\varepsilon.
\end{aligned}
\]
Since $\varepsilon>0$ is arbitrary, $U(P,f)-L(P,f)$ can be made arbitrarily small; hence
$f$ is Riemann integrable on $[a,b]$.


\section*{Problem 1 (Quickies) — Statements and Theory}

\paragraph{Statements (exact).}
\begin{enumerate}[(a)]
	\item Use the equivalence of norms on a finite dimensional vector space to show that for each $n$, there is a constant $C$ such that the following holds: for every polynomial $p$ of degree $\le n$ there is $x_0 \in [0,1/n]$ such that
	\[
	|p(x)| \le C\,|p(x_0)| \qquad \text{for every } x \in [0,1].
	\]
	\item If $(X,d)$ is a metric space and $A$ is a non-empty subset of $X$, show that the distance from $x\in X$ to $A$ defined by
	\[
	\rho(x)=\inf_{y\in A} d(x,y)
	\]
	is a Lipschitz function on $X$ with Lipschitz constant $\le 1$.
	\item If $(x_n)$, $(y_n)$ are Cauchy sequences in a metric space $(X,d)$, show that $\big(d(x_n,y_n)\big)$ is convergent (in $\mathbb{R}$).
\end{enumerate}

\paragraph{Theory (tools to use).}
\begin{itemize}
	\item \textbf{Equivalence of norms on finite-dimensional spaces.}
	If $V$ is finite dimensional, any two norms $\|\cdot\|_\alpha,\|\cdot\|_\beta$ on $V$ are equivalent:  
	$\exists\,c,C>0$ with $c\|v\|_\beta \le \|v\|_\alpha \le C\|v\|_\beta$ for all $v\in V$.
	In particular, on $\mathcal P_n$ (polynomials of degree $\le n$), the sup-norms
	\[
	\|p\|_{[0,1]}:=\sup_{x\in[0,1]}|p(x)|, \qquad
	\|p\|_{[0,1/n]}:=\sup_{x\in[0,1/n]}|p(x)|
	\]
	are equivalent. The compactness of the unit sphere
	$S=\{p\in\mathcal P_n:\|p\|_{[0,1/n]}=1\}$ and continuity of $p\mapsto \|p\|_{[0,1]}$
	give a finite $C=\max_{p\in S}\|p\|_{[0,1]}$, yielding
	$\|p\|_{[0,1]}\le C\,\|p\|_{[0,1/n]}$.
	Since a nonzero polynomial cannot vanish on infinitely many points, each sup is attained:
	$\|p\|_{[0,1/n]}=|p(x_0)|$ for some $x_0\in[0,1/n]$.
	
	\item \textbf{Distance to a set is $1$-Lipschitz.}
	For any $x,x'\in X$ and $y\in A$,
	\[
	\rho(x)\le d(x,y)\le d(x,x')+d(x',y)\quad\Rightarrow\quad
	\rho(x)-\rho(x')\le d(x,x').
	\]
	Swapping $x,x'$ gives $|\rho(x)-\rho(x')|\le d(x,x')$.
	
	\item \textbf{Distances of Cauchy sequences form a Cauchy sequence in $\mathbb R$.}
	For all $m,n$,
	\[
	\big|d(x_n,y_n)-d(x_m,y_m)\big|
	\le d(x_n,x_m)+d(y_n,y_m)
	\]
	by the triangle inequality. Thus if $(x_n)$ and $(y_n)$ are Cauchy, then $(d(x_n,y_n))$ is Cauchy in
	$\mathbb R$ and hence convergent (completeness of $\mathbb R$).
\end{itemize}

\paragraph{(a) Polynomials on $[0,1]$ via equivalence of norms.}
Let $\mathcal P_n=\{p:\deg p\le n\}$. Define
\[
\|p\|_{[0,1]}:=\sup_{x\in[0,1]}|p(x)|,\qquad
\|p\|_{[0,1/n]}:=\sup_{x\in[0,1/n]}|p(x)|.
\]
These are norms on $\mathcal P_n$ (a nonzero polynomial cannot vanish on an infinite set).
Since $\mathcal P_n$ is finite dimensional, the norms are equivalent: there exists $C=C(n)$ with
\[
\|p\|_{[0,1]}\le C\,\|p\|_{[0,1/n]}\qquad(\forall\,p\in\mathcal P_n).
\]
Indeed, on the compact sphere $S=\{p:\|p\|_{[0,1/n]}=1\}$ the continuous map
$p\mapsto \|p\|_{[0,1]}$ attains a maximum $C$. For $p\ne0$,
$\|p\|_{[0,1]}\le C\|p\|_{[0,1/n]}$. Choosing $x_0\in[0,1/n]$ with
$|p(x_0)|=\|p\|_{[0,1/n]}$ yields $|p(x)|\le C\,|p(x_0)|$ for all $x\in[0,1]$.

\paragraph{(b) Distance to a set is $1$-Lipschitz.}
For nonempty $A\subset X$, define $\rho(x)=\inf_{y\in A} d(x,y)$. For any $x,x'\in X$ and $y\in A$,
$\rho(x)\le d(x,y)\le d(x,x')+d(x',y)$, hence $\rho(x)-\rho(x')\le d(x,x')$; swapping $x,x'$
gives $|\rho(x)-\rho(x')|\le d(x,x')$.

\paragraph{(c) Distances of Cauchy pairs converge.}
If $(x_n)$ and $(y_n)$ are Cauchy in $(X,d)$, then
\[
\big|d(x_n,y_n)-d(x_m,y_m)\big|\le d(x_n,x_m)+d(y_n,y_m)\xrightarrow[n,m\to\infty]{}0,
\]
so $(d(x_n,y_n))$ is Cauchy in $\mathbb R$ and hence convergent.

\bigskip
\section*{Problem 2 — Statements and Theory}

\paragraph{Statements (exact).}
\begin{enumerate}[(a)]
	\item Is the set $(1,2]$ an open subset of the metric space $\mathbb{R}$ with the metric $d(x,y)=|x-y|$? Is it closed?
	What if we replace the metric space $\mathbb{R}$ with the space $[0,2]$, the space $(1,3)$ or the space $(1,2]$, in each case with the metric $d$?
	\item Let $X$ be a set equipped with the discrete metric, and $Y$ any metric space. Describe all open subsets of $X$, closed subsets of $X$, compact subsets of $X$, Cauchy sequences in $X$, continuous functions $f:X\to Y$ and continuous functions $f:Y\to X$.
\end{enumerate}

\paragraph{Theory (tools to use).}
\begin{itemize}
	\item \textbf{Subspace topology / metric.}
	If $M\subset Z$ and $Z$ is metric with $d_Z$, the subspace metric on $M$ is
	$d_M(x,y)=d_Z(x,y)$. A set $U\subset M$ is open in $M$ iff $U=M\cap O$ for some open $O\subset Z$.
	Closedness is analogous.
	
	\item \textbf{Open/closed in the line.}
	With $d(x,y)=|x-y|$ on $\mathbb R$, an interval of the form $(\alpha,\beta)$ is open, $[\alpha,\beta]$ is closed; a half-open interval like $(1,2]$ is neither open (no ball around $1$ stays inside) nor closed (its complement is not open).
	
	\item \textbf{Discrete metric.}
	The discrete metric is $d(x,y)=\mathbf 1_{x\ne y}$. Then every singleton $\{x\}$ is open; hence \emph{every subset of $X$ is open and also closed}. Compact subsets are exactly the finite subsets (an infinite discrete space is not compact). A sequence in $X$ is Cauchy iff it is eventually constant (for $\varepsilon=\tfrac12$). Every function $f:X\to Y$ is continuous (domain is discrete). For maps into a discrete space, $f:Y\to X$ is continuous iff for each $x\in X$, the fibre $f^{-1}(\{x\})$ is open in $Y$ (equivalently, $f$ is \emph{locally constant}; in particular, if $Y$ is connected then $f$ is constant).
\end{itemize}

\paragraph{(a) Openness/closedness of $(1,2]$ under various ambient spaces.}
With the usual metric $d(x,y)=|x-y|$:
\begin{itemize}
\item In $\mathbb R$, $(1,2]$ is neither open nor closed:
it is not open because the point $2\in(1,2]$ has no $\varepsilon$-ball contained in $(1,2]$;
it is not closed because $1$ is a limit point (e.g. $1+1/n\to1$) not belonging to the set.

	\item As a subspace of $[0,2]$, $(1,2]$ is both open ($[0,2]\cap (1,\infty)$) and closed (complement $[0,1]$ is open in $[0,2]$).
\item As a subspace of $(1,3)$, the set $(1,2]$ is \emph{closed} because it can be written as the
intersection of the subspace with a closed subset of $\mathbb{R}$:
\[
(1,2]=(1,3)\cap(-\infty,2].
\]
Since $(-\infty,2]$ is closed in $\mathbb{R}$, the intersection is closed in the subspace $(1,3)$.
Equivalently, the complement $(1,3)\setminus(1,2]=(2,3)$ is open in $(1,3)$ (it is $(1,3)\cap(2,\infty)$),
so $(1,2]$ is closed but not open in the subspace.

	\item In the space $(1,2]$ itself, it is clopen (the whole space).
\end{itemize}

\paragraph{(b) Discrete metric on $X$.}
If $X$ has the discrete metric $d(x,y)=\mathbf 1_{x\ne y}$ and $Y$ is any metric space, then:
\begin{itemize}
	\item Every subset of $X$ is open (hence also closed).
	\item The compact subsets of $X$ are exactly the finite subsets.
	\item A sequence in $X$ is Cauchy iff it is eventually constant.
	\item Every function $f:X\to Y$ is continuous.
\end{itemize}

\section*{Problem 3. Metrics and small balls}

For each of the following sets $X$, determine whether the given function $d$ defines a metric on $X$.
In each case where the function does define a metric, describe the open ball $B_\varepsilon(x)$
for $x\in X$ and $\varepsilon>0$ small.
\begin{enumerate}[(i)]
	\item $X=\mathbb{R}^n$; \quad $d(x,y)=\min\{|x_1-y_1|,\ |x_2-y_2|,\ \ldots,\ |x_n-y_n|\}$.
	\item $X=\mathbb{Z}$; \quad $d(x,x)=0$, and, for $x\neq y$, $d(x,y)=2^{n}$ where $x-y=2^{n}a$ with
	$n$ a non–negative integer and $a$ an odd integer.
	\item $X$ is the set of functions from $\mathbb{N}$ to $\mathbb{N}$; \quad $d(f,f)=0$, and, for $f\neq g$,
	$d(f,g)=2^{-n}$ for the least $n$ such that $f(n)\neq g(n)$.
	\item $X=\mathbb{C}$; \quad $d(z,w)=|z-w|$ if $z$ and $w$ lie on the same line through the origin,
	and $d(z,w)=|z|+|w|$ otherwise.
\end{enumerate}

\paragraph{(i) $X=\mathbb{R}^n$, $d(x,y)=\min\{|x_1-y_1|,\dots,|x_n-y_n|\}$.}
Not a metric: identity fails since $x\neq y$ may share a coordinate, giving $d(x,y)=0$.

\paragraph{(ii) $X=\mathbb{Z}$, $d(x,x)=0$, and for $x\neq y$, $d(x,y)=2^{n}$ where $x-y=2^{n}a$ with $n\ge0$, $a$ odd.}
Let $\nu_2(k)$ be the exponent of $2$ in $k\neq0$. Then $d(x,y)=2^{\nu_2(x-y)}$.
For all $x,y,z$,
\[
\nu_2(x-z)\ge \min\{\nu_2(x-y),\nu_2(y-z)\}\quad\Rightarrow\quad
d(x,z)\le \max\{d(x,y),d(y,z)\}.
\]
Hence $d$ is an ultrametric. Small balls:
\[
B_\varepsilon(x)=\{y\in\mathbb Z:\ d(x,y)<\varepsilon\}=
\begin{cases}
	\{x\}, & 0<\varepsilon\le 1,\\[2mm]
	\{y:\ x-y\ \text{odd}\}, & 1<\varepsilon\le 2,\\[1mm]
	\{y:\ x-y\not\equiv 0\pmod{4}\}, & 2<\varepsilon\le 4,\\[1mm]
	\vdots\\
	\{y:\ \nu_2(x-y)\le k\}=\mathbb Z\setminus (x+2^{k+1}\mathbb Z), & 2^{k}<\varepsilon\le 2^{k+1}.
\end{cases}
\]

\paragraph{(iii) $X=\mathbb{N}^{\mathbb{N}}$ (functions $\mathbb{N}\to\mathbb{N}$).}
\[
d(f,f)=0,\qquad
\text{and for } f\neq g,\quad d(f,g)=2^{-n},
\quad \text{where } n=\min\{k:\ f(k)\neq g(k)\}.
\]

This is a \emph{metric} (in fact, an ultrametric).

\begin{itemize}
	\item If $n_{fg}$ is the first index where $f$ and $g$ differ, then for any $h$,
	\[
	n_{fh}\ge \min\{n_{fg},n_{gh}\},
	\]
	and hence
	\[
	d(f,h)=2^{-n_{fh}}\le \max\{2^{-n_{fg}},2^{-n_{gh}}\}
	=\max\{d(f,g),d(g,h)\}.
	\]
\end{itemize}

\noindent\hrulefill

\paragraph{Small open balls (cylinder sets).}
Distances take values $2^{-1},2^{-2},\dots$.
For $\varepsilon\in(2^{-(m+1)},2^{-m}]$ ($m\ge0$),
\[
B_\varepsilon(f)=\{g:\ g(k)=f(k)\ \text{for all } k\le m\}.
\]
Equivalently, with $N=\lfloor\log_2(1/\varepsilon)\rfloor$,
\[
B_\varepsilon(f)=\{g:\ g|_{\{1,\dots,N\}}=f|_{\{1,\dots,N\}}\}.
\]


\paragraph{(iv) $X=\mathbb C$.}
\[
d(z,w)=
\begin{cases}
	|z-w|, & \text{if $z,w$ lie on the same line through the origin (i.e.\ } w\in \mathbb R z),\\[2mm]
	|z|+|w|, & \text{otherwise.}
\end{cases}
\]

This is a \emph{metric}.

\begin{itemize}
	\item \textbf{Symmetry/positivity.} Clear; and $d(z,w)=0 \Rightarrow z=w$.
	\item \textbf{Triangle inequality.}  
	If $z,w$ are not colinear with the origin then $d(z,w)=|z|+|w|$.  
	For any $u\in\mathbb C$,
	\[
	d(z,u)\ge |z-u|,\qquad d(u,w)\ge |u-w|
	\]
	(if colinear we have equality with the Euclidean distance; otherwise the value is
	$|\,\cdot\,|+|\,\cdot\,|$, which dominates $|\,\cdot-\cdot\,|$ via
	$|u|+|w|\ge |\,u-w\,|$). Hence
	\[
	d(z,u)+d(u,w)\ \ge\ |z-u|+|u-w|\ \ge\ |z-w|.
	\]
	When $z,w$ are colinear, the rightmost term equals $d(z,w)=|z-w|$, giving the inequality.
	When $z,w$ are non–colinear, $d(z,w)=|z|+|w| \le d(z,u)+d(u,w)$ since in every case
	$d(v,u)\ge \big||v|-|u|\big|$ and at least one summand is $\ge |v|+|u|$.
\end{itemize}

\noindent\hrulefill

\paragraph{Small open balls.} Fix $z\in\mathbb C$ and $\varepsilon>0$.
\[
B_\varepsilon(z)=
\underbrace{\{\,w\in\mathbb C:\ w\in \mathbb R z\ \text{and}\ |w-z|<\varepsilon\,\}}_{\text{points on the same line as }z}
\ \ \cup\ \
\underbrace{\{\,w\notin \mathbb R z:\ |z|+|w|<\varepsilon\,\}}_{\text{off-line points}}.
\]

Thus:
\begin{itemize}
	\item If $\varepsilon\le |z|$: the second set is empty; $B_\varepsilon(z)$ is the open interval
	of length $2\varepsilon$ on the line $\mathbb R z$ centered at $z$.
	\item If $\varepsilon>|z|$: in addition you get the Euclidean disk
	$\{\,w:\ |w|<\varepsilon-|z|\,\}$ consisting of all off-line points sufficiently
	close to the origin.
\end{itemize}




\section*{Problem 4. Let $(X,d)$ be a metric space.}

\paragraph{(a) Exact statement.}
Show that the union of any collection of open subsets of $X$ must be open (regardless of whether
the collection is finite, countable or uncountable), and that the intersection of any finite
collection of open subsets is again open. Formulate and prove similar properties about the
closed subsets of $X$.

\paragraph{(b) Exact statement.}
Let $E$ be a subset of $X$. Show that there is a unique largest open subset $E^\circ$ of $X$
contained in $E$, i.e.\ a unique open subset $E^\circ$ of $X$ such that $E^\circ\subset E$ and if
$G$ is any open subset of $X$ with $G\subset E$ then $G\subset E^\circ$. The set $E^\circ$ is
called the \emph{interior} of $E$ in $X$. Show also that there is a unique smallest closed subset
$\overline{E}$ of $X$ containing $E$, i.e.\ a unique closed subset $\overline{E}$ of $X$ with
$E\subset \overline{E}$ and if $F$ is any closed subset of $X$ with $E\subset F$ then
$\overline{E}\subset F$. The set $\overline{E}$ is called the \emph{closure} of $E$ in $X$.

\paragraph{(c) Exact statement.}
Show that
\[
E^\circ=\{\,x\in X:\ B_\varepsilon(x)\subset E\ \text{for some }\varepsilon>0\,\}
\qquad\text{and}\qquad
\overline{E}=\{\,x\in X:\ x_n\to x\ \text{for some sequence }(x_n)\subset E\,\}.
\]

\bigskip
\noindent\textbf{Theory (tools and facts).}
\begin{itemize}
	\item \textbf{Open/closed set algebra.}
	\begin{itemize}
		\item Arbitrary unions of open sets are open; finite intersections of open sets are open.
		\item Arbitrary intersections of closed sets are closed; finite unions of closed sets are closed.
		\item Complement duality: $U$ open $\Leftrightarrow$ $X\setminus U$ closed.
	\end{itemize}
	
	\item \textbf{Open balls.} For $x\in X$ and $\varepsilon>0$,
	\[
	B_\varepsilon(x)=\{y\in X:\ d(x,y)<\varepsilon\}
	\]
	is open; every open set is a union of open balls.
	
	\item \textbf{Interior.}
	\[
	E^\circ=\bigcup\{\,G\subset X:\ G\ \text{open and } G\subset E\,\}
	\]
	is open, satisfies $E^\circ\subset E$, and is the largest open subset of $E$.
	Metric characterisation:
	\[
	x\in E^\circ\ \Longleftrightarrow\ \exists\,\varepsilon>0:\ B_\varepsilon(x)\subset E.
	\]
	
	\item \textbf{Closure.}
	\[
	\overline{E}=\bigcap\{\,F\subset X:\ F\ \text{closed and } E\subset F\,\}
	\]
	is closed, contains $E$, and is the smallest closed set containing $E$.
	Equivalent characterisations:
	\[
	x\in\overline{E}\ \Longleftrightarrow\ \forall\,\varepsilon>0,\ B_\varepsilon(x)\cap E\neq\varnothing
	\ \Longleftrightarrow\ \exists\ (x_n)\subset E\ \text{with } x_n\to x.
	\]
	
	\item \textbf{Limit points.} $x$ is a limit (accumulation) point of $E$ iff every ball
	$B_\varepsilon(x)$ intersects $E\setminus\{x\}$; then $x\in\overline{E}$.
\end{itemize}

\paragraph{(a) Unions/intersections of open/closed sets.}
\begin{itemize}
	\item \emph{Arbitrary union of open sets is open.}
	If $x\in\bigcup_{i\in I}U_i$ with each $U_i$ open, then $x\in U_j$ for some $j$,
	so $\exists\,\varepsilon>0$ with $B_\varepsilon(x)\subset U_j\subset\bigcup_{i\in I}U_i$.
	\item \emph{Finite intersection of open sets is open.}
	If $x\in\bigcap_{k=1}^m U_k$, pick $\varepsilon_k>0$ with $B_{\varepsilon_k}(x)\subset U_k$.
	Then $B_{\min_k \varepsilon_k}(x)\subset\bigcap_{k=1}^m U_k$.
	\item \emph{Closed sets (duals).}
	Arbitrary intersections of closed sets are closed, and finite unions of closed sets are closed
	(by complements of the two statements above).
\end{itemize}

\noindent\hrulefill

\paragraph{(b) Interior and closure via extremal properties.}
Define
\[
E^\circ=\bigcup\{G\subset X:\ G\ \text{open and}\ G\subset E\},\qquad
\overline E=\bigcap\{F\subset X:\ F\ \text{closed and}\ E\subset F\}.
\]
Then $E^\circ$ is open, $E^\circ\subset E$, and if $G$ is open with $G\subset E$ then $G\subset E^\circ$;
hence $E^\circ$ is the unique largest open subset of $X$ contained in $E$.
Also $\overline E$ is closed, $E\subset\overline E$, and if $F$ is closed with $E\subset F$,
then $\overline E\subset F$; hence $\overline E$ is the unique smallest closed subset of $X$
containing $E$.

\noindent\hrulefill

\paragraph{(c) Metric characterisations.}
\begin{align*}
	E^\circ&=\{x\in X:\ \exists\,\varepsilon>0\ \text{with}\ B_\varepsilon(x)\subset E\},\\
	\overline E&=\{x\in X:\ \exists\ (x_n)\subset E\ \text{with}\ x_n\to x\}.
\end{align*}
\emph{Proof of the first:} If $x\in E^\circ$, then $x$ lies in some open $G\subset E$, so
$B_\varepsilon(x)\subset G\subset E$ for some $\varepsilon>0$.
Conversely, if $B_\varepsilon(x)\subset E$, then $x$ lies in the open set $B_\varepsilon(x)\subset E$,
hence $x\in E^\circ$.

\emph{Proof of the second:} If $x\in\overline E$, then every ball $B_{1/n}(x)$ meets $E$; pick
$x_n\in E\cap B_{1/n}(x)$ to get $x_n\to x$. Conversely, if $x_n\in E$ and $x_n\to x$, then every
neighbourhood of $x$ contains some $x_n\in E$, so every ball around $x$ meets $E$, hence
$x\in\overline E$.

\section*{Problem 5}
Let $V$ be a normed space, $x\in V$ and $r>0$. Prove that the closure of the open ball
$B_r(x)$ is the closed ball $D_r(x)=\{y\in V:\|x-y\|\le r\}$.
Give an example showing that in a general metric space $(X,d)$ the closure of $B_r(x)$
need not equal $D_r(x)=\{y\in X:\ d(x,y)\le r\}$.

\paragraph{Solution.}
Since $y\mapsto \|x-y\|$ is continuous on $V$, the set
$D_r(x)=\{y:\|x-y\|\le r\}$ is closed and contains $B_r(x)=\{y:\|x-y\|<r\}$. Hence
\[
\overline{B_r(x)}\subset D_r(x).
\]
For the reverse inclusion, let $y\in D_r(x)$. If $\|x-y\|<r$ then $y\in B_r(x)\subset
\overline{B_r(x)}$. If $\|x-y\|=r$, define for $t\in(0,1)$
\[
y_t=(1-t)y+tx.
\]
Then
\[
\|x-y_t\|=\|(1-t)(x-y)\|=(1-t)\|x-y\|=(1-t)r<r,
\]
so $y_t\in B_r(x)$ and $y_t\to y$ as $t\downarrow 0$. Thus $y\in\overline{B_r(x)}$.
Therefore $D_r(x)\subset \overline{B_r(x)}$, and we conclude
\[
\overline{B_r(x)}=D_r(x).
\]

\paragraph{Counterexample in a general metric space.}
Let $X$ be any nonempty set with the discrete metric $d(u,v)=\mathbf 1_{u\ne v}$.
Take $r=1$ and $x\in X$. Then $B_1(x)=\{x\}$, hence $\overline{B_1(x)}=\{x\}$,
while
\[
D_1(x)=\{y\in X:\ d(x,y)\le 1\}=X.
\]
Thus $\overline{B_1(x)}\ne D_1(x)$ in this metric space.

\paragraph{Notation.}
For a metric space $(X,d)$, a center $x\in X$ and $r>0$,
\[
B_r(x):=\{y\in X:\ d(x,y)<r\},\qquad
D_r(x):=\{y\in X:\ d(x,y)\le r\}.
\]

\paragraph{Counterexample 1 (discrete metric).}
Let $X$ be any nonempty set and $d(u,v)=\mathbf 1_{u\ne v}$.
For $r=1$ and any $x\in X$,
\[
B_1(x)=\{x\},\qquad \overline{B_1(x)}=\{x\},\qquad
D_1(x)=X.
\]
Hence $\overline{B_1(x)}\neq D_1(x)$.

\noindent\hrulefill

\paragraph{Counterexample 2 (bounded metric on $\mathbb{R}$).}
Let $X=\mathbb{R}$ with $d(x,y)=\min\{1,|x-y|\}$.
For any $x\in\mathbb{R}$ and $r=1$,
\[
B_1(x)=\{y:\ |x-y|<1\}=(x-1,x+1),\qquad
\overline{B_1(x)}=[x-1,x+1],\qquad
D_1(x)=\{y:\ \min(1,|x-y|)\le 1\}=\mathbb{R}.
\]
Thus $\overline{B_1(x)}\neq D_1(x)$.

\noindent\hrulefill

\paragraph{Counterexample 3 (star graph metric).}
Let $X=\{x\}\cup\{y_n:n\in\mathbb{N}\}$ and define
\[
d(x,y_n)=1,\qquad d(y_n,y_m)=2\ (n\ne m),\qquad d(z,z)=0.
\]
This is a metric (shortest–path metric on a tree).
For $r=1$,
\[
B_1(x)=\{x\},\qquad \overline{B_1(x)}=\{x\},\qquad
D_1(x)=\{x\}\cup\{y_n:n\in\mathbb{N}\}.
\]
Hence $\overline{B_1(x)}\neq D_1(x)$.

\section*{Problem 6}
In lectures we proved that if $E$ is a compact subset of $\mathbb{R}^n$ (Euclidean metric), then any
continuous $f:E\to\mathbb{R}$ has bounded image. Prove the converse: if $E\subset\mathbb{R}^n$
and every continuous $f:E\to\mathbb{R}$ has bounded image, then $E$ is compact.

\paragraph{Solution.}
By Heine--Borel, it suffices to show that $E$ is closed and bounded.

\emph{Boundedness.} The function $x\mapsto \|x\|$ is continuous on $\mathbb{R}^n$; by hypothesis
its restriction to $E$ is bounded, so $E$ is bounded.

\emph{Closedness.} Suppose $E$ is not closed. Then there exists $p\in \overline{E}\setminus E$.
Define
\[
g:E\to\mathbb{R},\qquad g(x)=\frac{1}{\|x-p\|}.
\]
Since $p\notin E$, the denominator is positive on $E$, and $g$ is continuous as a composition of
continuous functions. Because $p\in\overline{E}$, there is a sequence $(x_k)\subset E$ with
$x_k\to p$, hence $\|x_k-p\|\to0$ and $g(x_k)\to\infty$, so $g$ is unbounded on $E$.
This contradicts the hypothesis that every continuous $E\to\mathbb{R}$ is bounded. Therefore
$E$ is closed.

Being both closed and bounded in $\mathbb{R}^n$, the set $E$ is compact by Heine--Borel.
\qed

\paragraph{Remarks and variants.}
\begin{itemize}
	\item The hypothesis “every continuous $f:E\to\mathbb R$ is bounded’’ is called
	\emph{pseudocompactness}. In metric spaces (in particular in $\mathbb R^n$),
	pseudocompact $\Longleftrightarrow$ compact; our proof used Heine--Borel.
	
	\item \emph{Alternative closedness proof.}
	If $E$ were not closed, pick $p\in\overline E\setminus E$ and consider the continuous
	distance-to-$p$ map $h(x)=\|x-p\|$. Then $h(E)\subset(0,\infty)$ is not closed in $\mathbb R$,
	so $1/h$ is unbounded on $E$—the same contradiction as above.
	
	\item \emph{Sequential compactness route.}
	From boundedness, every sequence in $E$ lies in a bounded subset of $\mathbb R^n$,
	hence has a convergent subsequence $(x_{n_k})\to q\in\mathbb R^n$.
	If $q\notin E$, the function $x\mapsto 1/\|x-q\|$ would be an unbounded continuous map on $E$,
	contradiction. Thus every sequence in $E$ has a convergent subsequence whose limit lies in $E$,
	so $E$ is sequentially compact, hence compact.
	
	\item \emph{Stronger converse (extreme values).}
	If, in addition, every continuous $f:E\to\mathbb R$ attains its maximum on $E$, then $E$ is
	compact: closedness follows as above, boundedness from $x\mapsto\|x\|$, and now Heine--Borel.
\end{itemize}

\section*{Problem 7. Topological vs.\ metric properties}

Each item makes sense for a metric space $(X,d)$. Decide which are \emph{topological}
(depend only on the open sets) and justify.

\begin{enumerate}[(i)]
	\item \textbf{Boundedness of a subset of $X$} — \emph{Not topological.}
	
	Example: On $\mathbb R$, the metrics $d(x,y)=|x-y|$ and $d'(x,y)=\min\{1,|x-y|\}$ generate the same
	topology. Then $(\mathbb R,d')$ is bounded (diameter $\le 1$) while $(\mathbb R,d)$ is unbounded.
	
	\item \textbf{Closedness of a subset of $X$} — \emph{Topological.}
	
	Closed sets are complements of open sets; hence determined by the topology alone.
	
	\item \textbf{“Closed and bounded” for a subset of $X$} — \emph{Not topological.}
	
	Closed is topological, bounded is not. With $d,d'$ as in (i), $\mathbb R$ is closed in itself and
	bounded in $(\mathbb R,d')$ but not in $(\mathbb R,d)$.
	
	\item \textbf{Total boundedness of $X$} — \emph{Not topological.}
	
	Let $h:\mathbb R\to(0,1)$ be a homeomorphism, and define $d_h(x,y)=|h(x)-h(y)|$. Then $d_h$
	induces the usual topology on $\mathbb R$, but $(\mathbb R,d_h)$ is isometric to $(0,1)$ and thus
	totally bounded, while $(\mathbb R,|\cdot|)$ is not.
	
	\item \textbf{Completeness of $X$} — \emph{Not topological.}
	
	Let $h:(0,1)\to\mathbb R$ be a homeomorphism and $\rho(x,y)=|h(x)-h(y)|$. Then $\rho$ induces the
	usual topology on $(0,1)$, but $(0,1,\rho)$ is complete (isometric to $\mathbb R$), whereas
	$(0,1,|\cdot|)$ is not complete.
	
	\item \textbf{“$X$ is complete and totally bounded$\,\,$”} — \emph{Topological (for metric spaces).}
	
	In metric spaces,
	\[
	\text{complete and totally bounded}\quad\Longleftrightarrow\quad\text{compact}.
	\]
	Compactness is a topological property; hence the conjunction is topological.
	(Neither completeness nor total boundedness alone is topological; see (iv)–(v).)
\end{enumerate}



\section*{Problem 8}
Use the Contraction Mapping Theorem to show that the equation $\cos x=x$ has a unique real
solution. Find this solution to reasonable accuracy (radians), and justify the claimed accuracy.

\paragraph{Solution.}
Set $g(x)=\cos x$. Any solution of $\cos x=x$ is a fixed point of $g$.

\emph{Contraction on $[-1,1]$.}
Since $|\cos x|\le 1$, any fixed point must lie in $[-1,1]$. On this interval
\[
|g'(x)|=|-\sin x|\le \sin 1=:L<1,
\]
and $g([-1,1])=[\cos 1,1]\subset [-1,1]$. Thus $g$ is a contraction on $[-1,1]$
with constant $L=\sin 1\approx 0.84$. By the Contraction Mapping Theorem there is a unique
fixed point $x^\ast\in[-1,1]$, and for any $x_0\in[-1,1]$ the iteration $x_{n+1}=\cos x_n$
converges to $x^\ast$.

\emph{Numerics.}
Starting from $x_0=1$ gives
\[
\begin{aligned}
	x_1&=\cos 1=0.5403023059,\quad
	x_2\approx 0.8575532158,\quad
	x_3\approx 0.6542897905,\quad
	x_4\approx 0.7934803587,\\
	x_5&\approx 0.7013687736,\ \ldots,\
	x_{10}\approx 0.7442373549,\
	x_{15}\approx 0.7383689990,\
	x_{20}\approx 0.7391843998.
\end{aligned}
\]
Hence $x^\ast\approx 0.7390851332$.

\emph{Error bound.}
On the smaller invariant interval $I=[0.73,0.75]$ one has
$g(I)\subset I$ and $L_I=\sup_{x\in I}|g'(x)|=\sin(0.75)\approx 0.6816$.
For a contraction,
\[
|x^\ast-x_n|\le \frac{L_I}{1-L_I}\,|x_n-x_{n-1}|.
\]
At $n=20$, $|x_{20}-x_{19}|\approx 2.47\times 10^{-4}$, yielding
$|x^\ast-x_{20}|\lesssim 5.3\times 10^{-4}$. Thus $x^\ast=0.73918\pm 0.00053$
(correct to three decimal places). Additional iterations improve the bound,
e.g.\ $x^\ast\approx 0.73908513$ (8 d.p.).

\section*{Problem 9 - Weighted sup norm; Picard operator as a contraction}
Let $I=[0,R]$ be an interval and let $C(I)$ be the space of continuous functions on $I$.
Show that, for any $\alpha\in\mathbb{R}$, we may define a norm by
\[
\|f\|_\alpha=\sup_{x\in I}|f(x)|e^{-\alpha x},
\]
and that the norm $\|\cdot\|_\alpha$ is Lipschitz equivalent to the uniform norm
\[
\|f\|=\sup_{x\in I}|f(x)|.
\]
Now suppose that $\phi:\mathbb{R}^2\to\mathbb{R}$ is continuous and Lipschitz in the second variable.
Consider the map $T$ from $C(I)$ to itself sending $f$ to
\[
(Tf)(x)=y_0+\int_0^x \phi(t,f(t))\,dt.
\]
Give an example to show that $T$ need not be a contraction under the uniform norm.
Show, however, that $T$ is a contraction under the norm $\|\cdot\|_\alpha$ for some $\alpha$, and hence deduce that the differential equation
\[
f'(x)=\phi(x,f(x))
\]
has a unique solution on $I$ satisfying $f(0)=y_0$.

\bigskip
\hrule
\bigskip

\paragraph{Theory.}

\textbf{1. Weighted supremum norm.}
For $\alpha\in\mathbb{R}$ define
\[
\|f\|_\alpha=\sup_{x\in[0,R]}|f(x)|e^{-\alpha x}.
\]
This is a norm on $C(I)$. It is \emph{Lipschitz equivalent} to the usual supremum norm
\(\|f\|_\infty=\sup_{x\in[0,R]}|f(x)|\)
because on $[0,R]$, the weight $e^{-\alpha x}$ satisfies
\[
e^{-|\alpha|R}\le e^{-\alpha x}\le e^{|\alpha|R}.
\]
Hence for all $f\in C(I)$,
\[
e^{-|\alpha|R}\|f\|_\infty \le \|f\|_\alpha \le e^{|\alpha|R}\|f\|_\infty,
\]
showing equivalence of norms.

\bigskip

\textbf{2. The Picard operator.}
Assume $\phi:\mathbb{R}^2\to\mathbb{R}$ is continuous and Lipschitz in its second variable:
there exists $L>0$ such that
\[
|\phi(t,y)-\phi(t,z)|\le L|y-z| \quad \forall\, t,y,z.
\]
Then the operator
\[
(Tf)(x)=y_0+\int_0^x \phi(t,f(t))\,dt
\]
maps $C(I)$ to itself, and a fixed point $f=Tf$ is precisely a solution of the integral form
\[
f(x)=y_0+\int_0^x \phi(t,f(t))\,dt,
\]
equivalent to the differential equation $f'(x)=\phi(x,f(x))$ with $f(0)=y_0$.

\bigskip

\textbf{3. Contraction property.}
For $f,g\in C(I)$,
\[
|(Tf-Tg)(x)|\le \int_0^x |\phi(t,f(t))-\phi(t,g(t))|\,dt
\le L\int_0^x |f(t)-g(t)|\,dt.
\]
Thus
\[
\|Tf-Tg\|_\infty\le LR\,\|f-g\|_\infty.
\]
If $LR\ge1$, $T$ is not a contraction in the uniform norm.

\bigskip

\textbf{4. Weighted norm case.}
For the weighted norm $\|\cdot\|_\alpha$,
\[
e^{-\alpha x}|(Tf-Tg)(x)|
\le L\int_0^x e^{-\alpha(x-t)}|f(t)-g(t)|e^{-\alpha t}\,dt
\le L\|f-g\|_\alpha \int_0^x e^{-\alpha s}\,ds
\le \frac{L}{\alpha}\|f-g\|_\alpha.
\]
Hence $\|Tf-Tg\|_\alpha\le \frac{L}{\alpha}\|f-g\|_\alpha$.  
Choosing $\alpha>L$ ensures that $T$ is a contraction in $\|\cdot\|_\alpha$.

\bigskip

\textbf{5. Existence and uniqueness.}
By the Banach Fixed Point Theorem, $T$ admits a unique fixed point $f\in C(I)$.
Thus the differential equation $f'(x)=\phi(x,f(x))$, $f(0)=y_0$, has a unique continuous solution on $I$.
\qed

Let $I=[0,R]$ and $C(I)$ be the space of continuous functions on $I$. For $\alpha\in\mathbb R$ set
\[
\|f\|_\alpha=\sup_{x\in I}|f(x)|e^{-\alpha x},\qquad \|f\|_\infty=\sup_{x\in I}|f(x)|.
\]
\paragraph{(a) Equivalence of norms.}
Since $e^{-\alpha x}\in [e^{-|\alpha|R},e^{|\alpha|R}]$ on $I$, for all $f\in C(I)$,
\[
e^{-|\alpha|R}\,\|f\|_\infty \ \le\ \|f\|_\alpha \ \le\ e^{|\alpha|R}\,\|f\|_\infty.
\]
Hence $\|\cdot\|_\alpha$ is a norm Lipschitz–equivalent to the uniform norm.

\paragraph{(b) Picard operator.}
Assume $\phi:\mathbb R^2\to\mathbb R$ is continuous and Lipschitz in its second variable:
$|\phi(t,y)-\phi(t,z)|\le L|y-z|$ for all $t,y,z$.
Define $T:C(I)\to C(I)$ by
\[
(Tf)(x)=y_0+\int_0^x \phi(t,f(t))\,dt .
\]

\emph{Not a contraction in $\|\cdot\|_\infty$ in general.}
For $f,g\in C(I)$,
\[
\|Tf-Tg\|_\infty\le L\!\int_0^R |f(t)-g(t)|\,dt \le LR\,\|f-g\|_\infty .
\]
If $LR\ge1$, the Lipschitz factor is not $<1$. Example: $I=[0,2]$, $\phi(t,y)=y$ ($L=1$) gives
$\|Tf-Tg\|_\infty\le 2\,\|f-g\|_\infty$.

\emph{$T$ is a contraction in $\|\cdot\|_\alpha$ for suitable $\alpha>0$.}
For $x\in[0,R]$,
\[
\begin{aligned}
	e^{-\alpha x}\,|(Tf-Tg)(x)|
	&\le L\int_0^x e^{-\alpha(x-t)}\big(|f(t)-g(t)|e^{-\alpha t}\big)\,dt \\
	&\le L\|f-g\|_\alpha \int_0^x e^{-\alpha s}\,ds
	\ \le\ \frac{L}{\alpha}\,\|f-g\|_\alpha .
\end{aligned}
\]
Therefore $\|Tf-Tg\|_\alpha \le \frac{L}{\alpha}\|f-g\|_\alpha$. Choosing $\alpha>L$ makes $T$ a contraction.
By Banach’s Fixed Point Theorem, $T$ has a unique fixed point $f\in C(I)$, i.e.
\[
f(x)=y_0+\int_0^x \phi(t,f(t))\,dt ,
\]
which is the unique solution on $I$ of the IVP $f'(x)=\phi(x,f(x))$, $f(0)=y_0$.
\qed

\section*{Problem 10}
Let $(X,d)$ be a non-empty complete metric space. Suppose $f:X\to X$ is a contraction and
$g:X\to X$ is a function which commutes with $f$, i.e.\ $f(g(x))=g(f(x))$ for all $x\in X$.
Show that $g$ has a fixed point. Must this fixed point be unique?

\bigskip
\paragraph{Theory.}
\begin{itemize}
	\item \textbf{Banach Fixed Point Theorem (Contraction Mapping).}
	If $(X,d)$ is complete and $f:X\to X$ is a contraction, i.e.
	$\exists\,0<L<1$ such that $d(fx,fy)\le L\,d(x,y)$ for all $x,y\in X$, then
	\begin{enumerate}
		\item $f$ has a \emph{unique} fixed point $p\in X$;
		\item for any $x_0\in X$, the Picard iteration $x_{n+1}=f(x_n)$ converges to $p$ with the a priori bound
		$d(x_n,p)\le \dfrac{L^n}{1-L}\,d(x_1,x_0)$.
	\end{enumerate}
	
	\item \textbf{Commutation and fixed points.}
	If $f\circ g=g\circ f$ and $p$ is a fixed point of $f$ (i.e.\ $f(p)=p$), then
	\[
	f\big(g(p)\big)=g\big(f(p)\big)=g(p),
	\]
	so $g(p)$ is also a fixed point of $f$. Hence $g$ maps the fixed-point set $\mathrm{Fix}(f)$
	into itself.
	
	\item \textbf{Consequence when $f$ is a contraction.}
	Since $\mathrm{Fix}(f)=\{p\}$ is a singleton by Banach’s theorem, the previous item forces
	$g(p)=p$. Thus $p$ is a fixed point of $g$.
	
	\item \textbf{Uniqueness warning.}
	The fixed point of $g$ need not be unique. Example: on $X=\mathbb{R}$, let
	$f(x)=x/2$ (a contraction) and $g(x)=x$ (identity). Then $f\circ g=g\circ f$,
	$g$ has \emph{every} point as a fixed point.
\end{itemize}

\subsection*{Examples for Problem 10 (commuting with a contraction)}

Let $(X,d)$ be complete and let the contraction be the “halving” map $f=\tfrac12\,\mathrm{Id}$ (so $f(x)=\tfrac12 x$ on linear spaces, or $f(h)=\tfrac12 h$ pointwise on function spaces). In each case below $f\circ g=g\circ f$, and we list the fixed points of $g$.

\paragraph{A. $g$ has \emph{many} fixed points (non-unique).}
\[
X=\mathbb{R},\quad f(x)=\tfrac12 x,\quad g(x)=x.
\]
Then $f\circ g=g\circ f=\tfrac12 x$, and $\mathrm{Fix}(g)=\mathbb{R}$ (every point is fixed).

\paragraph{B. $g$ has a \emph{unique} fixed point but is not a contraction.}
\[
X=\mathbb{R}^2,\quad f(x)=\tfrac12 x,\quad 
g(x)=Ax,\ \ A=\begin{pmatrix}2&0\\[2pt]0&\tfrac12\end{pmatrix}.
\]
Clearly $f$ commutes with any linear $A$, so $f\circ g=g\circ f$.  
Fixed points solve $Ax=x$, i.e.\ $(2-1)x_1=0$, $(\tfrac12-1)x_2=0$, hence $\mathrm{Fix}(g)=\{(0,0)\}$.
Moreover $\|A\|=2>1$, so $g$ is not a contraction.

\paragraph{C. $g$ has a \emph{finite} fixed-point set.}
\[
X=\mathbb{R},\quad f(x)=\tfrac12 x,\quad g(x)=-x.
\]
Then $f(g(x))=\tfrac12(-x)=-(\tfrac12 x)=g(f(x))$.  
Fixed points satisfy $-x=x\Rightarrow x=0$, so $\mathrm{Fix}(g)=\{0\}$.

\paragraph{D. $g$ has \emph{uncountably many} fixed points.}
\[
X=\mathbb{R}^2,\quad f(x)=\tfrac12 x,\quad g=\Pi_L \text{ (orthogonal projection onto a line $L$).}
\]
Since $\tfrac12\Pi_L=\Pi_L\tfrac12$, we have $f\circ g=g\circ f$.  
Fixed points of $g$ are precisely the points on $L$, so $\mathrm{Fix}(g)=L$ (uncountable).

\paragraph{E. Nonlinear $g$ on a function space; $g$ has several fixed points.}
\[
X=C([0,1]),\quad \|h\|_\infty\ \text{metric},\quad f(h)=\tfrac12 h,\quad g(h)=h^3\ \text{(pointwise cube)}.
\]
Then $f(g(h))=\tfrac12 h^3=g(f(h))=(\tfrac12 h)^3\cdot 4=\tfrac12 h^3$.  
Fixed points satisfy $h^3=h$ pointwise. By continuity, $h$ must be a constant in $\{-1,0,1\}$, so
\[
\mathrm{Fix}(g)=\{-\mathbf{1},\ \mathbf{0},\ \mathbf{1}\}.
\]

\paragraph{F. Why commutation is essential (counterexample without it).}
\[
X=\mathbb{R},\quad f(x)=\tfrac12 x\ \text{(contraction)},\quad g(x)=x+1.
\]
Here $g$ has no fixed point, and $f\circ g(x)=\tfrac12(x+1)\neq \tfrac12 x+1=g\circ f(x)$, so the
commutation hypothesis fails.

\section*{Problem 11 — Strict contractions on complete vs compact spaces; non-expansive maps}

\noindent\textbf{Task.}
\begin{enumerate}\setlength{\itemsep}{4pt}
	\item Give a non-empty complete metric space $(X,d)$ and a map $f:X\to X$ with
	$d(f(x),f(y))<d(x,y)$ for all $x\neq y$, but \emph{no fixed point}.
	\item Now assume $X\subset\mathbb{R}^n$ is non-empty and \emph{compact} (Euclidean metric).
	Show that such an $f$ \emph{must} have a fixed point (and decide uniqueness).
	\item If $g:X\to X$ only satisfies $d(g(x),g(y))\le d(x,y)$ for all $x,y$, must $g$ have a fixed point?
\end{enumerate}

\bigskip
\hrule
\bigskip

\noindent\textbf{Theory you’ll use.}
\begin{itemize}\setlength{\itemsep}{6pt}
	\item A metric space is \emph{complete} if every Cauchy sequence converges.
	\item On a compact metric space, any continuous map $h:X\to\mathbb{R}$ attains its minimum
	(apply this to $h(x)=d(f(x),x)$).
	\item \emph{Edelstein’s fixed point theorem.} On a compact metric space, a map
	$f$ with $d(f(x),f(y))<d(x,y)$ for all $x\neq y$ has a \emph{unique} fixed point.
	\item If $a_n=d(x_{n+1},x_n)$ and $a_{n+1}<a_n$, then $(a_n)$ is strictly decreasing
	and converges to some $L\ge0$.
\end{itemize}

\bigskip
\hrule
\bigskip

\noindent\textbf{Part 1 — Strictly contractive, complete, no fixed point (example).}

\medskip
\noindent Let $X=(0,\infty)$ and define
\[
d(x,y):=\big|\ln x-\ln y\big|\qquad (x,y>0).
\]
Then $(X,d)$ is complete because $(0,\infty)$ with this metric is isometric to
$(\mathbb{R},|\cdot|)$ via $x\mapsto \ln x$.

\medskip
\noindent Define $f(x)=x+1$. For $x\neq y$, by the mean value theorem there exist
$\xi$ between $x+1$ and $y+1$ and $\eta$ between $x$ and $y$ with
\[
\big|\ln(x+1)-\ln(y+1)\big|=\frac{|x-y|}{\xi},
\qquad
\big|\ln x-\ln y\big|=\frac{|x-y|}{\eta}.
\]
Since $\xi>\eta>0$, we have
\[
d\big(f(x),f(y)\big)=\frac{|x-y|}{\xi}\;<\;\frac{|x-y|}{\eta}=d(x,y)
\quad\text{for all }x\ne y.
\]
But $f(x)=x+1$ has no fixed point. Hence we have a complete space with a strict
contraction (in the pointwise sense) without a fixed point.

\bigskip
\hrule
\bigskip

\noindent\textbf{Part 2 — Compact $X\subset\mathbb{R}^n$: existence and uniqueness.}

\medskip
\noindent Let $X\subset\mathbb{R}^n$ be non-empty and compact (Euclidean metric).
Suppose $f:X\to X$ satisfies $d(f(x),f(y))<d(x,y)$ for all $x\ne y$.

\medskip
\noindent\emph{Existence.}
Pick any $x_0\in X$ and iterate $x_{k+1}=f(x_k)$. Set $a_k=d(x_{k+1},x_k)$.
Then
\[
a_{k+1}=d\big(f(x_{k+1}),f(x_k)\big)\;<\;d(x_{k+1},x_k)=a_k,
\]
so $(a_k)$ decreases to some $L\ge0$. By compactness, there is a subsequence
$x_{n_j}\to p\in X$. Then $x_{n_j+1}=f(x_{n_j})\to f(p)$, and by continuity of $d$,
\[
L=\lim_{j\to\infty} d\big(x_{n_j+1},x_{n_j}\big)
= d\big(f(p),p\big).
\]
If $L>0$, applying the strict inequality to $(p,f(p))$ gives
\[
d\big(f(f(p)),f(p)\big)\;<\;d\big(p,f(p)\big)=L,
\]
but the left-hand side equals $\lim_j d(x_{n_j+2},x_{n_j+1})=L$, a contradiction.
Hence $L=0$ and $f(p)=p$.

\medskip
\noindent\emph{Uniqueness.}
If $p\ne q$ are fixed points, then
\[
d\big(f(p),f(q)\big)=d(p,q)\;<\;d(p,q),
\]
impossible. Therefore the fixed point is \emph{unique}.

\bigskip
\hrule
\bigskip

\noindent\textbf{Part 3 — Non-expansive maps ($\le$) need not have fixed points.}

\medskip
\noindent Let $X=S^1\subset\mathbb{R}^2$ be the unit circle with the Euclidean metric,
and let $g$ be rotation by a nonzero angle. Then
\[
d\big(g(x),g(y)\big)=d(x,y)\le d(x,y)\quad\text{for all }x,y\in S^1,
\]
but $g$ has no fixed point. Thus the condition $d(g(x),g(y))\le d(x,y)$
does \emph{not} ensure a fixed point, even on compact sets.

\bigskip
\hrule
\bigskip

\section*{Problem 12}
(a) Let $B=B_1(0)\subset\mathbb{R}^n$ be the closed unit ball and
$F:[0,1]\times B\to\mathbb{R}^n$ be continuous. Suppose there is $K>0$ with
\[
\|F(t,x)-F(t,y)\|\le K\|x-y\|\qquad \forall\,t\in[0,1],\ \forall\,x,y\in B.
\]
Given $x_1,x_2\in\mathbb{R}^n$, by Picard--Lindelöf there exist
$\varepsilon\in(0,1]$ and differentiable $f_1,f_2:[0,\varepsilon]\to B$ such that
$\dot f_j(t)=F(t,f_j(t))$ and $f_j(0)=x_j$ for $j=1,2$. Show that
\[
\|f_1(t)-f_2(t)\|\le \|x_1-x_2\|\,e^{Kt}\qquad (0\le t\le\varepsilon).
\]

(b) Now assume a Hölder condition: there are $K>0$ and $\alpha\in(0,1)$ with
\[
\|F(t,x)-F(t,y)\|\le K\|x-y\|^\alpha\qquad \forall\,t\in[0,1],\ \forall\,x,y\in B.
\]
Show for the same $f_1,f_2$ that
\[
\|f_1(t)-f_2(t)\|
\le \Big(\,\|x_1-x_2\|^{1-\alpha}+(1-\alpha)Kt\,\Big)^{\frac{1}{1-\alpha}}
\qquad (0\le t\le\varepsilon).
\]
What does this imply about the set of solutions to
$\dot f=F(t,f(t))$, $f(0)=x_0$ on some interval $[0,\varepsilon]$?

\bigskip
\hrule
\bigskip

\paragraph{Theory.}
\begin{itemize}\setlength{\itemsep}{6pt}
	\item \textbf{Grönwall (differential form).} If $u\ge0$ is absolutely continuous and
	$u'(t)\le Ku(t)$ a.e., then $u(t)\le u(0)e^{Kt}$.
	\item \textbf{Osgood / generalized Grönwall.} If $u'(t)\le Ku(t)^\alpha$ with
	$0<\alpha<1$, then
	$u(t)^{1-\alpha}\le u(0)^{1-\alpha}+(1-\alpha)Kt$.
\end{itemize}

\bigskip

\paragraph{Solution.}
Let $w(t)=f_1(t)-f_2(t)$ and $u(t)=\|w(t)\|$.
Since $\dot w(t)=F(t,f_1(t))-F(t,f_2(t))$ and $u$ is absolutely continuous,
\[
u'(t)\le \|\dot w(t)\|=\|F(t,f_1(t))-F(t,f_2(t))\|.
\]

\emph{(a) Lipschitz case.}
From the hypothesis, $u'(t)\le Ku(t)$.
By Grönwall, $u(t)\le u(0)e^{Kt}=\|x_1-x_2\|e^{Kt}$.
In particular, if $x_1=x_2$ then $u(0)=0$ and $u\equiv0$ (uniqueness).

\emph{(b) Hölder case.}
Now $u'(t)\le K u(t)^\alpha$.
For $u(t)>0$, divide by $u^\alpha$ and integrate:
\[
\frac{d}{dt}\big(u^{1-\alpha}(t)\big)=(1-\alpha)u(t)^{-\alpha}u'(t)
\le (1-\alpha)K.
\]
Hence
\[
u(t)^{1-\alpha}\le u(0)^{1-\alpha}+(1-\alpha)Kt,
\]
i.e.
\[
\boxed{\ \|f_1(t)-f_2(t)\|
	\le \Big(\,\|x_1-x_2\|^{1-\alpha}+(1-\alpha)Kt\,\Big)^{\frac{1}{1-\alpha}} \ }.
\]

\paragraph{Consequence for uniqueness.}
When $x_1=x_2$, the bound reads
$\|f_1(t)-f_2(t)\|\le\big((1-\alpha)Kt\big)^{1/(1-\alpha)}$, which does not force
$\|f_1(t)-f_2(t)\|\equiv0$. Thus for $\alpha\in(0,1)$ uniqueness may fail.
(Example: $\dot y=|y|^\alpha$ with $y(0)=0$ admits both $y\equiv0$ and
$y(t)=\big((1-\alpha)Kt\big)^{1/(1-\alpha)}$.)
\qed


\section*{Problem 13$^{\ast}$}
Let $(X,d)$ be a non-empty complete metric space and let $f:X\to X$ be a function such that
for each positive integer $n$ we have
\begin{enumerate}[(i)]
	\item if $d(x,y)<n+1$ then $d(f(x),f(y))<n$;
	\item if $d(x,y)<1/n$ then $d(f(x),f(y))<1/(n+1)$.
\end{enumerate}
Must $f$ have a fixed point?

\bigskip
\hrule
\bigskip

\paragraph{Theory.}
(Boyd--Wong / Matkowski contraction.)
If there is a function $\phi:[0,\infty)\to[0,\infty)$ such that
$\phi(r)<r$ for all $r>0$ and
\[
d(fx,fy)\le \phi\big(d(x,y)\big)\qquad\forall\,x,y\in X,
\]
then $f$ has a unique fixed point on any complete metric space.

\bigskip
\paragraph{Solution.}
Define a piecewise function
\[
\phi(t)=
\begin{cases}
	\lceil t\rceil - 1, & t\ge 1,\\[3pt]
	\dfrac{1}{n+1}, & \text{if } \dfrac{1}{n+1}<t\le \dfrac{1}{n},\quad n\in\mathbb N.
\end{cases}
\]
Then $\phi(t)<t$ for all $t>0$.  


By (i) and (ii), we have $d(fx,fy)\le \phi(d(x,y))$.


Thus $f$ is a generalized contraction in the sense of Boyd--Wong, so
$f$ has a (necessarily unique) fixed point.

Uniqueness also follows since
$d(f(p),f(q))<d(p,q)$ is impossible if both $p,q$ are fixed.

\bigskip
\hrule
\bigskip

\section*{Problem 14$^{\ast}$}
Let $K$ be a compact subset of $\mathbb{R}$ and let $p\in K$. Construct a metric $d$ on
$K_1:=K\setminus\{p\}$ such that $(K_1,d)$ is complete and the topology generated by $d$ on $K_1$
is the same as the topology generated by the Euclidean metric on $K_1$.

\paragraph{Construction.}
For $x,y\in K_1$ set
\[
d(x,y)\;=\;|x-y|\;+\;\Bigl|\frac{1}{x-p}-\frac{1}{y-p}\Bigr|.
\]

\paragraph{Same topology.}
The map $h:K_1\to\mathbb{R}^2$, $h(x)=(x,\,(x-p)^{-1})$ is an embedding.
The metric $d$ is the pullback of the product metric on $\mathbb{R}^2$ via $h$, hence induces
on $K_1$ the subspace topology from $\mathbb{R}$ (the first coordinate), i.e. the Euclidean subspace topology.

\paragraph{Completeness.}
Let $(x_n)$ be $d$-Cauchy in $K_1$. Then $(x_n)$ is Cauchy in $|\cdot|$, and
$\big(\,(x_n-p)^{-1}\big)$ is Cauchy in $|\cdot|$, hence bounded. Therefore there exists $\delta>0$
with $|x_n-p|\ge \delta$ for all $n$.

Because $K$ is compact, $(x_n)$ has a subsequence $x_{n_j}\to x\in K$; the uniform gap to $p$
implies $x\neq p$, so $x\in K_1$. 

Since $(x_n)$ is Cauchy in $|\cdot|$, the whole sequence
converges to the same $x$.

 Moreover, $(x_n)$ is Cauchy for the second summand, hence
$\bigl((x_n-p)^{-1}\bigr)$ converges in $\mathbb{R}$ to $(x-p)^{-1}$. Thus $d(x_n,x)\to 0$,
showing $(K_1,d)$ is complete. \hfill$\Box$


\section*{Problem 15$^{\ast}$}
\textbf{(a)} Let $(X,d)$ be a totally bounded metric space. Show that any sequence $(x_k)$ in $X$ has a Cauchy subsequence.

\textbf{(b)} Show that a metric space is compact if and only if it is complete and totally bounded.

\bigskip

\paragraph{Theory and definitions.}
\begin{itemize}\setlength{\itemsep}{6pt}
	\item \emph{Totally bounded}: for every $\varepsilon>0$, there exist $x_1,\dots,x_N\in X$ such that
	\[
	X \subset \bigcup_{i=1}^N B(x_i,\varepsilon).
	\]
	\item \emph{Cauchy sequence}: $(y_n)$ is Cauchy if for every $\varepsilon>0$ there exists $N$ such that
	$d(y_m,y_n)<\varepsilon$ for all $m,n\ge N$.
	\item \emph{Complete}: every Cauchy sequence converges in $X$.
	\item \emph{Compact}: every open cover has a finite subcover (equivalently in metric spaces: sequentially compact).
\end{itemize}

\bigskip

\paragraph{(a) Existence of a Cauchy subsequence in a totally bounded space.}
Fix the scale sequence $\varepsilon_m := 2^{-m}$ for $m\in\mathbb{N}$.

\medskip

\emph{Step 1 (first thinning).} By total boundedness, $X$ can be covered by finitely many balls of radius $\varepsilon_1$.
One of these balls contains infinitely many terms of $(x_k)$. Pass to that infinite subsequence and relabel it $(x^{(1)}_k)$.

\medskip

\emph{Step 2 (iterated thinning).} Again by total boundedness, cover $X$ by finitely many balls of radius $\varepsilon_2$.
One such ball contains infinitely many terms of $(x^{(1)}_k)$. Pass to that infinite subsequence and relabel it $(x^{(2)}_k)$.

\medskip

\emph{Step 3 (continue).} Proceed inductively: after step $m$ we have an infinite subsequence
$(x^{(m)}_k)$ contained in a single ball of radius $\varepsilon_m$.

\medskip

\emph{Step 4 (diagonal choice).} Define the diagonal subsequence $y_m := x^{(m)}_m$.
Then $y_m$ lies inside a ball of radius $\varepsilon_m$ that contains all but finitely many terms of the previous stage.
In particular,
\[
d(y_{m+1},y_m) \;<\; \varepsilon_m \qquad \text{for all } m\in\mathbb{N}.
\]

\medskip

\emph{Step 5 (Cauchy property).} Let $\varepsilon>0$ and choose $M$ with $\varepsilon_M<\varepsilon/2$.
For $n>m\ge M$, the points $y_m,y_n$ both lie in a ball of radius $\varepsilon_m\le \varepsilon_M<\varepsilon/2$,
so $d(y_m,y_n)\le 2\varepsilon_M<\varepsilon$. Hence $(y_m)$ is Cauchy.

\medskip

Thus every sequence in a totally bounded metric space admits a Cauchy subsequence.

\bigskip

\paragraph{(b) Compact $\Longleftrightarrow$ complete $+$ totally bounded.}

\emph{($\Rightarrow$) Compact $\Rightarrow$ complete and totally bounded.}
Compact metric spaces are sequentially compact. A Cauchy sequence has a convergent subsequence;
Cauchy $+$ convergent subsequence forces the whole sequence to converge, hence completeness.
Total boundedness follows since otherwise some $\varepsilon$-net would require infinitely many balls,
contradicting compactness (use a standard Lebesgue number/covering argument).

\medskip

\emph{($\Leftarrow$) Complete and totally bounded $\Rightarrow$ compact.}
Let $(x_k)$ be any sequence in $X$. By (a), it has a Cauchy subsequence; by completeness this subsequence converges in $X$.
Hence $X$ is sequentially compact. In metric spaces, sequential compactness is equivalent to compactness.
Therefore $X$ is compact.

\hfill$\Box$

\section*{Problem 1 (Quickies)}

\paragraph{(a)}
Let $F:[0,1]\times\mathbb{R}^m\to\mathbb{R}$ be continuous and 
$a=(a_0,\dots,a_{m-1})\in\mathbb{R}^m$.
Suppose that $F$ is uniformly Lipschitz in the $\mathbb{R}^m$ variables near $a$, 
i.e. for some $K>0$ and an open $U\subset\mathbb{R}^m$ containing $a$,
\[
|F(t,x)-F(t,y)|\le K\|x-y\|\quad\text{for all }t\in[0,1],\;x,y\in U.
\]
Show using the Picard–Lindelöf theorem that there exists $\varepsilon>0$ such that
the $m$th order IVP
\[
f^{(m)}(t)=F\bigl(t,f(t),f^{(1)}(t),\ldots,f^{(m-1)}(t)\bigr),\qquad
f^{(j)}(0)=a_j\;(0\le j\le m-1)
\]
has a unique $C^m$ solution $f:[0,\varepsilon)\to\mathbb{R}$.

\paragraph{(b)}
Let $f:\mathbb{R}^2\to\mathbb{R}$ and $a\in\mathbb{R}^2$.
If all directional derivatives $D_u f(a)$ exist and depend linearly on $u$, 
does it follow that $f$ is differentiable at $a$?

\paragraph{(c)}
Let $f:\mathbb{R}^2\to\mathbb{R}$.
If $f$ is differentiable at $0$ and its partial derivatives exist in a neighbourhood of $0$,
does it follow that one partial derivative is continuous at $0$?

\paragraph{(d)}
Let $f:[a,b]\to\mathbb{R}^2$ be continuous and differentiable on $(a,b)$.
Does there exist $c\in(a,b)$ such that
\[
f(b)-f(a)=f'(c)(b-a)?
\]

\bigskip
\hrule
\bigskip

\paragraph{Theory.}
\begin{itemize}\setlength{\itemsep}{6pt}
	\item \textbf{Picard--Lindelöf theorem.}
	If $F(t,x)$ is continuous in $t$ and Lipschitz in $x$, then $x'(t)=F(t,x(t))$, $x(0)=a$ 
	has a unique local solution.
	\item Linearity of directional derivatives does not imply differentiability unless continuity of the derivative is also guaranteed.
	\item Differentiability at a point does not imply that partial derivatives are continuous.
	\item The mean value theorem does not extend in this form to vector-valued functions.
\end{itemize}

\bigskip

\paragraph{Solution.}

\subparagraph{(a)} 
Set $y_1=f,\,y_2=f',\,\dots,\,y_m=f^{(m-1)}$.
Then the system
\[
\begin{cases}
	y_1'=y_2,\\
	y_2'=y_3,\\
	\vdots\\
	y_{m-1}'=y_m,\\
	y_m'=F(t,y_1,\dots,y_m)
\end{cases}
\]
defines $Y' = G(t,Y)$ with $G$ continuous and Lipschitz in $Y$.
By Picard–Lindelöf, there exists $\varepsilon>0$ and a unique $C^1$ function 
$Y:[0,\varepsilon)\to\mathbb{R}^m$.
Hence $f=y_1\in C^m$ solves the given $m$th order problem uniquely.

\subparagraph{(b)} 
No. For example,
\[
f(x,y)=
\begin{cases}
	\dfrac{x^2y}{x^2+y^2}, & (x,y)\ne(0,0),\\[4pt]
	0,&(x,y)=(0,0),
\end{cases}
\]
has all directional derivatives at $0$ and they depend linearly on $u$, 
but $f$ is not differentiable at $0$.

\subparagraph{(c)} 
No. For instance,
\[
f(x,y)=
\begin{cases}
	x^2\sin(1/x), & x\ne0,\\[4pt]
	0,&x=0,
\end{cases}
\]
is differentiable at $0$ but $f_x$ is not continuous there.

\subparagraph{(d)} 
No. The scalar mean value theorem fails for vector-valued maps. Example:
\[
f(t)=(\cos t,\sin t),\quad a=0,\;b=2\pi.
\]
Then $f(b)-f(a)=(0,0)$ while $f'(t)=(-\sin t,\cos t)$ never vanishes,
so there is no $c$ with $f(b)-f(a)=f'(c)(b-a)$.
\hfill$\Box$

\section*{Problem 2}
Let $F:[a,b]\times\mathbb{R}^n\to\mathbb{R}^n$ be continuous, $x_0\in\mathbb{R}^n$ and $R>0$.
Suppose that
\[
\sup_{[a,b]\times \overline{B_R(x_0)}}\|F\|\le \frac{R}{b-a}
\quad\text{and}\quad
\|F(t,x)-F(t,y)\|\le K\|x-y\|
\]
for some $K$ and all $t\in[a,b]$, $x,y\in\overline{B_R(x_0)}$.
We showed in lecture that for each $t_0\in[a,b]$ there is a unique
$f\in C([a,b];\overline{B_R(x_0)})$ solving
\[
f(t)=x_0+\int_{t_0}^{t}F\big(s,f(s)\big)\,ds,\qquad t\in[a,b].
\]
Show that this $f$ is in fact the unique function in $C([a,b];\mathbb{R}^n)$ solving the integral equation.

\paragraph{Theory.}
\begin{itemize}\setlength{\itemsep}{6pt}
	\item (Picard--Lindel\"of / Gr\"onwall) If $F$ is continuous in $t$ and Lipschitz in $x$, the integral equation has a unique solution as long as the solution stays in the domain.
	\item A uniform bound on $\|F\|$ controls how far a solution can move away from $x_0$ by integrating the bound.
\end{itemize}

\paragraph{Solution.}
Let $g\in C([a,b];\mathbb{R}^n)$ satisfy $g(t)=x_0+\int_{t_0}^{t}F(s,g(s))\,ds$.
Define
\[
\Lambda^+ := \{\,t\in[t_0,b] : \|g(\sigma)-x_0\|\le R\ \ \forall\,\sigma\in[t_0,t]\,\}.
\]
Then $\Lambda^+$ is nonempty and closed. If $\tau:=\sup\Lambda^+<b$, continuity of $g$ and the bound on
$\|F\|$ give
\[
\|g(\tau)-x_0\|
\le \int_{t_0}^{\tau}\|F(s,g(s))\|\,ds
\le (\tau-t_0)\sup\|F\|
\le (\tau-t_0)\frac{R}{b-a}<R,
\]
so by continuity there exists $\eta>0$ with $\|g(\sigma)-x_0\|<R$ for $\sigma\in[\tau,\tau+\eta]$,
contradicting the definition of $\tau$. Hence $\Lambda^+=[t_0,b]$.
A symmetric argument for $[a,t_0]$ shows $g([a,b])\subset\overline{B_R(x_0)}$.
By the lecture result and Lipschitz uniqueness, $g=f$.
Therefore $f$ is the unique solution in $C([a,b];\mathbb{R}^n)$.
\hfill$\Box$

\bigskip\hrule\bigskip

\section*{Problem 3}
Let $U\subset\mathbb{R}^n$ be open, $f:U\to\mathbb{R}^n$ and $a\in U$.
A differentiable curve passing through $a$ is a differentiable map
$\gamma:(-1,1)\to\mathbb{R}^n$ with $\gamma(0)=a$.
If $f\circ\gamma$ is differentiable at $0$ for every differentiable curve $\gamma$ passing through $a$,
does it follow that $f$ is differentiable at $a$?

\paragraph{Theory.}
A pointwise form of Boman's theorem (the $C^1$ case) states that if $f\circ\gamma$ is differentiable
at $0$ for every $C^1$ curve $\gamma$ with $\gamma(0)=a$, then $f$ is Fr\'echet differentiable at $a$.

\paragraph{Solution.}
Yes. Define $L:\mathbb{R}^n\to\mathbb{R}^n$ by $Lv:=\frac{d}{dt}\big|_{t=0}(f\circ\gamma_v)(t)$ with
$\gamma_v(t)=a+tv$. Linearity of $L$ follows from additivity and homogeneity of derivatives along
the curves $\gamma_{u+v}$ and $\gamma_{\alpha v}$.
For $h\to0$, consider the $C^1$ curve $\gamma_h(t)=a+th$ (for small $t$). Differentiability of
$f\circ\gamma_h$ at $0$ yields
\[
f(a+h)-f(a)-Lh=o(\|h\|),
\]
which is the Fr\'echet differentiability of $f$ at $a$ with $Df(a)=L$.
\hfill$\Box$

\bigskip\hrule\bigskip

\section*{Problem 4}
Define $f:\mathbb{R}^3\to\mathbb{R}^2$ by
\[
f(x,y,z)=\big(e^{x+y+z},\ \cos(z^2y)\big).
\]
Without making use of partial derivatives, show that $f$ is everywhere differentiable and find $Df(a)$ at each $a\in\mathbb{R}^3$.
Then compute the partial derivatives of $f$ and, using appropriate results on partial derivatives, give an alternative proof.

\paragraph{Theory.}
Sums, products and compositions of $C^1$ maps are $C^1$; the chain rule gives
$D(\phi\circ g)(a)=D\phi(g(a))\circ Dg(a)$.

\paragraph{Solution (without partials).}
Write $f=(f_1,f_2)$ with $f_1=e^{\,\ell}$ where $\ell(x,y,z)=x+y+z$ is linear, and
$f_2=\cos\circ h$ where $h(x,y,z)=yz^2$.
Hence $f$ is $C^1$ everywhere and, for $a=(x,y,z)$ and $u=(u_x,u_y,u_z)$,
\[
Df_1(a)[u]=e^{x+y+z}(u_x+u_y+u_z),\qquad
Dh(a)[u]=z^2u_y+2zy\,u_z,
\]
\[
Df_2(a)[u]=-\sin(z^2y)\,(z^2u_y+2zy\,u_z).
\]
Thus the Jacobian is
\[
Df(a)=
\begin{pmatrix}
	e^{x+y+z} & e^{x+y+z} & e^{x+y+z}\\[3pt]
	0 & -\,z^2\sin(z^2y) & -\,2zy\sin(z^2y)
\end{pmatrix}.
\]

\paragraph{Alternative proof (via partial derivatives).}
\[
\partial_x f_1=\partial_y f_1=\partial_z f_1=e^{x+y+z},\qquad
\partial_x f_2=0,\ \ \partial_y f_2=-z^2\sin(z^2y),\ \ \partial_z f_2=-2zy\sin(z^2y).
\]
These partials are continuous, hence $f\in C^1$, and the above Jacobian agrees with the matrix of partials.
\hfill$\Box$

\section*{Problem 5 -- Normalization map $f(x)=x/\|x\|$}

\paragraph{Statement.}
Let $f:\mathbb{R}^3\to\mathbb{R}^3$ be given by
\[
f(x)=
\begin{cases}
	\dfrac{x}{\|x\|}, & x\ne0,\\[4pt]
	0,&x=0.
\end{cases}
\]
Show that $f$ is differentiable for $x\ne0$ and
\[
Df(x)(h)=\frac{h}{\|x\|}-\frac{x(x\cdot h)}{\|x\|^3},
\]
and that $Df(x)(h)\perp x$.

\paragraph{Theory.}
The map $x\mapsto\|x\|$ is $C^\infty$ on $\mathbb{R}^3\setminus\{0\}$.
The derivative of $x\mapsto x/\|x\|$ projects $h$ onto the tangent plane of the sphere $S^2$.

\paragraph{Solution.}
For $x\ne0$,
\[
Df(x)(h)=\frac{d}{dt}\bigg|_{t=0}\frac{x+th}{\|x+th\|}
=\frac{h}{\|x\|}-\frac{x(x\cdot h)}{\|x\|^3}.
\]
At $x=0$, $f$ is not continuous since $\|f(x)\|=1$ for $x\ne0$ but $f(0)=0$.
Hence $f$ is not differentiable at $0$.

\paragraph{Orthogonality.}
\[
x\cdot Df(x)(h)=\frac{x\cdot h}{\|x\|}-\frac{\|x\|^2(x\cdot h)}{\|x\|^3}=0.
\]
Thus $Df(x)(h)\perp x$. Geometrically, the derivative gives the projection onto the tangent space of the unit sphere at $f(x)$.

\paragraph{Example.}
For $x=(1,0,0)$, $Df(x)(h)=h-(h_1,0,0)=(0,h_2,h_3)$, tangent to $S^2$ at $(1,0,0)$.
\hfill$\Box$

\bigskip\hrule\bigskip

\section*{Problem 6 -- Differentiability of $|x||y|$ and $\dfrac{xy}{\sqrt{x^2+y^2}}$}

\paragraph{(a) Function $f(x,y)=|x||y|$.}
Partial derivatives:
\[
f_x(x,y)=\operatorname{sgn}(x)|y|,\qquad
f_y(x,y)=|x|\operatorname{sgn}(y).
\]
Hence $f$ is differentiable wherever $x\ne0$ and $y\ne0$.
On the axes, $\operatorname{sgn}$ is discontinuous, so $f$ fails to be differentiable there.

\paragraph{(b) Function $g(x,y)=\dfrac{xy}{\sqrt{x^2+y^2}}$, $g(0,0)=0$.}
In polar form, $x=r\cos\theta$, $y=r\sin\theta$:
\[
g(x,y)=r\cos\theta\sin\theta=\tfrac{r}{2}\sin(2\theta).
\]
Then $g\to0$ as $r\to0$, so $g$ is continuous at $0$.
However,
\[
\frac{|g(x,y)|}{\sqrt{x^2+y^2}}=|\cos\theta\sin\theta|
\]
depends on $\theta$, hence the derivative at $0$ does not exist.
Thus $g$ is differentiable everywhere except at $(0,0)$.
\hfill$\Box$

\bigskip\hrule\bigskip

\section*{Problem 7 -- Separate continuity vs. joint continuity}

\paragraph{Statement.}
Let $f:U\subset\mathbb{R}^2\to\mathbb{R}$ be such that
$f(\cdot,y)$ and $f(x,\cdot)$ are continuous for each fixed $y$ and $x$.
Show that $f$ need not be jointly continuous.
If $f$ is Lipschitz in each variable separately, deduce that $f$ is continuous.

\paragraph{Theory.}
Separate continuity does not imply joint continuity.
If each section is Lipschitz with bounded constants, $f$ becomes globally Lipschitz, hence jointly continuous.

\paragraph{Solution.}
A standard counterexample:
\[
f(x,y)=
\begin{cases}
	\dfrac{xy}{x^2+y^2}, & (x,y)\ne(0,0),\\[4pt]
	0,&(x,y)=(0,0).
\end{cases}
\]
For fixed $x$ or $y$, $f$ is continuous, but along $y=x$,
$f(x,x)=1/2\ne0$, so $f$ is not continuous at $(0,0)$.

If $f$ satisfies
\[
|f(x_1,y)-f(x_2,y)|\le L|x_1-x_2|,\qquad
|f(x,y_1)-f(x,y_2)|\le M|y_1-y_2|,
\]
then for all $(x_1,y_1),(x_2,y_2)\in U$,
\[
|f(x_1,y_1)-f(x_2,y_2)|\le L|x_1-x_2|+M|y_1-y_2|,
\]
so $f$ is globally Lipschitz and therefore continuous.

\paragraph{Examples.}
\begin{itemize}
	\item $f(x,y)=x+y$ is Lipschitz in both variables and continuous.
	\item $f(x,y)=|x|+|y|$ is separately Lipschitz and globally continuous.
\end{itemize}
\hfill$\Box$

\section*{Problem 8 -- Chain rule for higher derivatives}

\paragraph{Statement.}
Let $f,g:\mathbb{R}\to\mathbb{R}$ be twice continuously differentiable.
Show that
\[
(f\circ g)''(x) = f''(g(x))(g'(x))^2 + f'(g(x))\,g''(x),
\]
and find the expression for the third derivative.

\paragraph{Theory.}
The chain rule gives $(f\circ g)'(x)=f'(g(x))g'(x)$.
Applying the product rule again yields higher-order formulas,
which generalize to Faà di Bruno’s formula for arbitrary order.

\paragraph{Solution.}
Differentiate:
\[
(f\circ g)'(x)=f'(g(x))g'(x),
\]
then
\[
(f\circ g)''(x)=f''(g(x))(g'(x))^2+f'(g(x))g''(x).
\]
Differentiating once more,
\[
(f\circ g)'''(x)=f'''(g(x))(g'(x))^3
+3f''(g(x))g'(x)g''(x)
+f'(g(x))g'''(x).
\]
This pattern continues, involving combinations of derivatives of $f$ and $g$.

\paragraph{Example.}
For $f(u)=u^2$, $g(x)=e^x$, we have $(f\circ g)(x)=e^{2x}$.
Then
\[
(f\circ g)''(x)=4e^{2x},
\]
and by the formula,
$f''(g(x))(g'(x))^2+f'(g(x))g''(x)=2(e^x)^2+2e^x e^x=4e^{2x}$, confirming the result.
\hfill$\Box$

\bigskip\hrule\bigskip

\section*{Problem 9 -- Weighted sup norm; Picard operator as a contraction}
Let $F:[a,b]\times\mathbb{R}^n\to\mathbb{R}^n$ be continuous and Lipschitz in $x$
with constant $L$.
Define $(Tf)(t)=x_0+\int_a^t F(s,f(s))\,ds$ on $C([a,b];\mathbb{R}^n)$.
Under the weighted norm
\[
\|f\|_\alpha=\sup_{t\in[a,b]}e^{-\alpha(t-a)}\|f(t)\|,
\]
we have
\[
\|Tf-Tg\|_\alpha\le \frac{L}{\alpha}\|f-g\|_\alpha.
\]
Thus $T$ is a contraction whenever $\alpha>L$.

\paragraph{Example.}
If $F(t,x)=x$ on $[0,2]$, then $L=1$.
In the standard sup norm, $L(b-a)=2>1$, so $T$ is not a contraction.
With $\alpha=2$, $\frac{L}{\alpha}=\frac12<1$, so $T$ is a contraction.
\hfill$\Box$

\bigskip\hrule\bigskip

\section*{Problem 10 -- Inverse Function Theorem}

\paragraph{Statement.}
Let $f:\mathbb{R}^n\to\mathbb{R}^n$ be continuously differentiable near $a$,
and suppose $Df(a)$ is invertible.
Show that there exist neighborhoods $U\ni a$ and $V\ni f(a)$ such that
$f:U\to V$ is a diffeomorphism.

\paragraph{Theory.}
If $Df(a)$ is invertible, then near $a$ the linearization
$f(x)\approx f(a)+Df(a)(x-a)$ behaves like a bijection.
The Banach fixed-point theorem applied to
$T_y(x)=x-Df(a)^{-1}(f(x)-y)$ proves local invertibility.

\paragraph{Solution.}
Let $A=Df(a)$.  For $y$ near $f(a)$ define $T_y(x)=x-A^{-1}(f(x)-y)$.
If $Df(x)$ is close to $A$ near $a$, then $T_y$ is a contraction.
By the contraction mapping theorem, there is a unique $x$ with $T_y(x)=x$, i.e. $f(x)=y$.
Thus $f$ has a $C^1$ local inverse near $a$.
Differentiating $f(f^{-1}(y))=y$ gives
\[
Df^{-1}(y)=[Df(f^{-1}(y))]^{-1}.
\]

\paragraph{Examples.}
\begin{enumerate}\setlength{\itemsep}{3pt}
	\item $f(x)=x^3$ on $\mathbb{R}$: $Df(x)=3x^2$, invertible for $x\ne0$.  
	Locally invertible near $x=1$ with inverse $f^{-1}(y)=y^{1/3}$.
	\item $f(x,y)=(e^x\cos y,\,e^x\sin y)$ on $\mathbb{R}^2$:
	$\det Df(x,y)=e^{2x}\ne0$, hence locally invertible at every point.
\end{enumerate}
\hfill$\Box$

% --- Optional once in your preamble for quick vertical gaps ---
% \newcommand{\gap}{\vspace{6pt}}

\section*{Problem 7 — Separate continuity vs.\ joint continuity; Lipschitz criterion}

\textbf{Statement.}\gap

Let $f:U\subset\mathbb{R}^2\to\mathbb{R}$ satisfy:
for each fixed $y$, $f(\cdot,y)$ is continuous; for each fixed $x$, $f(x,\cdot)$ is continuous.
\begin{enumerate}\setlength{\itemsep}{6pt}
	\item Give an example where $f$ is \emph{not} jointly continuous on $U$.
	\item Suppose there exist $L,M\ge0$ such that
	\[
	|f(x_1,y)-f(x_2,y)|\le L|x_1-x_2|, \qquad
	|f(x,y_1)-f(x,y_2)|\le M|y_1-y_2|
	\]
	for all relevant points in $U$ (constants independent of the other variable).
	Show that $f$ is jointly continuous on $U$.
	\item Deduce: if $D_1 f$ exists and is bounded on $U$, and $f(x,\cdot)$ is continuous for each fixed $x$,
	then $f$ is continuous on $U$.
\end{enumerate}

\bigskip\hrule\bigskip

\textbf{Theory \& Solution.}\gap

\textbf{(1) Counterexample: separate $\not\Rightarrow$ joint).}\;
\[
f(x,y)=
\begin{cases}
	\dfrac{xy}{x^2+y^2}, & (x,y)\neq(0,0),\\[4pt]
	0,&(x,y)=(0,0).
\end{cases}
\]
For fixed $x$ or fixed $y$, $f\to 0$ at the origin; along $y=x$ we have $f(x,x)=\tfrac12$, so $f$ is not continuous at $(0,0)$.

\medskip

\textbf{(2) Uniform one–variable Lipschitz $\Rightarrow$ joint continuity).}\;
For $(x_1,y_1),(x_2,y_2)$,
\[
\begin{aligned}
	|f(x_1,y_1)-f(x_2,y_2)|
	&\le |f(x_1,y_1)-f(x_2,y_1)| + |f(x_2,y_1)-f(x_2,y_2)|\\
	&\le L|x_1-x_2| + M|y_1-y_2|,
\end{aligned}
\]
hence $f$ is globally Lipschitz on $U$ and therefore continuous.

\medskip

\textbf{(3) Bounded $D_1 f$ gives the needed uniform Lipschitz in $x$).}\;
If $|D_1 f|\le L$ on $U$ and $y$ is fixed, the 1D mean value theorem in $x$ yields
$|f(x_1,y)-f(x_2,y)|\le L|x_1-x_2|$ (constant independent of $y$).
Combine with continuity in $y$ and apply (2). \hfill$\square$

\bigskip\hrule\bigskip

\textbf{Examples.}\gap

\begin{itemize}\setlength{\itemsep}{6pt}
	\item $f(x,y)=|x|\,\sin y$: here $D_1 f=\operatorname{sgn}(x)\sin y$ is bounded; $f(x,\cdot)$ is continuous. Conclusion: $f$ is continuous on $U$.
	\item The counterexample above shows separate continuity alone is insufficient.
\end{itemize}

% ==========================================================

\section*{Problem 8 — One partial continuous near $a$ + the other exists at $a$ $\Rightarrow$ differentiability at $a$}

\textbf{Statement.}\gap

Let $f:\mathbb{R}^2\to\mathbb{R}$ and $a=(x_0,y_0)\in\mathbb{R}^2$.
Assume $D_1 f$ exists on an open ball $B(a,r)$ and is continuous at $a$, and $D_2 f(a)$ exists.
Show that $f$ is differentiable at $a$.

\bigskip\hrule\bigskip

\textbf{Theory \& Solution.}\gap

\textbf{(Increment decomposition and line integration).}\;
For $h=(h_1,h_2)$ small,
\[
\begin{aligned}
	f(x_0+h_1,y_0+h_2)-f(x_0,y_0)
	&=\big[f(x_0+h_1,y_0+h_2)-f(x_0,y_0+h_2)\big] \\
	&\quad + \big[f(x_0,y_0+h_2)-f(x_0,y_0)\big] \\
	&= \int_{0}^{h_1} D_1 f(x_0+t,y_0+h_2)\,dt \;+\; \big(f(x_0,y_0+h_2)-f(x_0,y_0)\big).
\end{aligned}
\]
Continuity of $D_1 f$ at $a$ gives
\[
\int_{0}^{h_1} D_1 f(x_0+t,y_0+h_2)\,dt
= D_1 f(a)\,h_1 + o(\|h\|).
\]
Since $D_2 f(a)$ exists,
\[
f(x_0,y_0+h_2)-f(x_0,y_0)= D_2 f(a)\,h_2 + o(\|h\|).
\]
Therefore
\[
f(a+h)-f(a) = D_1 f(a)\,h_1 + D_2 f(a)\,h_2 + o(\|h\|),
\]
so $f$ is Fr\'echet differentiable at $a$ with gradient $(D_1 f(a),D_2 f(a))$. \hfill$\square$

\bigskip\hrule\bigskip

\textbf{Example.}\gap

$f(x,y)=x^2|y|$: $D_1 f=2x|y|$ is continuous; $D_2 f(0,0)=0$ exists; hence $f$ is differentiable at $(0,0)$.

% ==========================================================

\section*{Problem 9 — Operator norm: basic and useful properties}

\textbf{Statement.}\gap

With Euclidean norms, let $\|A\|_{op}=\sup_{\|x\|=1}\|Ax\|$ for $A\in\mathcal{L}(\mathbb{R}^n;\mathbb{R}^m)$.
Prove:
\begin{enumerate}\setlength{\itemsep}{6pt}
	\item $\displaystyle \|A\|_{op}=\sup_{x\neq 0}\frac{\|Ax\|}{\|x\|}$.
	\item For $A\in\mathcal{L}(\mathbb{R}^n;\mathbb{R}^m)$, $B\in\mathcal{L}(\mathbb{R}^m;\mathbb{R}^p)$,
	\ $\|B\circ A\|_{op}\le \|B\|_{op}\,\|A\|_{op}$.
	\item If $A\in\mathcal{L}(\mathbb{R}^n;\mathbb{R})$, then $\exists\,a\in\mathbb{R}^n$ with $Ax=\langle a,x\rangle$ and $\|A\|_{op}=\|a\|$.
	\item If $A\in\mathcal{L}(\mathbb{R};\mathbb{R}^m)$, then $\exists\,a\in\mathbb{R}^m$ with $Ax=xa$ and $\|A\|_{op}=\|a\|$.
	\item If $(A_{ij})$ is the matrix of $A$ in the standard bases, then
	\[
	\frac{1}{\sqrt{n}}\Big(\sum_{j=1}^n\sum_{i=1}^m A_{ij}^2\Big)^{1/2}
	\ \le\ \|A\|_{op}\ \le\
	\Big(\sum_{j=1}^n\sum_{i=1}^m A_{ij}^2\Big)^{1/2},
	\]
	and the right inequality is an equality iff $A=0$ or $\operatorname{rank}(A)=1$.
\end{enumerate}

\bigskip\hrule\bigskip

\textbf{Theory \& Solution.}\gap

\textbf{(1) Homogeneity).}\;
$\|Ax\|/\|x\|$ depends only on $x/\|x\|$; taking $x$ on the unit sphere yields the same supremum.

\medskip

\textbf{(2) Submultiplicativity).}\;
For $\|x\|=1$, $\|B(Ax)\|\le \|B\|_{op}\|Ax\|\le \|B\|_{op}\|A\|_{op}$.
Taking the supremum gives $\|B\circ A\|_{op}\le \|B\|_{op}\|A\|_{op}$.

\medskip

\textbf{(3) Linear functionals as inner products).}\;
Finite-dimensional Riesz: $\exists a$ with $Ax=\langle a,x\rangle$ and by Cauchy–Schwarz
$\|Ax\|\le \|a\|\,\|x\|$, hence $\|A\|_{op}=\|a\|$.

\medskip

\textbf{(4) Maps $\mathbb{R}\to\mathbb{R}^m$).}\;
Let $a:=A(1)$. Then $Ax=A(x\cdot 1)=xa$ and
$\|Ax\|=\|x a\|=|x|\,\|a\|$, so $\|A\|_{op}=\|a\|$.

\medskip

\textbf{(5) Frobenius vs.\ operator norm; bounds and equality).}\;
Let $\|A\|_{F} := \big(\sum_{i,j} A_{ij}^2\big)^{1/2}$. For any unit $x$,
\[
\|Ax\|^2=\sum_{i=1}^{m}\Big(\sum_{j=1}^{n} A_{ij}x_j\Big)^2
\ \le\ \sum_{i=1}^{m}\Big(\sum_{j=1}^{n} A_{ij}^2\Big)\Big(\sum_{j=1}^{n}x_j^2\Big)
\ =\ \|A\|_F^2,
\]
so $\|A\|_{op}\le \|A\|_F$.
For the lower bound, take $x=\frac{1}{\sqrt{n}}(1,\dots,1)$ to obtain
$\|Ax\|^2\ge \tfrac{1}{n}\sum_{i,j}A_{ij}^2$, hence $\|A\|_{op}\ge \|A\|_F/\sqrt{n}$.
Equality $\|A\|_{op}=\|A\|_F$ holds iff $A$ has exactly one nonzero singular value,
i.e.\ $A=0$ or $\operatorname{rank}(A)=1$. \hfill$\square$

\bigskip\hrule\bigskip

\textbf{Examples.}\gap

\begin{itemize}\setlength{\itemsep}{6pt}
	\item Rank–one matrix $A=uv^\top$: $\|A\|_{op}=\|u\|\,\|v\|=\|A\|_F$ (right equality case).
	\item Identity $A=I_n$: $\|A\|_{op}=1$, $\|A\|_F=\sqrt{n}$; the bounds read
	$1\le 1 \le \sqrt{n}$.
\end{itemize}

\bigskip\hrule\bigskip

% --- Optional once in your preamble for quick vertical gaps ---
% \newcommand{\gap}{\vspace{6pt}}

\section*{Problem 10 — Global diffeomorphism under $\|Df - I\| \le \mu < 1$}

\textbf{Statement.}\gap

Let $f:\mathbb{R}^n \to \mathbb{R}^n$ be $C^1$ and assume
\[
\|Df(x) - I\| \le \mu \quad \text{for some } \mu\in(0,1) \text{ and all } x.
\]
Show:
\begin{enumerate}\setlength{\itemsep}{6pt}
	\item $f$ maps open sets to open sets.
	\item For all $x,y$,
	\[
	(1-\mu)\,\|x-y\| \;\le\; \|f(x)-f(y)\| \;\le\; (1+\mu)\,\|x-y\|.
	\]
	Hence $f$ is injective and $f(\mathbb{R}^n)$ is closed.
	\item Conclude that $f$ is a diffeomorphism of $\mathbb{R}^n$.
	\item Explain why the conclusion can fail if one only assumes $\|Df(x)-I\|\le 1$.
\end{enumerate}

\bigskip\hrule\bigskip

\textbf{Theory \& Solution.}\gap

\textbf{(Two–sided Lipschitz via line integration).}\;
Let $h:=y-x$ and $g(t):=f(x+th)$. Then
\[
f(y)-f(x)=\int_{0}^{1} Df(x+th)\,h\,dt
= h + \int_{0}^{1} (Df(x+th)-I)\,h\,dt.
\]
Hence
\[
\|f(y)-f(x)-h\|
\;\le\; \int_{0}^{1} \|Df(x+th)-I\|\,\|h\|\,dt
\;\le\; \mu\,\|h\|.
\]
By the triangle inequality,
\[
(1-\mu)\,\|h\|\ \le\ \|f(y)-f(x)\|\ \le\ (1+\mu)\,\|h\|.
\]
\emph{Injectivity} follows from the lower bound (if $f(x)=f(y)$ then $x=y$).

\medskip

\textbf{(Local diffeomorphism).}\;
Since $\|Df(x)-I\|<1$, each $Df(x)$ is invertible (Neumann series).
By the inverse function theorem, $f$ is a local $C^1$ diffeomorphism and therefore
\emph{maps open sets to open sets}.

\medskip

\textbf{(Closed range).}\;
If $f(x_k)\to y$ and $(x_k)$ is bounded, the lower Lipschitz bound gives
$\|x_k-x_\ell\|\le \tfrac{1}{1-\mu}\,\|f(x_k)-f(x_\ell)\|$, so $(x_k)$ is Cauchy
and converges to some $x$ with $f(x)=y$. Thus $f(\mathbb{R}^n)$ is closed.

\medskip

\textbf{(Surjectivity and global diffeo).}\;
The image $f(\mathbb{R}^n)$ is nonempty, open (local diffeo) and closed (previous step)
in the connected space $\mathbb{R}^n$, hence $f(\mathbb{R}^n)=\mathbb{R}^n$.
Thus $f$ is bijective; by the inverse function theorem $f^{-1}$ is $C^1$.
Therefore $f$ is a global $C^1$ \emph{diffeomorphism}.

\bigskip\hrule\bigskip

\textbf{Why not with “$\le 1$” only?}\gap

If $\|Df-I\|\le 1$, invertibility of $Df$ can fail somewhere.
In $n=1$, let $\phi$ be a $C^\infty$ bump with $0\le \phi \le 1$ and set
$f'(x)=1-\phi(x)$ so that $|f'(x)-1|\le 1$ but $f'(0)=0$.
Then $f$ is not locally invertible near $0$, and the two–sided Lipschitz bound can fail.

\bigskip\hrule\bigskip

\textbf{Examples.}\gap

\begin{itemize}\setlength{\itemsep}{6pt}
	\item \emph{Scalar case:} $f(x)=x-\tanh x$ on $\mathbb{R}$.
	Then $f'(x)=1-\operatorname{sech}^2 x\in(0,1)$, so $\|Df-I\|=\sup_x|f'(x)-1|<1$.
	The bounds give $(1-\mu)|x-y|\le|f(x)-f(y)|\le(1+\mu)|x-y|$,
	and $f$ is a global $C^1$ diffeomorphism of $\mathbb{R}$.
	
	\item \emph{Vector case:} $f(x)=x+g(x)$ with $\|Dg(x)\|\le \mu<1$ (e.g.\ $g(x)=\varepsilon\sin x$ componentwise).
	Then $\|Df-I\|\le\mu$ and all conclusions apply; $f$ is bi-Lipschitz and globally invertible.
\end{itemize}

% Optional helper for small vertical spacing
% \newcommand{\gap}{\vspace{6pt}}

\section*{Problem 11 — Differentiability of the determinant function}

\textbf{Statement.}\gap

Let $\mathcal{M}_n$ be the space of $n\times n$ real matrices with a norm.  
Show that the determinant function
\[
\det:\mathcal{M}_n \longrightarrow \mathbb{R}
\]
is differentiable at the identity matrix $I$ with
\[
D\det(I)(H)=\operatorname{tr}(H).
\]
Deduce that $\det$ is differentiable at any invertible matrix $A$ with
\[
D\det(A)(H)=\det(A)\,\operatorname{tr}(A^{-1}H).
\]
Show further that $\det$ is twice differentiable at $I$ and find $D^2\det(I)$ as a bilinear map.

\bigskip\hrule\bigskip

\textbf{Theory \& Solution.}\gap

\textbf{(Derivative at the identity).}\;
For $A(t)=I+tH$, expand the determinant polynomially:
\[
\det(I+tH)=1+t\,\operatorname{tr}(H)+o(t).
\]
Hence $D\det(I)(H)=\operatorname{tr}(H)$.

\medskip

\textbf{(Derivative at a general invertible matrix).}\;
Write $\det(A+H)=\det(A)\det(I+A^{-1}H)$.  
By the previous step,
\[
\det(I+A^{-1}H)=1+\operatorname{tr}(A^{-1}H)+o(\|H\|),
\]
so
\[
D\det(A)(H)=\det(A)\,\operatorname{tr}(A^{-1}H).
\]

\medskip

\textbf{(Second derivative at $I$).}\;
Expand further:
\[
\det(I+tH)=1+t\,\operatorname{tr}(H)
+\tfrac{1}{2}\big((\operatorname{tr}H)^2-\operatorname{tr}(H^2)\big)t^2+o(t^2).
\]
Therefore
\[
D^2\det(I)(H,K)
=\operatorname{tr}(H)\operatorname{tr}(K)-\operatorname{tr}(HK),
\]
which is symmetric and bilinear. \hfill$\square$

\bigskip\hrule\bigskip

\textbf{Examples.}\gap
\begin{itemize}\setlength{\itemsep}{6pt}
	\item For $n=2$,
	\[
	\det\!\begin{pmatrix}1+t a & t b\\ t c & 1+t d\end{pmatrix}
	=1+t(a+d)+t^2(ad-bc),
	\]
	giving $D\det(I)(H)=a+d=\operatorname{tr}(H)$ and
	$D^2\det(I)(H,H)=2(ad-bc)=(\operatorname{tr}H)^2-\operatorname{tr}(H^2)$.
\end{itemize}

% ==========================================================
\section*{Problem 12 — Matrix squaring and local square root}

\textbf{Statement.}\gap

Define $f:\mathcal{M}_n\to\mathcal{M}_n$ by $f(A)=A^2$.  
Show that $f$ is continuously differentiable on $\mathcal{M}_n$.  
Deduce that there exists $\varepsilon>0$ such that, on the open ball
$B_\varepsilon(I)\subset\mathcal{M}_n$, there is a continuous map
\[
g:B_\varepsilon(I)\longrightarrow \mathcal{M}_n
\]
satisfying $g(A)^2=A$ for all $A\in B_\varepsilon(I)$.
Discuss whether one can define a continuous square-root function on the whole $\mathcal{M}_n$.

\bigskip\hrule\bigskip

\textbf{Theory \& Solution.}\gap

\textbf{(Derivative of $f$).}\;
For $A,H\in\mathcal{M}_n$,
\[
f(A+H)=(A+H)^2=A^2+AH+HA+H^2,
\]
hence
\[
Df(A)(H)=AH+HA.
\]
This is linear in $H$ and depends continuously on $A$,
so $f$ is $C^1$ on $\mathcal{M}_n$.

\medskip

\textbf{(Local inverse near $I$).}\;
At $A=I$, $Df(I)(H)=2H$, which is invertible as a linear map on $\mathcal{M}_n$.  
By the inverse function theorem, there exist $\varepsilon>0$
and a $C^1$ map $g:B_\varepsilon(I)\to\mathcal{M}_n$ such that
$g(I)=I$ and $f(g(A))=A$, i.e.\ $g(A)^2=A$ for all $A$ in $B_\varepsilon(I)$.

\medskip

\textbf{(Global impossibility).}\;
A continuous square-root map cannot exist globally on $\mathcal{M}_n$:
for example, $\det(g(A))^2=\det(A)$ forces $\det(g(A))$ to change sign
discontinuously when $\det(A)$ crosses $0$.

\bigskip\hrule\bigskip

\textbf{Examples.}\gap
\begin{itemize}\setlength{\itemsep}{6pt}
	\item For scalars ($n=1$), $f(x)=x^2$ and $Df(1)=2$.
	By the inverse function theorem, $x\mapsto\sqrt{x}$ exists and is smooth near $1$,
	but no continuous square root exists on all of $\mathbb{R}$.
	\item For matrices, the same holds: locally near $I$, each $A$ close to $I$ has a
	unique square root $g(A)$ obtained by a convergent binomial series.
\end{itemize}

% ==========================================================
\section*{Problem 13 — Local inversion on the curve $x^3+y^3-3xy=0$}

\textbf{Statement.}\gap

Let $C=\{(x,y)\in\mathbb{R}^2:\; x^3+y^3-3xy=0\}$ and define
\[
F:\mathbb{R}^2\longrightarrow\mathbb{R}^2,\qquad F(x,y)=(x,\,x^3+y^3-3xy).
\]
Show that $F$ is locally $C^1$-invertible at every point of
$C\setminus\{(0,0),(\tfrac{2}{3},\tfrac{2}{3})\}$.
That is, for such $(x_0,y_0)$ there are open sets
$U,V$ with $F:U\to V$ a $C^1$ bijection.
Find the derivative of the inverse and deduce that near each such point
the curve $C$ is the graph of a $C^1$ function $y=g(x)$.

\bigskip\hrule\bigskip

\textbf{Theory \& Solution.}\gap

Compute the Jacobian:
\[
DF(x,y)=
\begin{pmatrix}
	1 & 0\\[4pt]
	3x^2-3y & 3y^2-3x
\end{pmatrix},
\qquad
J(x,y)=\det DF(x,y)=3(y^2-x^2+y-x).
\]
The determinant vanishes when $y^2-x^2+y-x=0$, i.e.\
along the lines $y=x$ and $y=1-x$.

\medskip

\textbf{(Points where $DF$ is invertible).}\;
On $C$, the points where $J(x,y)\ne0$
correspond to $(x,y)$ not lying on those lines simultaneously,
which occurs except at $(0,0)$ and $(\tfrac{2}{3},\tfrac{2}{3})$.
Hence by the inverse function theorem, $F$ is locally a $C^1$ diffeomorphism there.

\medskip

\textbf{(Derivative of the inverse).}\;
For invertible $DF(x_0,y_0)$,
\[
D(F^{-1})(F(x_0,y_0)) = [DF(x_0,y_0)]^{-1}.
\]
Thus the slope of the corresponding branch of $C$,
viewed as $y=g(x)$, is given by
\[
g'(x_0)= -\,\frac{F_y(x_0,y_0)}{F_x(x_0,y_0)}
= \frac{3y_0^2-3x_0}{3x_0^2-3y_0}.
\]

\bigskip\hrule\bigskip

\textbf{Examples.}\gap
\begin{itemize}\setlength{\itemsep}{6pt}
	\item At $(1,0)$, $g'(1)=\tfrac{-3}{3-0}=-1$.
	\item At $(0,1)$, $g'(0)=\tfrac{3-0}{0-3}=-1$ (symmetric branch).
	Thus near these points $C$ is a smooth curve crossing the axes with slope $-1$.
\end{itemize}


\begin{figure}[!ht]
	\centering
	\includegraphics[width=.7\linewidth]{figures/MATH-ALLS-2_curve_C_xy_plane.png}
	\caption{\textbf{Source plane (x,y).} The curve $C:\; x^3+y^3-3xy=0$.}
\end{figure}

\begin{figure}[!ht]
	\centering
	\includegraphics[width=.7\linewidth]{figures/MATH-ALLS-2_image_FC_uv_plane.png}
	\caption{\textbf{Target plane (u,v).} Image $F(C)$ with $F(x,y)=(x,x^3+y^3-3xy)$. On $C$ we have $v=0$, so $F(C)$ lies on the $u$–axis.}
\end{figure}


% Optional helper once in the preamble
% \newcommand{\gap}{\vspace{6pt}}

\section*{Problem 14\textsuperscript{*} — Maximum principles for C\textsuperscript{2} functions}


\textbf{(i) Local maximum: gradient and Hessian.}\gap

Let $f$ be a real-valued $C^2$ function on an open $U\subset\mathbb{R}^2$.
If $f$ has a local maximum at $a\in U$ (i.e.\ $\exists\,\rho>0$ with $B_\rho(a)\subset U$ and $f(x)\le f(a)$ for $x\in B_\rho(a)$),
then
\[
\nabla f(a)=0
\quad\text{and}\quad
H(a):=\big(D_{ij}f(a)\big)\ \text{is negative semidefinite.}
\]

\bigskip\hrule\bigskip

\textbf{Proof.}\gap
For any unit $v$, consider $g(t)=f(a+tv)$ for small $t$.
Then $g$ has a local maximum at $0$, hence $g'(0)=0$ and $g''(0)\le 0$.
But $g'(0)=\nabla f(a)\!\cdot v$ and $g''(0)=v^\top H(a)\,v$.
Thus $\nabla f(a)\!\cdot v=0$ for all unit $v$, giving $\nabla f(a)=0$, and $v^\top H(a)\,v\le0$ for all $v$, so $H(a)\preceq 0$.
\hfill$\square$

\bigskip\hrule\bigskip

\textbf{(ii) Weak maximum principle for $\Delta + aD_1 + bD_2 + c$.}\gap

Let $U\subset\mathbb{R}^2$ be bounded and open, and let $f:\overline U\to\mathbb{R}$ be continuous on $\overline U$ and $C^2$ in $U$.
Assume
\[
\Delta f + a\,D_1 f + b\,D_2 f + c\,f \ \ge\ 0 \quad \text{in } U,
\]
where $\Delta f=D_{11}f+D_{22}f$, and $a,b,c$ are real-valued on $U$ with $c<0$ on $U$.
If $f$ is positive somewhere in $\overline U$, then
\[
\sup_{\overline U} f \;=\; \sup_{\partial U} f .
\]

\bigskip\hrule\bigskip

\textbf{Proof.}\gap
Let $M=\sup_{\overline U} f$.
If $M\le 0$, the claim is trivial.
Assume $M>0$.
If $M$ is attained at $x_0\in U$, then (i) gives $\nabla f(x_0)=0$ and $\Delta f(x_0)\le 0$.
Hence
\[
\big(\Delta f + aD_1 f + bD_2 f + c f\big)(x_0)
= \Delta f(x_0) + c(x_0) f(x_0)
\le c(x_0) M \;<\; 0,
\]
contradicting the assumed inequality.
Therefore every positive maximum is on $\partial U$, i.e.\ $\sup_{\overline U} f=\sup_{\partial U} f$.
\hfill$\square$

\bigskip\hrule\bigskip

\textbf{Uniqueness for the Dirichlet problem.}\gap

If $a,b,c$ are as above with $c<0$, $\varphi:\partial U\to\mathbb{R}$ is continuous, and $g:\mathbb{R}^2\to\mathbb{R}$ is given,
then there is \emph{at most one} continuous $f$ on $\overline U$, $C^2$ in $U$, solving
\[
\Delta f + aD_1 f + bD_2 f + c f = g \quad\text{in } U,
\qquad
f=\varphi \quad\text{on } \partial U.
\]
Indeed, if $f_1,f_2$ solve it, $h=f_1-f_2$ satisfies $\Delta h + aD_1 h + bD_2 h + c h=0$ in $U$ and $h=0$ on $\partial U$.
Applying the maximum principle to $h$ and $-h$ shows $h\equiv 0$.

\bigskip\hrule\bigskip

\textbf{Examples.}\gap

\begin{itemize}\setlength{\itemsep}{6pt}
	\item If $a=b=0$ and $c=-\lambda<0$, the operator is $\Delta - \lambda$.
	Any interior positive maximum would give $(\Delta - \lambda)f(x_0)\le -\lambda f(x_0)<0$, a contradiction.
	\item In one dimension ($U=(0,1)$), the statement reduces to:
	if $f\in C^2(0,1)\cap C[0,1]$ and $f'' - \lambda f \ge 0$ with $\lambda>0$, then
	$\max_{[0,1]} f = \max\{f(0),f(1)\}$.
\end{itemize}

\begin{figure}[!ht]
	\centering
	\includegraphics[width=\linewidth]{figures/MATH-ALLS-2_F_mapping_xy_to_uv_side_by_side.png}
	\caption{\textbf{Mapping $F(x,y)=(x,x^3+y^3-3xy)$.}
		The folium of Descartes $C$ in the $(x,y)$-plane (left) is mapped to the $u$-axis in the $(u,v)$-plane (right).}
\end{figure}

% --- Examples for Problem 14 (Maximum principles) ---

\noindent\textbf{Example A (harmonic):} $u(x,y)=x^2-y^2$ on the unit disk
$D=\{x^2+y^2<1\}$.
\begin{itemize}
	\item \textbf{Laplacian:} $\Delta u = u_{xx}+u_{yy}=2-2=0 \;\Rightarrow\; u$ is \emph{harmonic}.
	\item \textbf{Maximum principle (harmonic):} If $u$ is harmonic on a bounded domain and continuous on $\overline D$, then $u$
	attains its max/min \emph{only on $\partial D$} unless $u$ is constant.
	\item \textbf{Boundary values (polar coords).} With $x=\cos\theta$, $y=\sin\theta$ on $\partial D$,
	\[
	u(\cos\theta,\sin\theta)=\cos^2\theta-\sin^2\theta=\cos(2\theta).
	\]
	Hence on $\partial D$: $\max u=1$ at $\theta=0,\pi$ (points $(\pm1,0)$),
	and $\min u=-1$ at $\theta=\tfrac{\pi}{2},\tfrac{3\pi}{2}$ (points $(0,\pm1)$).
	\item \textbf{Interior critical point:} $\nabla u=(2x,-2y)=0 \Rightarrow (0,0)$.
	The Hessian $D^2u=\begin{pmatrix}2&0\\[2pt]0&-2\end{pmatrix}$ is indefinite $\Rightarrow$ \emph{saddle}, not extremum.
\end{itemize}

\begin{figure}[!ht]
	\centering
	\includegraphics[width=.72\linewidth]{figures/MATH-ALLS-2_problem14_subharmonic_heatmap.png}
	\caption{\textbf{Harmonic example.} $u(x,y)=x^2-y^2$ on the unit disk ($\Delta u=0$).
		Extrema occur only on $\partial D$: $\max u=1$ at $(\pm1,0)$; $\min u=-1$ at $(0,\pm1)$.
		The interior critical point $(0,0)$ is a saddle.}
\end{figure}

\medskip\hrule\medskip

\noindent\textbf{Example B (subharmonic):} $v(x,y)=x^2+y^2$ on the unit disk $D$.
\begin{itemize}
	\item \textbf{Laplacian:} $\Delta v = 2+2=4 \ge 0 \;\Rightarrow\; v$ is \emph{subharmonic}.
	\item \textbf{Maximum principle (subharmonic):} A subharmonic function attains its \emph{maximum} on $\partial D$ (unless constant).
	Interior maxima are forbidden.
	\item \textbf{Values:} On $\partial D$ ($r=1$): $v=1$ $\Rightarrow$ \emph{global maximum}. At the center: $v(0,0)=0$
	$\Rightarrow$ \emph{global minimum} (allowed for subharmonic).
\end{itemize}

\begin{figure}[!ht]
	\centering
	\includegraphics[width=.72\linewidth]{figures/MATH-ALLS-2_problem14_subharmonic_heatmap.png}
	\caption{\textbf{Subharmonic example.} $v(x,y)=x^2+y^2$ on the unit disk ($\Delta v=4$).
		$v$ increases radially and peaks on $\partial D$; the minimum is at the center.}
\end{figure}


\end{document}
